\documentclass[11pt,a4paper]{article}

\usepackage[utf8]{inputenc}
\usepackage[T1]{fontenc}
\usepackage{lmodern}
\usepackage[margin=1.0in]{geometry}
\usepackage{amsmath,amssymb,amsfonts,amsthm}
\usepackage{mathtools,bm,booktabs,array,enumitem,listings,xcolor,tabularx}
\usepackage{hyperref}
\hypersetup{colorlinks=true,linkcolor=blue!60!black,citecolor=blue!60!black,urlcolor=blue!60!black,
  pdftitle={Hierarchical Bio-Kagome Solid-State Metamaterials (v5.5)},
  pdfauthor={Nicholas Clifford Maino III}}

\newtheorem{theorem}{Theorem}[section]
\newtheorem{lemma}[theorem]{Lemma}
\newtheorem{proposition}[theorem]{Proposition}
\newtheorem{corollary}[theorem]{Corollary}
\newtheorem{conjecture}[theorem]{Conjecture}
\theoremstyle{definition}
\newtheorem{definition}[theorem]{Definition}
\theoremstyle{remark}
\newtheorem{remark}[theorem]{Remark}

\lstdefinestyle{cleanstyle}{
  basicstyle=\ttfamily\scriptsize,breaklines=true,frame=single,
  numbers=left,numberstyle=\tiny\color{gray},
  keywordstyle=\color{blue!70!black}\bfseries,
  commentstyle=\color{green!50!black}\itshape,
  stringstyle=\color{red!60!black},showstringspaces=false,tabsize=4,captionpos=b,
  columns=fullflexible,keepspaces=true}
\lstdefinestyle{consolestyle}{
  basicstyle=\ttfamily\scriptsize,breaklines=true,frame=single,columns=fullflexible,keepspaces=true,
  breakindent=2em,captionpos=b}

\title{\textbf{Hierarchical Bio-Kagome Solid-State Metamaterials:\\
Anisotropic Phonon-Electron Decoupling, Compact Localized Flat-Band States,\\
and Identifiability of Three- and Four-Chord Schur-Complement Boundary Tomography}}
\author{\textbf{Nicholas Clifford Maino III}\\
Independent Computational \& Materials Research, Melbourne, Victoria, Australia\\
\texttt{License: CC-BY-4.0}}
\date{October 31, 2026 (Revision 5.5; supersedes Revisions 4.0 and 5.0--5.4)}

\begin{document}
\maketitle

\begin{abstract}
\noindent We formulate a two-tier ($N=2$) hierarchical bio-composite metamaterial (\textit{BioLux Lattice}), a triaxially woven hemp and chitin--mycelium scaffold with Cu-graphene/BaTiO$_3$ nodes and h-BN thermal relays, designed as a passive solid-state switch. We derive the exact $3\times3$ Bloch spectrum of the breathing Kagome lattice and an exact compact localized state at energy $t_1+t_2$, and certify the positive orientation of all 132{,}019 perturbed annulus triangles. We then determine what the diagonal of the perimeter Schur complement and chord line sums of the cell states can identify. The diagonal boundary response is strictly monotone in every cell state at the 69 perimeter nodes with an interior neighbour, but perimeter nodes bonded to the same interior node carry identical data, so the boundary Jacobian has rank at most 45 at every state. With four chord families the linearised operator has exact rank 112 of 169; reading each chord in $k$ interleaved parts raises it to 163 for $k=2$ and 169 for $k=3$ at $R_{\text{hex}} = 8$. The kernels of the split chords contain every translate of the switching component dilated by $k$, and are spanned by them in every case computed. Closed ridge loops screen such components: inside a closed 36-cell loop the undilated pairs are separated by at most $6.3\times10^{-4}$ noise standard deviations at 34 dB, so every decoder errs with probability at least 0.4998. Because the chord kernels grow quadratically with the array size while the boundary rank grows linearly, no fixed split order gives full rank for every size. The four chord directions have cross ratio 2, which the Gardner--Gritzmann criterion excludes: two mirror-image convex 8-cell ridges share all chord data, at 102 placements. The switching-component results are instances of the Hajdu--Tijdeman theory of discrete tomography; the counts on the hexagonal domain and the comparison with the boundary rank are new.
\end{abstract}

\noindent \textbf{Keywords:} Mechanical Metamaterials, Breathing Kagome, Compact Localized States, Laser-Induced Graphene, Dirichlet-to-Neumann, Schur Complement, Identifiability, Discrete Tomography.

\vspace{0.5em}\hrule\vspace{0.5em}

% ============================================================================
\section{Introduction}
% ============================================================================

Conventional interactive surfaces and biometric arrays rely on continuous active capacitive polling, multilayer PCBs, active-matrix TFT backplanes, and 2D affine transformations --- incurring continuous static power ($P_{\text{static}} > 0$), harness mass, and EMI vulnerability. We formalize \textit{physical-computational displacement}: replacing active silicon logic and coordinate transforms by passive mechanical constraints, material phase transitions, and graph-theoretic boundary topologies. Specifically: (i) \textbf{rotational affine eradication} via a kinematic dual-finger mold ($\theta_{\text{yaw}} = 0^\circ \Rightarrow \mathbf{R} = \mathbf{I}_2$); (ii) \textbf{hardware interrupt gating} via a $0$--$3$ piezoelectric film; and (iii) \textbf{passive boundary tomography} via chord currents along three lattice directions (four in Revision~4.0) and an augmented Schur-complement decoder~\cite{calderon1980,schur1917}, whose identifiability we determine exactly (Section~\ref{sec:inv}). Section~\ref{sec:nonlin} treats the nonlinear problem (monotonicity, the rank of the boundary Jacobian, screening by closed ridge loops, seeded ridge classes), Section~\ref{sec:split} chords read in interleaved parts, and Section~\ref{sec:convex} convex ridges; related work is in Section~\ref{sec:related}.

\subsection{Related work}\label{sec:related}
\noindent\textbf{Scope and prior-art grounding.} All identifiability, monotonicity, data-collapse, and split-chord theorems in Section~\ref{sec:inv} are substrate-agnostic: they hold for any passive resistor-shunt network sharing this Breathing Kagome graph topology and conductance regime, independent of the illustrative bio-composite material parameters modelled in Sections~2--4. The switching-component constructions in Theorems~\ref{thm:ghost}(5) and~\ref{thm:splitker}(1) are specific instances of Hajdu--Tijdeman algebraic discrete tomography~\cite{hajdu2001} (extending Ryser's switching operation~\cite{ryser1957}), while the convex-uniqueness cross-ratio obstruction in Proposition~\ref{prop:convex} applies the criterion of Gardner and Gritzmann~\cite{gardner1997} to hexagonal-grid tomography~\cite{matej1998,nagy2016}. Within this classical framework, our specific contributions are the closed-form dilated kernel count $T_k(R)$ on the hexagonal domain (Theorem~\ref{thm:splitker}), the global nonlinear boundary data-collapse bound $\operatorname{rank}\mathbf{J}(\mathbf{s})\le 45$ (Theorem~\ref{thm:collapse}), the asymptotic split-chord rank breakdown for fixed $k$ (Theorem~\ref{thm:splitscale}), and the explicit realization of the Gardner--Gritzmann cross-ratio-2 convex obstruction on the $R_{\text{hex}}=8$ lattice (Proposition~\ref{prop:convex}).

\emph{Discrete tomography.} The chord data of Section~\ref{sec:inv} are discrete X-rays: line sums of a binary function on a lattice set in a few lattice directions. Functions with zero line sums are the switching components of discrete tomography, going back to Ryser's interchanges for rows and columns~\cite{ryser1957}. Hajdu and Tijdeman characterised them completely: on a rectangular grid every function with zero line sums along a finite set of directions is a combination of translates of the product of the binomials belonging to the directions~\cite{hajdu2001}, and they applied this to rows, columns, diagonals and anti-diagonals~\cite{hajdu2001laa}, which are the four chord families here in axial coordinates. Stolk and Batenburg describe the structure of the line sums for finite convex sets~\cite{stolk2010}. The switching component $\mathbf{z}$ of Theorem~\ref{thm:ghost}(5) and its dilations in Theorem~\ref{thm:splitker} are instances of this theory; what is specific here is the count of their translates in the hexagonal domain, the computed spanning, and the comparison with the rank of the boundary data.

\emph{Uniqueness.} Gardner and Gritzmann proved that four suitably chosen lattice directions determine convex lattice sets by their X-rays, that three never suffice, and that any seven do; in particular four directions suffice whenever the cross ratio of their slopes avoids $\{4/3,3/2,2,3,4\}$~\cite{gardner1997}. The four chord families have cross ratio 2, so convex ridges are not guaranteed to be determined, and Section~\ref{sec:convex} gives the explicit pair.

\emph{Triangular and hexagonal grids.} Binary tomography on the hexagonal grid~\cite{matej1998} and on the triangular grid with projections along its lanes~\cite{nagy2016} has been studied for reconstruction algorithms; the three-family chords here are such lane sums.

\emph{Resistor networks and impedance tomography.} For circular planar resistor networks the conductances are recoverable from the full response matrix exactly for critical networks~\cite{curtis1998,colin1996}. Here only the diagonal of the response matrix is measured and the unknowns are binary cell states, so those results do not apply directly. Lemma~\ref{lem:monotone}(1) without strictness is Rayleigh's monotonicity law~\cite{doyle1984} for the response matrix, and its continuum counterpart underlies the monotonicity methods of impedance tomography~\cite{tamburrino2002,harrach2013}.

% ============================================================================
\section{Mass-Consistent Stoichiometry and $N=2$ Scaling}
% ============================================================================

\subsection{Five-Component Closure}
Accounting for hyphal microporosity ($\phi \approx 87\%$) and Kagome void fraction $V_5$~\cite{islam2017,danninger2022}:
\begin{equation}
\rho_{\text{eff}} = \sum_{i=1}^5 V_i \rho_i = (0.220)(1480) + (0.450)(180) + (0.050)(4200) + (0.044)(1450) + (0.236)(1.2) = 680.7\ \text{kg/m}^3.
\end{equation}

\begin{table}[h!]\centering\small
\begin{tabular}{@{}llcccp{4.5cm}@{}}
\toprule
\textbf{Phase} & \textbf{Constituent} & $\rho_i$ (kg/m³) & $V_i$ (\%) & $W_i$ (\%) & \textbf{Function} \\
\midrule
I   & Mercerized \textit{C. sativa} bast + CNF & 1480 & 22.0 & 47.83 & Isostatic Kagome frame~\cite{maxwell1864,guest2003} \\
II  & \textit{T. versicolor} chitin-glucan foam & 180 & 45.0 & 11.90 & $z$-axis dielectric, LOI = 34\% \\
III & Cu-Gr nodes + 20 vol\% $\text{BaTiO}_3$ film & 4200 & 5.0 & 30.85 & Piezo wake-up, tunneling, antimicrobial \\
IV  & h-BN + C18 PCM & 1450 & 4.4 & 9.37 & Phonon-relay, latent heat \\
V   & Inter-strut air voids & 1.2 & 23.6 & 0.04 & Compression accommodation \\
\midrule
\textbf{Total} & & \textbf{680.7} & \textbf{100.0} & \textbf{100.00} & \textbf{Exact closure} \\
\bottomrule
\end{tabular}
\caption{Five-component volumetric and gravimetric closure ($W_i$ rounded; unrounded values sum to 100).}
\end{table}

\subsection{Two-Tier Scaling}
Following Fleck--Deshpande--Ashby~\cite{fleck2010} and Gibson--Ashby~\cite{gibson1997}: Tier 1 macro-Kagome ($a_1 = 2$ mm) is stretching-dominated, $b_1 = 1.0$, giving $E_{xy} = \tfrac{1}{3}E_{\text{hemp}}\bar\rho_1^{1.0}$. Tier 2 micro-Kagome ($a_2 = 50\,\mu$m) is bending-dominated in the porous chitin-CNF core, $b_2 = 1.5$, giving $E_{z,\text{hier}} = E_{\text{chitin-CNF}}\bar\rho_1^{1.0}\bar\rho_2^{1.5} \approx 45$ MPa.

% ============================================================================
\section{Multidomain Coupled Physics}
% ============================================================================

\subsection{Mechanical Isostasy and Piezoelectric Wake-Up}
Maxwell--Calladine counting for $d=2$, $z=4$ gives $M = 0$~\cite{maxwell1864,guest2003}, localizing strain to the loaded cell. The Phase III piezo element is a $4\,\mu$m-thick BaTiO$_3$/chitin top-cap film (20 vol\% loading). For $\phi_p = 0.20$, $\varepsilon_{\text{BT}} = 1700$, $\varepsilon_m = 3.2$~\cite{yamada1982}:
\begin{equation}
\varepsilon_{\text{eff}} = \varepsilon_m \frac{1 + 2\phi_p\beta}{1-\phi_p\beta},\quad \beta = \frac{\varepsilon_{\text{BT}}-\varepsilon_m}{\varepsilon_{\text{BT}}+2\varepsilon_m} = 0.994,\quad \varepsilon_{\text{eff}} = 5.58.
\end{equation}
With empirical $d_{33,\text{eff}} = 2.0$ pC/N, $g_{33,\text{eff}} = d_{33,\text{eff}}/(\varepsilon_0\varepsilon_{\text{eff}}) = 40.5\times 10^{-3}$ V·m/N. Under $F_z = 1.5$ N, $A_{\text{node}} = 3.46$ mm² (one $a_1$ unit cell), $h_{\text{node}} = 4\,\mu$m:
\begin{equation}
V_{\text{wake}} = g_{33,\text{eff}}\frac{F_z}{A_{\text{node}}}h_{\text{node}} = 70.2\ \text{mV} > V_{\text{th}} = 50.0\ \text{mV},
\end{equation}
triggering the nanopower comparator ($P_{\text{static}} \ll 1\,\mu$W) while compressing the LIG tunneling barrier~\cite{simmons1963}.

\subsection{Electronic Domain: Exact Bloch Spectrum, CLS, and Bulk Gap}
The $3\times 3$ Bloch Hamiltonian with Bravais vectors $\mathbf{a}_1 = a(1,0)$, $\mathbf{a}_2 = a(\tfrac{1}{2},\tfrac{\sqrt3}{2})$, $\mathbf{a}_3 = \mathbf{a}_2 - \mathbf{a}_1$, intracell ($\triangle^{\uparrow}$) hopping $t_1 = \tfrac12 t_\Sigma(1+\delta_b)$ and intercell ($\triangle^{\downarrow}$) hopping $t_2 = \tfrac12 t_\Sigma(1-\delta_b)$ is~\cite{ezawa2018,kempkes2019}
\begin{equation}
\mathcal{H}(\mathbf{k}) = -\begin{bmatrix}
0 & t_1+t_2e^{-i\mathbf{k}\cdot\mathbf{a}_1} & t_1+t_2e^{-i\mathbf{k}\cdot\mathbf{a}_2}\\
\text{c.c.} & 0 & t_1+t_2e^{-i\mathbf{k}\cdot\mathbf{a}_3}\\
\text{c.c.} & \text{c.c.} & 0
\end{bmatrix}.
\end{equation}

\begin{theorem}[Exact Bloch Spectrum]\label{thm:bloch}
With $\Lambda_{\mathbf{k}} = 3 + 2\sum_i\cos(\mathbf{k}\cdot\mathbf{a}_i)\in[0,9]$:
\begin{equation}
E^{(1)}_{\mathbf{k}} = +(t_1+t_2),\qquad E^{(2,3)}_{\mathbf{k}} = -\frac{t_1+t_2}{2} \pm \sqrt{\tfrac{9}{4}(t_1-t_2)^2 + t_1 t_2\Lambda_{\mathbf{k}}}.
\end{equation}
In particular the two dispersive bands are separated by the bulk gap $\Delta E_{\mathbf{K}} = 3|t_1-t_2|$, attained at $\mathbf{K}$ ($\Lambda_{\mathbf{K}}=0$).
\end{theorem}
\begin{proof}
Expanding $\det(E\mathbf{I}_3 - \mathcal{H}(\mathbf{k})) = 0$ gives $E^3 - C_1 E - C_0 = 0$ with $C_1 = 3(t_1^2+t_2^2) + t_1t_2(\Lambda_{\mathbf{k}}-3)$, $C_0 = -2(t_1^3+t_2^3)-t_1t_2(t_1+t_2)(\Lambda_{\mathbf{k}}-3)$. Substituting $E_0 = t_1+t_2$ annihilates the cubic. Dividing by $(E-E_0)$ gives
$E^2 + (t_1+t_2)E - 2t_1^2 - 2t_2^2 + 5t_1t_2 - t_1t_2\Lambda_{\mathbf{k}} = 0$,
whose discriminant is $9(t_1-t_2)^2 + 4t_1t_2\Lambda_{\mathbf{k}}$, yielding $E^{(2,3)}_{\mathbf{k}}$. The splitting $2\sqrt{\cdot}$ is minimal where $\Lambda_{\mathbf{k}}$ is, i.e.\ at $\mathbf{K}$. Verified at 400 random $\mathbf{k}$ to within $2.3\times10^{-15}$ (App.~\ref{app:code}).
\end{proof}

\begin{theorem}[$\Gamma$-Degeneracy and Exact Real-Space CLS]\label{thm:cls}
(i) At $\Gamma$ ($\Lambda_\Gamma = 9$): $E^{(2)}_\Gamma = -(t_1+t_2)/2 + \tfrac{3}{2}(t_1+t_2) = +(t_1+t_2) = E^{(1)}_\Gamma$; the flat and upper dispersive bands touch quadratically via the Bergman--Wu--Balents mechanism~\cite{bergman2008}. (ii) For any internal hexagonal plaquette $H$ with perimeter $\partial H = \{v_1,\dots,v_6\}$,
$|\Psi^{(H)}_{\text{CLS}}\rangle = 6^{-1/2}\sum_{j=1}^6(-1)^j|v_j\rangle$
is an exact eigenstate of the full real-space Hamiltonian at energy $t_1+t_2$ with strictly zero amplitude outside $\partial H$ for all $t_1,t_2>0$.
\end{theorem}
\begin{proof}
(i) Direct evaluation at $\Lambda_\Gamma=9$. (ii) For any exterior vertex $v_{\text{ext}}$ sharing an edge $(v_{2m-1},v_{2m})$ with $\partial H$, the hopping amplitudes $(-t)(-1/\sqrt6) + (-t)(+1/\sqrt6) = 0$ cancel identically for $t \in \{t_1,t_2\}$. For each $v_j\in\partial H$, its two perimeter neighbors have opposite signs, yielding $(H\Psi_{\text{CLS}})(v_j) = (t_1+t_2)\Psi_{\text{CLS}}(v_j)$. Numerical residual $2.7\times10^{-17}$ (App.~\ref{app:code}).
\end{proof}

\begin{remark}[Which phase the design is in]\label{rem:phase}
$\Delta E_{\mathbf{K}}$ is invariant under $t_1\leftrightarrow t_2$, so Theorem~\ref{thm:bloch} cannot by itself identify a higher-order topological phase. For $\delta_b>0$ ($t_1>t_2$) the gap stays open along $t_2\to 0$, which connects the lowest band to decoupled $\triangle^{\uparrow}$ trimers with Wannier centres at the cell centres: the trivial atomic phase for a sample cut along complete $\triangle^{\uparrow}$ cells. Numerically (App.~\ref{app:code}) the $R_{\text{hex}}=8$ flake (217 cells) has exactly 217 states below the gap and none inside it at $\delta_b=+0.18$, whereas at $\delta_b=-0.18$ it has 204 states below the gap and 21 in-gap (hence boundary-localized) states --- the filling anomaly of the obstructed phase. The HOTI designation used up to Revision~3.8 therefore applies to $\delta_b<0$, not to the design value; corner-mode phenomenology of the breathing Kagome HOTI~\cite{ezawa2018} likewise requires the intracell hopping to be the weaker one.
\end{remark}

\subsection{Thermodynamic Domain}
h-BN nanosheets~\cite{zeng2015} with $G_K = 1.5\times 10^8$ W/m²K give $k_{xx,\text{eff}} \approx 185$ W/m·K vs.\ $k_{zz} \approx 0.08$ ($k_{xx}/k_{zz}\approx 2312$). C18 PCM buffers transients.

% ============================================================================
\section{UV-LIG and Galvo Diffeomorphism}
% ============================================================================

Cold photochemical graphitization window~\cite{lin2014}: $F_{\text{th,LIG}} < \frac{2E_p}{\pi\omega_0^2}\exp(-2\alpha_{\text{eff}}z_p) < F_{\text{ablate,h-BN}}$. At $\lambda = 355$ nm, $\tau_p = 12$ ps: $\triangle^{\uparrow}$ islands at 100\% AOM ($F = 0.420$ J/cm²) yield $R_s = 18\,\Omega$/sq ($g_1 = 2.168\times 10^{-2}$ S, i.e.\ strut aspect ratio $L/W = 2.56$); $\triangle^{\downarrow}$ bridges at 35\% AOM ($F = 0.147$ J/cm²) yield $R_s \approx 10^7\,\Omega$/sq ($g_2^{(0)} = 10^{-7}$ S at $L/W=1$, compressing to $10^{-3}$ S under load).

\begin{theorem}[Exact $C^2$ Radial Metric Integral]\label{thm:radial}
For $\rho_c = 0.80$, $u = (\rho-\rho_c)/(1-\rho_c)$, $\Delta_{\text{rel}} = (a_1-a_2)/a_2 = 39.0$, and quintic Hermite $a(\rho) = a_2 + (a_1-a_2)(6u^5-15u^4+10u^3)$:
\begin{equation}
R(\rho) = \begin{cases}
\rho, & \rho \le \rho_c\\
\rho + \Delta_{\text{rel}}(1-\rho_c)(u^6 - 3u^5 + \tfrac{5}{2}u^4), & \rho_c<\rho\le 1\\
\rho_c + (1-\rho_c)(1+\tfrac12\Delta_{\text{rel}}) + (1+\Delta_{\text{rel}})(\rho-1), & \rho>1.
\end{cases}
\end{equation}
\end{theorem}
\begin{proof}
$R' = 1 + \Delta_{\text{rel}}(6u^5-15u^4+10u^3) = a(\rho)/a_2$; $R'' = 30\Delta_{\text{rel}}(1-\rho_c)^{-1}u^2(1-u)^2\to 0$ at $u\in\{0,1\}$; the three branches and their first derivatives agree at $\rho_c$ and at $1$.
\end{proof}

The vertex map is $\mathbf{\Phi}(\mathbf{r}) = \mathbf{\Phi}_0(\mathbf{r}) + \eta_r a_2 \frac{R(\rho)}{\rho}W(u)\,\mathbf{u}(\psi)$, with $\mathbf{\Phi}_0(\mathbf{r}) = \frac{R(\rho)}{\rho}\mathbf{r}$, $\rho = \|\mathbf{D}^{-1}\mathbf{r}\|_2$, $\mathbf{D} = \operatorname{diag}(r_x,r_y) = \operatorname{diag}(10,12.5)$ mm, bump $W(u) = 16u^2(1-u)^2$ on $0<u<1$ (zero elsewhere), $\mathbf{u}(\psi) = (\cos\psi,\sin\psi)^T$ and $\eta_r = 0.080$. Triangles have reference side $\ell = \alpha_b a_2$ ($\triangle^{\uparrow}$, $\alpha_b = \tfrac12(1-\delta_b) = 0.41$) or $\ell = (1-\alpha_b)a_2$ ($\triangle^{\downarrow}$).

\begin{theorem}[Intracell Cancellation and Certified Simplicial Orientation]\label{thm:jac}
Let the transition phase be the composite \emph{Eisenstein-Golden wavevector projection}
\begin{equation}
\psi_{m,n,k} = 2\pi\varphi^{-1}\left(m^2 + mn + n^2\right) + \mathbf{q}_\varphi \cdot \mathbf{r}_{m,n,k} + \frac{2\pi(k-1)}{3},
\end{equation}
with $\mathbf{q}_\varphi = \frac{2\pi}{a_2}\left(\varphi^{-1}, \varphi^{-2}\right)$, $\|\mathbf{q}_\varphi\|_2 = 4.5650/a_2$ and $\varphi = \frac{1+\sqrt{5}}{2}$. Write $\mathbf{E}_0 = [\mathbf{P}_2^0-\mathbf{P}_1^0, \mathbf{P}_3^0-\mathbf{P}_1^0]$ for the edge matrix of a triangle under $\mathbf{\Phi}_0$ and $\Delta\mathbf{E} = [\bm{\delta}_2-\bm{\delta}_1, \bm{\delta}_3-\bm{\delta}_1]$ for its perturbation.
\begin{enumerate}[noitemsep]
    \item \textbf{Exact intracell cancellation.} In every cell, $I_{m,n} = m^2+mn+n^2$ cancels from all intracell differences, $\psi_{m,n,j} - \psi_{m,n,k} = \mathbf{q}_\varphi \cdot (\mathbf{r}_{m,n,j} - \mathbf{r}_{m,n,k}) + \frac{2\pi(j-k)}{3}$, so the intracell offsets are cell-independent constants:
    \begin{align}
    \psi_{m,n,2}-\psi_{m,n,1} &\equiv 2\pi\bigl(\alpha_b\varphi^{-1} + \tfrac13\bigr) = 211.22^\circ,\\
    \psi_{m,n,3}-\psi_{m,n,1} &\equiv 2\pi\bigl(\tfrac{\alpha_b}{2}(\varphi^{-1}+\sqrt3\varphi^{-2}) + \tfrac23\bigr) = 334.44^\circ \pmod{360^\circ}.
    \end{align}
    The $\triangle^{\uparrow}$ chord pair is therefore a rigid rotation of one fixed matrix $\mathbf{C}_{\uparrow}$, $\|\mathbf{C}_{\uparrow}\|_2 = 1.9382$.
    \item \textbf{First-order bound.} Let $\kappa = r_y/r_x$ and $\mu = \tfrac12(\kappa+\kappa^{-1}) = 1.025$. To first order in $\ell/((1-\rho_c)r_x)$,
    \begin{equation}
    \|\mathbf{E}_0^{-1}\Delta\mathbf{E}\|_2 \le \frac{4\mu\eta_r a_2}{\ell} = 0.800\ (\triangle^{\uparrow}),\quad 0.556\ (\triangle^{\downarrow}).
    \end{equation}
    The first-order suprema are $0.5451$ ($\triangle^{\uparrow}$) and $4\eta_r/(1-\alpha_b) = 0.5424$ ($\triangle^{\downarrow}$, attained at $u=\tfrac12$, $\theta = \pm\pi/2$).
    \item \textbf{Exhaustive certificate.} Over all 65{,}879 $\triangle^{\uparrow}$ and 66{,}140 $\triangle^{\downarrow}$ triangles carrying a non-zero perturbation (at least one vertex with $0<u<1$), exact evaluation gives
    \begin{equation}
    \max_{\triangle^{\uparrow}} \|\mathbf{E}_0^{-1}\Delta\mathbf{E}\|_2 = 0.55089, \qquad \max_{\triangle^{\downarrow}} \|\mathbf{E}_0^{-1}\Delta\mathbf{E}\|_2 = 0.52212,
    \end{equation}
    $\min\det(\mathbf{I}_2+\mathbf{E}_0^{-1}\Delta\mathbf{E}) = 0.57854$ ($\triangle^{\uparrow}$), $0.69278$ ($\triangle^{\downarrow}$) and $\det\mathbf{E}_0>0$ throughout. Hence every mapped triangle is strictly positively oriented, with $\min \det(\mathbf{E}_0+\Delta\mathbf{E})/\det_{\text{ref}} = 0.9988$ ($\triangle^{\uparrow}$), $0.9993$ ($\triangle^{\downarrow}$), $\det_{\text{ref}} = \tfrac{\sqrt3}{2}\ell^2$.
\end{enumerate}
\end{theorem}
\begin{proof}
(1) $I_{m,n}$ does not depend on $k$. With $\mathbf{r}_{m,n,2}-\mathbf{r}_{m,n,1} = \alpha_b a_2(1,0)$ and $\mathbf{r}_{m,n,3}-\mathbf{r}_{m,n,1} = \alpha_b a_2(\tfrac12,\tfrac{\sqrt3}{2})$, $\mathbf{q}_\varphi\cdot(\cdot)$ equals $2\pi\alpha_b\varphi^{-1}$ and $2\pi\alpha_b(\varphi^{-1}+\sqrt3\varphi^{-2})/2$. Since $\mathbf{u}(\psi_1+\Delta)-\mathbf{u}(\psi_1) = \mathbf{R}_{\psi_1}(\mathbf{u}(\Delta)-\mathbf{e}_1)$, the $\triangle^{\uparrow}$ chords are $\mathbf{R}_{\psi_{m,n,1}}\mathbf{C}_{\uparrow}$.

(2) Write $f = R/\rho$ and $\hat{\mathbf{e}} = (\cos\theta,\sin\theta)$ with $\mathbf{r} = \rho\mathbf{D}\hat{\mathbf{e}}$. Then $\nabla\rho = \mathbf{D}^{-1}\hat{\mathbf{e}}$ and
$D\mathbf{\Phi}_0 = f(\mathbf{I}+\gamma\mathbf{Q})$ with $\gamma = R'/f - 1 \ge 0$ (because $R'$ is non-decreasing, $R(\rho)=\int_0^\rho R' \le \rho R'(\rho)$) and $\mathbf{Q} = \mathbf{D}\hat{\mathbf{e}}\hat{\mathbf{e}}^T\mathbf{D}^{-1}$, an idempotent. Hence $f\,(D\mathbf{\Phi}_0)^{-1} = \mathbf{I}-\beta\mathbf{Q}$, $\beta = \gamma/(1+\gamma)\in[0,1)$. By the Kantorovich inequality $\|\mathbf{Q}\|_2 = \|\mathbf{D}\hat{\mathbf{e}}\|\,\|\mathbf{D}^{-1}\hat{\mathbf{e}}\| \le \mu$, and since $\|\mathbf{I}-\mathbf{Q}\|_2 = \|\mathbf{Q}\|_2$ for any idempotent $\mathbf{Q}\notin\{\mathbf{0},\mathbf{I}\}$~\cite{szyld2006}, $\|\mathbf{I}-\beta\mathbf{Q}\|_2 \le (1-\beta)+\beta\|\mathbf{Q}\|_2 \le \mu$. To first order in the cell size, $\mathbf{E}_0 = D\mathbf{\Phi}_0\mathbf{E}_{\text{ref}}$ and $\Delta\mathbf{E} = \eta_r a_2 fW\,\mathbf{X}$, where $\mathbf{X}$ has two unit-circle chord columns ($\|\mathbf{X}\|_2 \le \|\mathbf{X}\|_F \le 2\sqrt2$). Thus $\mathbf{E}_0^{-1}\Delta\mathbf{E} = \eta_r W\frac{a_2}{\ell}\,\mathbf{T}(\mathbf{I}-\beta\mathbf{Q})\mathbf{X}$ with $\mathbf{T} = (\mathbf{E}_{\text{ref}}/\ell)^{-1}$, $\|\mathbf{T}\|_2 = \sqrt2$ (both reference triangles have singular values $\sqrt{3/2},\sqrt{1/2}$), giving $\sqrt2\cdot\mu\cdot2\sqrt2\,\eta_r a_2/\ell$. For $\triangle^{\downarrow}$ the three phases are free, so the supremum over phases is $2\sqrt2\|\mathbf{T}(\mathbf{I}-\beta\mathbf{Q})\|_2$; at $\theta=\pm\pi/2$, $\mathbf{Q}\mathbf{e}_1 = \mathbf{0}$ and $\mathbf{e}_1$ is the top singular direction of $\mathbf{T}$, so the value $4\eta_r/(1-\alpha_b)$ is attained at $u=\tfrac12$; a grid over $(u,\theta)$ finds no larger value. The $\triangle^{\uparrow}$ supremum ($\mathbf{X} = \mathbf{R}_{\psi}\mathbf{C}_{\uparrow}$) is evaluated on a refined $(u,\theta,\psi)$ grid (App.~\ref{app:code}).

(3) For a real $2\times2$ matrix $\mathbf{M}$ with $\|\mathbf{M}\|_2<1$, every eigenvalue satisfies $|\lambda|<1$; real eigenvalues give $1+\lambda>0$ and a complex pair gives $|1+\lambda|^2>0$, so $\det(\mathbf{I}+\mathbf{M})>0$ and $\det(\mathbf{E}_0+\Delta\mathbf{E}) = \det\mathbf{E}_0\det(\mathbf{I}+\mathbf{E}_0^{-1}\Delta\mathbf{E})>0$. The values in (3) come from evaluating every triangle in closed form (App.~\ref{app:code}). A forward error bound that allows $64u$ relative error per mapped coordinate is at least $1.7\times10^{10}$ times smaller than each determinant, so rounding cannot change any sign. The $\triangle^{\uparrow}$ maximum $0.5509$ exceeds its first-order supremum $0.5451$ by $1.1\%$ through $\mathcal{O}(\ell)$ amplitude-gradient and curvature terms; the first-order numbers in (2) are explanatory, and the certificate for the actual lattice is the exhaustive evaluation.
\end{proof}

% ============================================================================
\section{Full $3N\times 3N$ Graph Laplacian and Boundary Inversion}\label{sec:inv}
% ============================================================================

\subsection{Full 3-Basis Breathing Kagome Laplacian, Forward Model and Chords}
The network $\mathcal{G}=(\mathcal{V},\mathcal{E})$ on the hexagonal domain $\max(|m|,|n|,|m+n|)\le R_{\text{hex}} = 8$ (217 cells, 651 nodes) carries $g_1$ on intracell ($\triangle^{\uparrow}$) bonds, $g_2^{(ij)}(\mathbf{s}) = g_2^{(0)} + \tfrac12(s_i+s_j)\Delta g_2$ on intercell ($\triangle^{\downarrow}$) bonds ($\Delta g_2 = g_2^{\text{on}}-g_2^{(0)}$), and node shunts $(y_{\text{base}} + s_i\,y_{\text{shunt,ridge}})/3$, where $s_i\in\{0,1\}$ marks ridge contact. Perimeter cells ($\max = R_{\text{hex}}$; 48 cells, 144 nodes) form $B$ ($V_B=0$ reference); the 169 remaining cells form $I$. Every node carries a positive shunt, so $L_{II}(\mathbf{s})$ is symmetric and strictly diagonally dominant with positive diagonal, hence positive definite (Gershgorin), and~\cite{schur1917,calderon1980}
\begin{equation}
\Lambda^{\text{full}}_{\partial}(\mathbf{s}) = L_{BB}(\mathbf{s}) - L_{BI}(\mathbf{s})[L_{II}(\mathbf{s})]^{-1}L_{IB}(\mathbf{s}).
\end{equation}
Design values: $g_1 = 2.168\times10^{-2}$, $g_2^{(0)} = 10^{-7}$, $g_2^{\text{on}} = 10^{-3}$, $y_{\text{shunt,ridge}} = 2.5\times10^{-3}$ and, in Revision~4.0, $y_{\text{base}} = 10^{-7}$ (all S). Revision~3.9 used $y_{\text{base}} = 10^{-6}$ S; that design is analysed below as the \emph{baseline}. $y_{\text{base}}$ is the leakage admittance from each cell to the reference ($10\,\text{M}\Omega$ per cell at $10^{-7}$ S): a requirement on the fabricated matrix, to be measured, not a derived material constant.

Chord data are line sums of the per-cell ridge signal, $P_\alpha = \sum_{c\in\alpha}(s_c\,y_{\text{shunt,ridge}} + \eta_c)$, over the lines of constant $m$, $n$ and $m+n$ (the three lattice directions, 15 lines each) and, in Revision~4.0, also of constant $m-n$ (29 lines at $30^\circ$ to $\mathbf{a}_1$, which is not a bond direction of the lattice). $\mathbf{C}$ is the corresponding 0/1 incidence matrix, $45\times169$ (three families) or $74\times169$ (four families).

\subsection{Exact Fréchet Derivative of the Schur Complement}
\begin{theorem}[Exact Fréchet Derivative and Quadratic Remainder]\label{thm:frechet}
Let $\mathbf{L}(\mathbf{y}) = \mathbf{L}_0 + \sum_{c\in I}\Delta y_c \mathbf{E}_c$, with $\mathbf{E}_c = \tfrac{1}{3}\sum_{k=1}^3\mathbf{e}_{3c+k}\mathbf{e}_{3c+k}^T$ and $\mathbf{G}_0 = -\mathbf{L}_{II,0}^{-1}\mathbf{L}_{IB,0}$. Then at $\mathbf{y}=\mathbf{0}$:
\begin{enumerate}
\item $D\Lambda_\partial(\mathbf{0})[\Delta\mathbf{y}] = \mathbf{G}_0^T\left(\sum_c \Delta y_c \mathbf{E}_c\right)\mathbf{G}_0$ (positive semidefinite for $\Delta\mathbf{y}\ge\mathbf{0}$).
\item $(b,b)$ diagonal entries: $A_{b,c} = \frac{1}{3}\sum_{k=1}^3 \big(G_{0,(3c+k),b}\big)^2$.
\item Remainder $\mathcal{R}_2(\Delta\mathbf{y}) = \Delta\Lambda_\partial - D\Lambda_\partial[\Delta\mathbf{y}] = -\mathbf{G}_0^T \Delta\mathbf{Y}(\mathbf{L}_{II,0}+\Delta\mathbf{Y})^{-1}\Delta\mathbf{Y}\mathbf{G}_0$, and for $\Delta\mathbf{y}\ge\mathbf{0}$, $\|\mathcal{R}_2\|_2 \le \|\mathbf{G}_0\|_2^2 (\Delta y_{\max}/3)^2 / \lambda_{\min}(\mathbf{L}_{II,0})$.
\end{enumerate}
\end{theorem}
\begin{proof}
Substituting the exact resolvent identity $(\mathbf{L}_{II,0}+\Delta\mathbf{Y})^{-1} = \mathbf{L}_{II,0}^{-1} - \mathbf{L}_{II,0}^{-1}\Delta\mathbf{Y}\mathbf{L}_{II,0}^{-1} + \mathbf{L}_{II,0}^{-1}\Delta\mathbf{Y}(\mathbf{L}_{II,0}+\Delta\mathbf{Y})^{-1}\Delta\mathbf{Y}\mathbf{L}_{II,0}^{-1}$ into $\Lambda_\partial = L_{BB} - L_{BI}L_{II}^{-1}L_{IB}$ yields $\Lambda_\partial(\Delta\mathbf{y}) - \Lambda_\partial(\mathbf{0}) = \mathbf{G}_0^T\Delta\mathbf{Y}\mathbf{G}_0 - \mathbf{G}_0^T\Delta\mathbf{Y}(\mathbf{L}_{II,0}+\Delta\mathbf{Y})^{-1}\Delta\mathbf{Y}\mathbf{G}_0$. Since $\|\mathbf{E}_c\|_2 = 1/3$ and the $\mathbf{E}_c$ have disjoint supports, $\|\Delta\mathbf{Y}\|_2 = \Delta y_{\max}/3$; for $\Delta\mathbf{Y}\succeq0$, $\|(\mathbf{L}_{II,0}+\Delta\mathbf{Y})^{-1}\|_2 \le 1/\lambda_{\min}(\mathbf{L}_{II,0})$. Verification (App.~\ref{app:code}): complex-step differentiation~\cite{squire1998} of $\operatorname{diag}\Lambda_\partial$ reproduces $\mathbf{A}_I$ in all 169 columns to within $10^{-9}$ componentwise, and the identity in (3) holds to within $2\times10^{-8}$ at a random admissible $\Delta\mathbf{y}$, for both values of $y_{\text{base}}$.
\end{proof}

\subsection{Augmented Stage-B Operator and its Identifiability}
In Stage B, we form the dimensionless \textbf{augmented perimeter-chord operator} $\widetilde{\mathbf{A}}_K = [\mathbf{A}_K^T, \, w_c\mathbf{C}_K^T]^T$, with 144 boundary rows plus one row per chord, where $\mathbf{A}_K$ holds the columns of $\mathbf{A}_I$ (Theorem~\ref{thm:frechet}(2)) for the candidate cells, $\mathbf{C}_K$ the corresponding columns of $\mathbf{C}$, and $w_c = \max|\Delta\Lambda_\partial|/y_{\text{shunt,ridge}}$ is a magnitude weight (both blocks are already dimensionless). Candidates are the cells whose chords all exceed $\tau = \tfrac12 y_{\text{shunt,ridge}}$. Revision~3.8 estimated the normalized state $\hat{\mathbf{s}}_K = \Delta\hat{\mathbf{y}}_K / y_{\text{shunt,ridge}}$ by
\begin{equation}\label{eq:ridge}
\hat{\mathbf{s}}_K = \left(\widetilde{\mathbf{A}}_K^T \widetilde{\mathbf{A}}_K + \lambda_{\text{reg}}\mathbf{I}_K\right)^{-1} \widetilde{\mathbf{A}}_K^T \begin{bmatrix} \operatorname{diag}\left(\Delta\mathbf{\Lambda}_{\partial,\text{noisy}}^{\text{full}}\right) / y_{\text{shunt,ridge}} \\ w_c \mathbf{P}_{\text{noisy}} / y_{\text{shunt,ridge}} \end{bmatrix},\qquad \lambda_{\text{reg}} = 10^{-10}.
\end{equation}
The maximum-likelihood decoder instead enumerates the binary patterns on the candidates whose chord sums equal the measured (rounded) counts and selects the one whose exact forward signature $\operatorname{diag}(\Lambda_\partial(\mathbf{s})-\Lambda_\partial(\mathbf{0}))$ is nearest to the data.

\begin{theorem}[Chord Gate, Exact Ranks and Identifiability of Stage B]\label{thm:ghost}
Let chord $\alpha$ contain $L_\alpha\le 2R_{\text{hex}}-1$ interior cells with per-cell noise $\eta_i \sim \mathcal{N}(0,\sigma_{\text{rel}}^2 y_{\text{shunt,ridge}}^2)$, and let the boundary-diagonal data carry i.i.d.\ noise of standard deviation $\sigma_\Lambda = \sigma_{\text{rel}}\max|\Delta\Lambda_\partial|$ ($\sigma_{\text{rel}} = 0.02$, i.e.\ 34 dB).
\begin{enumerate}
\item \textbf{Chord gate.} At threshold $\tau$, $p_{\text{chord}}\le Q\big(1/(2\sigma_{\text{rel}}\sqrt{L_\alpha})\big) \le Q(6.455) \le 5.42\times 10^{-11}$, so the family-wise error is at most $2.44\times 10^{-9}$ for 45 chords and $4.01\times 10^{-9}$ for 74 chords (summing the per-line bounds gives $2.48\times10^{-10}$ in both cases).
\item \textbf{Exact ranks.} For both values of $y_{\text{base}}$, $\operatorname{rank}\mathbf{A}_I = 45$. With three chord families $\operatorname{rank}\mathbf{C} = 42$ and $\operatorname{rank}\widetilde{\mathbf{A}}_I = 87$; with four, $\operatorname{rank}\mathbf{C} = 67$ and $\operatorname{rank}\widetilde{\mathbf{A}}_I = 112$. The kernels have dimension 82 and 57, so in neither design do boundary-diagonal and chord data determine a general interior state.
\item \textbf{Three families: ghost flips.} For the 12-cell test ridge of App.~\ref{app:code}, the noise-free gate returns $K=17$ candidates, $\operatorname{rank}\mathbf{C}_K = 13$, and exactly three binary patterns on these candidates (the true one and two ghosts) share all 45 chord sums. $\widetilde{\mathbf{A}}_K$ has full column rank 17 (exactly), but each ghost's boundary signature differs from the true one by only $d = 0.5265\,\sigma_\Lambda$ at $y_{\text{base}} = 10^{-6}$ S ($4.46\,\sigma_\Lambda$ at $10^{-7}$ S). Consequently every decoder errs with probability at least $Q(d/2) = 0.396$ ($0.0128$ at $10^{-7}$ S) on the true pattern or on a ghost.
\item \textbf{Four families: the test ridge is determined.} The four-family gate returns $K = 16$ candidates and $\mathbf{C}_K$ has full column rank 16 (exactly), so the 74 chord sums alone determine the test ridge among the candidates, for either $y_{\text{base}}$.
\item \textbf{Four families: no general uniqueness.} The switching component $\prod_{\mathbf{d}}(1-T^{\mathbf{d}})\,\delta_{(0,0)}$ over the four line directions equals $+1$ on $S_+$ and $-1$ on $S_-$,
    \begin{equation*}
    S_+ = \{(0,0),(1,2),(2,-1),(3,1)\},\qquad S_- = \{(0,1),(1,-1),(2,2),(3,0)\},
    \end{equation*}
    and has zero line sums in every family, so each of its translates yields two binary 4-cell patterns with identical chord data. Placed in isolation at least two rings deep, all 70 such pairs are separated by at least $30.09\,\sigma_\Lambda$ at the Revision-4.0 design point ($44.62\,\sigma_\Lambda$ at $y_{\text{base}} = 10^{-6}$ S). Enclosed by a closed loop of ridge cells (the 36 cells two rings inside the perimeter), the two resulting 40-cell ridges are separated by at most $6.3\times10^{-4}\,\sigma_\Lambda$ over all 44 placements, so every decoder errs with probability at least $0.4998$ on one ridge of each pair.
\end{enumerate}
\end{theorem}
\begin{proof}
(1) Gaussian tail integral and a union bound over the chords; the longest chord has $L_\alpha = 15$.

(2) \emph{Upper bounds.} Column $b$ of $\mathbf{G}_0$ is $-\mathbf{L}_{II,0}^{-1}\mathbf{L}_{IB,0}[:,b]$ and so depends on $b$ only through $\mathbf{L}_{IB,0}[:,b]$, which does not depend on $y_{\text{base}}$. Of the 144 boundary nodes, 75 have no interior neighbour (zero column), and the remaining 69 columns take only 45 distinct values; hence $\mathbf{A}_I$ has at most 45 distinct non-zero rows and $\operatorname{rank}\mathbf{A}_I\le45$. For three families, the vectors $\mathbf{w}^{1} = (\mathbf{1}_m,-\mathbf{1}_n,\mathbf{0})$, $\mathbf{w}^{2} = (\mathbf{1}_m,\mathbf{0},-\mathbf{1}_p)$ and $\mathbf{w}^{3} = (-m,-n,p)$, indexed by line, are linearly independent and satisfy $\mathbf{w}^T\mathbf{C} = \mathbf{0}$ (every cell lies on one line of each family, and $-m-n+(m+n) = 0$), so $\operatorname{rank}\mathbf{C}\le 42$. For both families the exact rank of the integer matrix $\mathbf{C}$ (42 and 67) is computed by fraction-free Bareiss elimination over $\mathbb{Z}$ with every division checked to be exact. Hence $\operatorname{rank}\widetilde{\mathbf{A}}_I \le 45+\operatorname{rank}\mathbf{C}$ ($w_c\ne0$ does not affect rank). \emph{Lower bounds.} All conductances are exact decimals, so $3\times10^7\,\mathbf{L}_0$ is an integer matrix for both values of $y_{\text{base}}$. Modulo the prime $p = 2^{31}-1$, $\mathbf{L}_{II,0}$ is invertible, so $\mathbf{G}_0$, $\mathbf{A}_I$ and $\mathbf{C}$ reduce consistently, and Gauss--Jordan elimination over $\mathbb{F}_p$ gives ranks 45, $\operatorname{rank}\mathbf{C}$ and $45+\operatorname{rank}\mathbf{C}$. A minor that is non-zero mod $p$ is non-zero over $\mathbb{Q}$, so these lower bounds meet the upper bounds.

(3) $K$ and the three patterns follow from enumerating all $2^{17}$ binary vectors on the candidates. $\operatorname{rank}\mathbf{C}_K = 13$ is certified mod $p$ from below and from above by four linearly independent vectors in $\ker\mathbf{C}_K$, the hexagonal rings with alternating signs around the cells $(a,0)$, $a = -2,\dots,1$; $\operatorname{rank}\widetilde{\mathbf{A}}_K = 17$ is certified mod $p$ as in (2). The boundary signatures are evaluated without cancellation as $\operatorname{diag}(\mathbf{G}(\mathbf{s})^T\Delta\mathbf{L}_{II}\mathbf{G}_0)$, valid since $\Delta\mathbf{L}_{IB} = \Delta\mathbf{L}_{BB} = \mathbf{0}$ when no ridge cell touches $B$; at the baseline they agree with an independent 256-bit ball-arithmetic evaluation to six significant digits (supplementary script). The chord data have the same distribution under the true and the ghost patterns, so only the boundary block distinguishes them. For two Gaussian hypotheses with mean separation $d$ and common noise $\sigma_\Lambda$, any test satisfies $P_0(\text{err})+P_1(\text{err}) \ge 2Q(d/2\sigma_\Lambda)$~\cite{tsybakov2009}, hence $\max(P_0(\text{err}),P_1(\text{err}))\ge Q(d/2\sigma_\Lambda)$.

(4) The mod-$p$ rank of $\mathbf{C}_K$ equals its number of columns, 16; a matrix with full column rank has trivial kernel, so no other real (a fortiori binary) vector on the candidates has the same chord sums.

(5) For a finitely supported $f$, $(1-T^{\mathbf{d}})f$ has zero sum along every line of direction $\mathbf{d}$ (the sum telescopes), and the four factors commute, so the product has zero line sums in all four families; expanding gives 16 terms $\pm\delta$, eight of which cancel in pairs. The separations are computed by the listing with the cancellation-free formula (no ridge cell is adjacent to $B$); for the least favourable enclosed placement (largest separation, $6.273\times10^{-4}\,\sigma_\Lambda$) they are confirmed in 256-bit ball arithmetic (supplementary script). The loop's contact shunts ($y_{\text{shunt,ridge}}/3$ per node) and switched bridges hold the enclosed region near the reference potential, so changes inside the loop barely reach $B$. The two-point bound of (3) then gives the error bound.
\end{proof}

\begin{remark}[Operating point, linearisation and Monte Carlo]\label{rem:mc}
Theorem~\ref{thm:frechet} is accurate for $\Delta y_{\max}\ll\lambda_{\min}(\mathbf{L}_{II,0})$, which is $3.39\times10^{-7}$ S at the baseline and $3.89\times10^{-8}$ S at $y_{\text{base}} = 10^{-7}$ S. At the Stage-B operating point $\Delta y = y_{\text{shunt,ridge}} = 2.5\times10^{-3}$ S the remainder bound evaluates to $0.13$ S and $16.1$ S respectively, against boundary signals of $5.90\times10^{-14}$ S and $1.65\times10^{-10}$ S, and the forward model also switches the bridges $g_2$, which Theorem~\ref{thm:frechet} does not cover. The first-order prediction $\mathbf{A}_I\mathbf{s}\,y_{\text{shunt,ridge}}$ therefore misses the data by a relative $21.9$ (baseline) and $1115$ ($10^{-7}$ S). Estimator~\eqref{eq:ridge} returns $|\hat s|\le 2.3\times10^{-10}$ at the baseline ($\le 1.5\times10^{-3}$ at $10^{-7}$ S) and recovers the test ridge in 0 of 2000 trials in every configuration of Table~\ref{tab:configs}. The maximum-likelihood decoder recovers it in 1239/2000 trials at the baseline (large-sample value $1-Q(0.263) = 0.604$, since the two ghosts nearly coincide), 1973/2000 with three families at $10^{-7}$ S, and 2000/2000 with four families. Only diagonal Schur entries are used; at the baseline, the full $144\times144$ Schur matrix raises the exact rank of the linearised Schur block from 45 to 169, but raises the ghost separation only to $0.794\,\sigma_\Lambda$ (decoder error still $\ge 0.345$).
\end{remark}

\begin{table}[h!]\centering\small
\begin{tabular}{@{}lcccc@{}}
\toprule
Configuration & $\operatorname{rank}\widetilde{\mathbf{A}}_I$ & $K$ & chord-consistent patterns & ML exact / 2000 \\
\midrule
Baseline (Rev.~3.9): $y_{\text{base}} = 10^{-6}$ S, 3 families & 87 & 17 & 3 & 1239 \\
$y_{\text{base}} = 10^{-7}$ S, 3 families & 87 & 17 & 3 & 1973 \\
$y_{\text{base}} = 10^{-6}$ S, 4 families & 112 & 16 & 1 & 2000 \\
Revision 4.0: $y_{\text{base}} = 10^{-7}$ S, 4 families & 112 & 16 & 1 & 2000 \\
\bottomrule
\end{tabular}
\caption{Stage-B identifiability of the 12-cell test ridge for the four combinations of the two design changes (2000 trials at 34 dB; App.~\ref{app:code}).}\label{tab:configs}
\end{table}

\begin{remark}[What each design change does]\label{rem:design}
For the test ridge the fourth chord family alone restores exact recovery; lowering $y_{\text{base}}$ alone raises maximum-likelihood recovery from 1239 to 1973 of 2000 by making the three-family ghosts visible at the boundary ($0.53\to4.46\,\sigma_\Lambda$). The benefit of a lower $y_{\text{base}}$ is pattern-dependent: it lowers the smallest isolated four-family ghost separation from $44.62$ to $30.09\,\sigma_\Lambda$, and neither change helps inside closed ridge loops (Theorem~\ref{thm:ghost}(5)). Revision~4.0 adopts both changes; the fourth family needs its own electrode routing because the $m-n$ lines do not follow lattice bonds.
\end{remark}

\subsection{Nonlinear Monotonicity, Data Collapse and Loop Screening}\label{sec:nonlin}
Throughout this subsection $\mathbf{s}\in[0,1]^{N_I}$ is a state of the interior cells and $s_B = \mathbf{0}$. For $b\in B$ let $\mathbf{w}_b(\mathbf{s})$ be the harmonic extension of $\mathbf{e}_b$: $\mathbf{w}_b = \mathbf{G}(\mathbf{s})\mathbf{e}_b$ on $I$, with $\mathbf{G}(\mathbf{s}) = -\mathbf{L}_{II}(\mathbf{s})^{-1}\mathbf{L}_{IB}(\mathbf{s})$, and $\mathbf{w}_b = \mathbf{e}_b$ on $B$. A boundary node $b$ is \emph{active} if it has an interior neighbour ($\mathbf{L}_{IB}[:,b]\neq\mathbf{0}$); 69 of the 144 boundary nodes are active and 75 are inactive. With $\mathbf{E}_c$ as in Theorem~\ref{thm:frechet} and $\mathbf{a}_e = \mathbf{e}_u-\mathbf{e}_v$ for a bond $e = (u,v)$,
\begin{equation}\label{eq:Mc}
\mathbf{M}_c = \frac{\partial\mathbf{L}}{\partial s_c} = y_{\text{shunt,ridge}}\,\mathbf{E}_c + \frac{\Delta g_2}{2}\sum_{e\in\mathcal{E}_2(c)}\mathbf{a}_e\mathbf{a}_e^T \succeq 0,
\end{equation}
where $\mathcal{E}_2(c)$ is the set of intercell bonds with an end in cell $c$. The 12-cell test ridge of App.~\ref{app:code} is denoted $\mathbf{s}^\star$, and $\mathbf{J}(\mathbf{s}) = \partial\operatorname{diag}\Lambda_\partial(\mathbf{s})/\partial\mathbf{s}\in\mathbb{R}^{144\times169}$ is the nonlinear boundary Jacobian.

\begin{lemma}[Monotonicity of the diagonal boundary response]\label{lem:monotone}
For every $\mathbf{s}\in[0,1]^{N_I}$, every interior cell $c$ and every $b\in B$,
\begin{equation}\label{eq:mono}
J_{bc}(\mathbf{s}) = \frac{\partial\Lambda_{\partial,bb}(\mathbf{s})}{\partial s_c} = \mathbf{w}_b(\mathbf{s})^T\mathbf{M}_c\,\mathbf{w}_b(\mathbf{s}).
\end{equation}
\begin{enumerate}
\item If $b$ is active, $J_{bc}(\mathbf{s}) \ge \tfrac13\,y_{\text{shunt,ridge}}\sum_{k=1}^3 G(\mathbf{s})_{3c+k,b}^2 > 0$; if $b$ is inactive, $J_{bc}(\mathbf{s}) = 0$.
\item If $\mathbf{s}\le\mathbf{s}'$ componentwise and $\mathbf{s}\neq\mathbf{s}'$, then $\Lambda_{\partial,bb}(\mathbf{s}') > \Lambda_{\partial,bb}(\mathbf{s})$ at all 69 active nodes. Distinct comparable states, in particular two states that differ in one cell, never have the same boundary-diagonal data.
\end{enumerate}
At the Revision-4.0 design point the minimum of $J_{bc}$ over the active nodes and all 169 cells is $1.633\times10^{-16}$ S at $\mathbf{s} = \mathbf{0}$, $1.908\times10^{-21}$ S at $\mathbf{s}^\star$ and $5.624\times10^{-34}$ S over 12 test states (baseline: $6.280\times10^{-33}$, $3.919\times10^{-34}$ and $6.631\times10^{-35}$ S).
\end{lemma}
\begin{proof}
Write $\Lambda_\partial = \mathbf{W}^T\mathbf{L}\mathbf{W}$ with $\mathbf{W} = [\mathbf{G};\mathbf{I}]$. Since $\mathbf{L}_{II}\mathbf{G}+\mathbf{L}_{IB} = \mathbf{0}$, $\mathbf{L}\mathbf{W} = [\mathbf{0};\Lambda_\partial]$, and $\partial\mathbf{W}/\partial s_c = [\partial\mathbf{G}/\partial s_c;\mathbf{0}]$; the two terms of $\partial\Lambda_\partial/\partial s_c$ that contain $\partial\mathbf{W}/\partial s_c$ therefore vanish, and $\partial\Lambda_\partial/\partial s_c = \mathbf{W}^T\mathbf{M}_c\mathbf{W}$, whose $(b,b)$ entry is~\eqref{eq:mono}. When cell $c$ is bonded to $B$, $\mathbf{M}_c$ also has entries in the $IB$ and $BB$ blocks, so the restriction $\mathbf{G}^T(\partial\mathbf{L}_{II}/\partial s_c)\mathbf{G}$ agrees with~\eqref{eq:mono} only for cells that are not bonded to $b$.

(1) $\mathbf{L}_{II}(\mathbf{s})$ has non-positive off-diagonal entries and is strictly diagonally dominant (every node carries a positive shunt), so it is a non-singular M-matrix; it is irreducible because the interior nodes form a connected graph. Hence $\mathbf{L}_{II}(\mathbf{s})^{-1}$ is entrywise positive~\cite{berman1994}. For active $b$ the column $-\mathbf{L}_{IB}(\mathbf{s})[:,b]$ is non-negative and non-zero, so $\mathbf{G}(\mathbf{s})[:,b] > \mathbf{0}$ entrywise, and dropping the bond terms of~\eqref{eq:Mc} gives the bound. For inactive $b$, $\mathbf{G}(\mathbf{s})[:,b] = \mathbf{0}$ and $\mathbf{w}_b = \mathbf{e}_b$, and no bond of cell $c$ ends at $b$, so $\mathbf{e}_b^T\mathbf{M}_c\mathbf{e}_b = 0$.

(2) Along $\mathbf{s}(t) = \mathbf{s}+t(\mathbf{s}'-\mathbf{s})\in[0,1]^{N_I}$, $\frac{d}{dt}\Lambda_{\partial,bb}(\mathbf{s}(t)) = \sum_c (s'_c-s_c)\,J_{bc}(\mathbf{s}(t))$ is a sum of non-negative terms of which at least one is positive by (1); integrate over $t\in[0,1]$.

Verification (App.~\ref{app:code}): the sign pattern, diagonal dominance, connectivity and $\mathbf{G}[:,\text{active}] > \mathbf{0}$ at 12 states ($\mathbf{0}$, $\mathbf{s}^\star$, six uniform random states and four random binary states) for both designs; \eqref{eq:mono} against complex-step derivatives in 24 columns at two states (relative error $\le 10^{-9}$); and, for (2), ten pairs of states that differ in one cell, with the integrand positive at six points of every segment and all 381 resolvable boundary increases positive.
\end{proof}
These minima lie many orders of magnitude below the noise level $\sigma_\Lambda$ of Theorem~\ref{thm:ghost}, so Lemma~\ref{lem:monotone}(2) is a statement about exact data, not about robustness to noise. Without strictness, part (1) is Rayleigh's monotonicity law~\cite{doyle1984}; strictness at the active nodes comes from the positivity of $\mathbf{L}_{II}(\mathbf{s})^{-1}$. It does not extend to incomparable states, which Theorem~\ref{thm:ghost}(3,5) already shows can be confused.

\begin{theorem}[Data collapse and rank of the nonlinear boundary Jacobian]\label{thm:collapse}
\begin{enumerate}
\item The 69 active nodes form 24 pairs and 21 singletons. The two nodes $b,b'$ of a pair are bonded to the same single interior node, and $\Lambda_{\partial,bb}(\mathbf{s}) = \Lambda_{\partial,b'b'}(\mathbf{s})$ for every $\mathbf{s}\in[0,1]^{N_I}$; at the 75 inactive nodes $\Lambda_{\partial,bb}$ does not depend on $\mathbf{s}$. Hence $\operatorname{rank}\mathbf{J}(\mathbf{s})\le 45$ at every state, for both designs.
\item $\operatorname{rank}\mathbf{J}(\mathbf{s}^\star) = 45$ (exact).
\item On the gate candidates of $\mathbf{s}^\star$: with four families $K = 16$ and $\operatorname{rank}\mathbf{C}_K = \operatorname{rank}\mathbf{J}(\mathbf{s}^\star)_K = \operatorname{rank}[\mathbf{J}(\mathbf{s}^\star)_K;\,w_c\mathbf{C}_K] = 16$ (ratio of extreme singular values $7.21\times10^{-3}$); with three families $K = 17$, $\operatorname{rank}\mathbf{C}_K = 13$ and $\operatorname{rank}\mathbf{J}(\mathbf{s}^\star)_K = \operatorname{rank}[\mathbf{J}(\mathbf{s}^\star)_K;\,w_c\mathbf{C}_K] = 17$ (ratio $4.92\times10^{-8}$); all ranks exact.
\end{enumerate}
\end{theorem}
\begin{proof}
(1) Every node has its two intracell bonds inside its own cell and its intercell bonds inside a single $\triangle^{\downarrow}$ triangle, so the interior neighbours of $b\in B$ are the interior vertices of its $\triangle^{\downarrow}$ triangle. If boundary nodes $b\neq b'$ have the same non-zero column of $\mathbf{L}_{IB}$, they have the same interior neighbours. Two interior neighbours $u,u'$ would force $\{u,u',b\} = \{u,u',b'\}$, so the common neighbour is a single node $u$ and the triangle is $\{u,b,b'\}$. For every $\mathbf{s}$ with $s_B = \mathbf{0}$, $b$ and $b'$ then carry identical bonds: $g_1$ twice inside their own cells, $g_2^{(0)}+\tfrac12 s_{c(u)}\Delta g_2$ to $u$, $g_2^{(0)}$ to each other, and the shunt $y_{\text{base}}/3$. So $\mathbf{L}_{IB}(\mathbf{s})[:,b] = \mathbf{L}_{IB}(\mathbf{s})[:,b']$ and $L_{BB}(\mathbf{s})_{bb} = L_{BB}(\mathbf{s})_{b'b'}$, and $\Lambda_{\partial,bb} = L_{BB,bb} - \mathbf{L}_{IB}[:,b]^T\mathbf{L}_{II}^{-1}\mathbf{L}_{IB}[:,b]$ takes the same value at $b$ and $b'$. At an inactive node $\mathbf{L}_{IB}[:,b] = \mathbf{0}$, and $L_{BB,bb}$ involves only bonds whose ends both have $s = 0$. Thus $\operatorname{diag}\Lambda_\partial(\mathbf{s})$ is determined by 45 functions of $\mathbf{s}$, and $\mathbf{J}(\mathbf{s})$ has at most 45 distinct non-zero rows. The listing counts the pairs, checks that each pair column has a single non-zero entry, and confirms the identities numerically on random ridges.

(2) At the binary state $\mathbf{s}^\star$ all conductances are exact decimals, so an integer multiple of $\mathbf{L}(\mathbf{s}^\star)$ is an integer matrix; the listing computes $\mathbf{W}(\mathbf{s}^\star)$ and the entries~\eqref{eq:mono} modulo $p = 2^{31}-1$, and Gauss--Jordan elimination over $\mathbb{F}_p$ gives rank 45, which meets the bound of (1).

(3) Lower bounds as in (2); the upper bound is $K$.
\end{proof}
Part (3) gives local injectivity of the candidate-restricted map near $\mathbf{s}^\star$, not uniqueness among binary states: with three families the two binary ghosts of Theorem~\ref{thm:ghost}(3) share the chord data of $\mathbf{s}^\star$ and are separated from it at the boundary by only $0.5265\,\sigma_\Lambda$ (baseline) and $4.46\,\sigma_\Lambda$ (design). With four families $\mathbf{C}_K$ alone already has full column rank, so the boundary block adds no rank.

\begin{theorem}[Screening by closed ridge loops]\label{thm:screen}
Let $\mathbf{z} = \prod_{\mathbf{d}\in D}(1-T^{\mathbf{d}})\,\delta_{(0,0)}$, $D = \{(0,1),(1,0),(1,-1),(1,1)\}$, be the four-family switching component of Theorem~\ref{thm:ghost}(5), and let $\mathcal{R}_r$ be the closed hexagonal ring of the $6r$ cells at radius $r$, at the Revision-4.0 design point.
\begin{enumerate}
\item $\mathbf{z}$ has 8 non-zero cells ($+1$ on $S_+$, $-1$ on $S_-$), zero line sums in all four families, and $\sum_{(m,n)} z_{(m,n)}\,m^a n^b = 0$ for $a+b\le3$; a moment of order 4 is non-zero.
\item For $r = 6$ (36 cells) and each of the 44 translates of $\mathbf{z}$ inside $\mathcal{R}_6$, the ridges $\mathcal{R}_6\cup S_+$ and $\mathcal{R}_6\cup S_-$ have identical chord data, and their boundary separations lie between $1.4445\times10^{-8}$ and $6.2732\times10^{-4}\,\sigma_\Lambda$ ($1.13\times10^{-17}$ and $4.91\times10^{-13}$ S). Every decoder therefore errs with probability at least $0.4998$ on one ridge of each pair. An error probability of at most $1\%$ needs $d\ge 2Q^{-1}(0.01)\,\sigma_\Lambda = 4.6527\,\sigma_\Lambda$, that is an SNR of at least $111.3$ dB for the most separated pair and $204.1$ dB for the least separated one, against 33.98 dB in this design.
\item For $r = 6,5,4,3$ (44, 24, 10 and 2 placements) the median separation is $2.34\times10^{-13}$, $2.83\times10^{-14}$, $6.58\times10^{-15}$ and $2.29\times10^{-15}$ S: it falls monotonically as the loop shrinks, by $0.67$ decades per ring. In noise units the medians are $2.99\times10^{-4}$, $2.83\times10^{-4}$, $3.64\times10^{-4}$ and $6.26\times10^{-4}\,\sigma_\Lambda$, which is not monotone. For every $r$ the loop reduces the separation of the same pair in isolation by a median factor between $1.0\times10^{-4}$ and $1.3\times10^{-4}$.
\item Inside $\mathcal{R}_6$ the separation collapses with the depth of the pair below the loop ($6$ minus the largest radius of its cells): the medians are $3.70\times10^{-4}$, $2.99\times10^{-4}$, $7.66\times10^{-8}$ and $1.44\times10^{-8}\,\sigma_\Lambda$ at depths 1, 2, 3 and 4 (20, 14, 8 and 2 placements).
\end{enumerate}
\end{theorem}
\begin{proof}
(1) Zero line sums are Theorem~\ref{thm:ghost}(5). Of the 16 terms of the expanded product, the identities $(0,1)+(1,0) = (1,1)$ and $(0,1)+(1,-1) = (1,0)$ pair eight with opposite signs, which cancel. For a polynomial $q$, $\sum_{\mathbf{x}} q(\mathbf{x})\,[(1-T^{\mathbf{d}})f](\mathbf{x}) = \sum_{\mathbf{x}}[(1-T^{-\mathbf{d}})q](\mathbf{x})\,f(\mathbf{x})$, and each difference lowers the degree of $q$ by at least one, so the four differences annihilate every $q$ of degree at most 3. The listing evaluates all moments of order at most 4.

(2)--(4) The separations are computed with the cancellation-free formula of Theorem~\ref{thm:ghost}(3), which applies because no ridge cell is adjacent to $B$ ($r\le R_{\text{hex}}-2$). Both extremes of (2) are enclosed in 256-bit ball arithmetic, directly from differences of Schur complements (App.~\ref{app:supp}); the enclosures agree with the listing to four digits. The error bound is the two-point bound of Theorem~\ref{thm:ghost}(3). Since $\sigma_\Lambda = \sigma_{\text{rel}}\max|\Delta\Lambda_\partial|$ and $\text{SNR} = -20\log_{10}\sigma_{\text{rel}}$, a separation of $d_0\,\sigma_\Lambda$ at $33.98$ dB becomes $d_0\,10^{(\text{SNR}-33.98)/20}\,\sigma_\Lambda$, which gives the thresholds. In noise units the separation need not fall when the absolute separation does, because $\sigma_\Lambda$ is proportional to the ridge's own largest boundary signal, which also shrinks as the loop moves inward.
\end{proof}

\begin{remark}[Ridge-class studies]\label{rem:classes}
Two seeded ridge classes are decoded at the Revision-4.0 design point (34 dB), 100 ridges per class and 200 trials per ridge. \emph{Strokes} (seed 2026): two or three strokes, each starting at a uniformly chosen cell of radius $\le5$ in a uniformly chosen lattice direction and visiting 4--8 cells, turning by $\pm60^\circ$ with probability $0.25$ per step (cells beyond radius 5 are skipped); ridges may branch, cross and close loops. \emph{Connected loop-free} (seed 1205), our reading of the Revision-5.0 class: start at a uniformly chosen cell of radius $\le5$ and repeatedly add a uniformly chosen cell of radius $\le5$ that is adjacent to exactly one ridge cell, until the ridge has 12 cells, so that its cell adjacency graph is a tree. For each ridge the four-family gate is applied, all binary patterns on the candidates with the ridge's 74 chord sums are enumerated exhaustively (depth-first search with line-sum propagation), and in each trial maximum likelihood selects among them; a trial counts as a success only if the noisy chords reproduce the candidate set and the chord counts and the decoder returns the true ridge.
\begin{center}\small
\begin{tabular}{@{}lcccc@{}}
\toprule
Class & unique chord data & non-unique rate (95\% CI) & ML exact & error rate (95\% CI) \\
\midrule
Strokes & 98/100 & $[0.243,\,7.04]\%$ & 20000/20000 & $[0,\,0.0185]\%$ \\
Connected loop-free, 12 cells & 95/100 & $[1.64,\,11.3]\%$ & 19978/20000 & $[0.0689,\,0.167]\%$ \\
\bottomrule
\end{tabular}
\end{center}
(Clopper--Pearson intervals.) The two non-unique stroke ridges have two and three chord-consistent patterns, separated by $10.687$, $13.443$ and $39.609\,\sigma_\Lambda$, and are decoded in 200/200 trials each. The five non-unique loop-free ridges have separations $26.044$, $4.742$, $2.376$, $67.632$ and $49.173\,\sigma_\Lambda$; all 22 errors occur on the two least separated (3 and 19 of 200, against $200\,Q(d/2) = 1.8$ and $23.5$), and none of their chord-consistent alternatives is itself connected and loop-free. Revision~5.0 reported $99/100$ and $19{,}886/20{,}000$ for its class without a sampler or code. For the loop-free class above both values lie outside the 95\% intervals; for the stroke class the error rate does.
\end{remark}

\begin{proposition}[Exact ranks for $R_{\text{hex}} = 2,\dots,8$]\label{prop:family}
With the baseline conductances ($y_{\text{base}} = 10^{-6}$ S) the ranks in Table~\ref{tab:family} are exact. Writing $R = R_{\text{hex}}$, for $2\le R\le8$ they follow $N_I = 3R^2-3R+1$, $9R-3$ active nodes in $6R-3$ classes of identical $\mathbf{L}_{IB}$ columns, $\operatorname{rank}\mathbf{A}_I = 6R-3$ ($R\ge3$), $\operatorname{rank}\mathbf{C} = 6R-6$ (three families) and $10R-13$ (four), and $\operatorname{rank}\widetilde{\mathbf{A}}_I = 12R-9$ with three families ($R\ge5$) and $16R-16$ with four ($R\ge6$). For smaller $R$ the augmented operator has full rank $N_I$, except with three families at $R = 4$, where its rank is 35 of 37.
\end{proposition}
\begin{proof}
Upper bounds as in Theorem~\ref{thm:ghost}(2): the number of distinct non-zero columns of $\mathbf{L}_{IB}$, the exact Bareiss ranks of $\mathbf{C}$, and $N_I$. Lower bounds by elimination modulo $p = 2^{31}-1$, and modulo $2^{31}-19$ where the bounds do not meet. They meet everywhere except with three families at $R = 4$, where both primes give 35 against the upper bound 37; exact rational elimination (App.~\ref{app:supp}) gives 35. The closed forms are checked by the listing for every $R\le8$.
\end{proof}
For $R\ge5$ with three families and $R\ge6$ with four, $\operatorname{rank}\widetilde{\mathbf{A}}_I = \operatorname{rank}\mathbf{A}_I+\operatorname{rank}\mathbf{C}$, so the boundary-diagonal and chord row spaces intersect trivially; at $R = 4$ with three families they share four dimensions. The identifiable dimension grows linearly in $R_{\text{hex}}$ and the number of interior cells quadratically; whether the closed forms persist for $R_{\text{hex}}>8$ is open.

\begin{conjecture}[Revision 5.0: connected loop-free ridges]\label{conj:tree}
If the ridges of two binary states $\mathbf{s},\mathbf{s}'$ are connected and loop-free with at most $R_{\text{hex}}+4$ cells, $\mathbf{C}\mathbf{s} = \mathbf{C}\mathbf{s}'$ (four families) and $\operatorname{diag}\Lambda_\partial(\mathbf{s}) = \operatorname{diag}\Lambda_\partial(\mathbf{s}')$, then $\mathbf{s} = \mathbf{s}'$.
\end{conjecture}
The sample of Remark~\ref{rem:classes} contains no two distinct connected loop-free ridges with the same chord data, so it does not test the boundary part of this conjecture.

\begin{table}[h!]\centering\small
\begin{tabular}{@{}rrrrrrrrr@{}}
\toprule
$R_{\text{hex}}$ & $N_I$ & active & classes & $\operatorname{rank}\mathbf{A}_I$ & $\operatorname{rank}\mathbf{C}$ (3) & $\operatorname{rank}\mathbf{C}$ (4) & $\operatorname{rank}\widetilde{\mathbf{A}}_I$ (3) & $\operatorname{rank}\widetilde{\mathbf{A}}_I$ (4) \\
\midrule
2 & 7 & 15 & 9 & 7 & 6 & 7 & 7 & 7 \\
3 & 19 & 24 & 15 & 15 & 12 & 17 & 19 & 19 \\
4 & 37 & 33 & 21 & 21 & 18 & 27 & 35 & 37 \\
5 & 61 & 42 & 27 & 27 & 24 & 37 & 51 & 61 \\
6 & 91 & 51 & 33 & 33 & 30 & 47 & 63 & 80 \\
7 & 127 & 60 & 39 & 39 & 36 & 57 & 75 & 96 \\
8 & 169 & 69 & 45 & 45 & 42 & 67 & 87 & 112 \\
\bottomrule
\end{tabular}
\caption{Exact ranks for $R_{\text{hex}} = 2,\dots,8$ ($y_{\text{base}} = 10^{-6}$ S; (3), (4): number of chord families). The listing prints $[35,37]$ for $R_{\text{hex}} = 4$ with three families; the value 35 is from exact rational elimination (App.~\ref{app:supp}).}\label{tab:family}
\end{table}

\subsection{Split Chords and Dilated Switching Components}\label{sec:split}
A chord line can also be read in $k$ interleaved parts. The \emph{mod-$k$ split chords} $\mathbf{C}^{(k)}$ have one row for each line of the four families and each residue $r$ modulo $k$, summing the cells of the line whose position along it ($n$ on lines of constant $m$, $m$ on the other three families) is congruent to $r$; empty rows are dropped. $\mathbf{C}^{(1)} = \mathbf{C}$ has 74 rows, the parity split $\mathbf{C}^{(2)}$ has 146 and the mod-3 split $\mathbf{C}^{(3)}$ has 216, so $\mathbf{C}^{(3)}$ needs three times the chord channels. Let $\mathbf{z}_k = \prod_{\mathbf{d}\in D}(1-T^{k\mathbf{d}})\,\delta_{(0,0)}$ be the $k$-dilation of the switching component of Theorem~\ref{thm:screen} ($+1$ on $kS_+$, $-1$ on $kS_-$), and write $R = R_{\text{hex}}$. Part (1) of the next theorem is an instance of the switching-component theory of discrete tomography~\cite{hajdu2001} (Section~\ref{sec:related}), applied to lines of step $k\mathbf{d}$; the count $T_k(R)$ on the hexagonal domain and the spanning in the computed cases are what is added here.

\begin{theorem}[Kernels of the split chords]\label{thm:splitker}
\begin{enumerate}
\item Every translate of $\mathbf{z}_k$ supported in $I$ lies in $\ker\mathbf{C}^{(k)}$, these translates are linearly independent, and their number is
\begin{equation}\label{eq:Tk}
T_k(R) = 3R^2-(10k+3)R+8k^2+5k+1\quad(R\ge 2k+1),\qquad T_k(R) = 0\quad(R\le 2k).
\end{equation}
Hence $\dim\ker\mathbf{C}^{(k)}\ge T_k(R)$.
\item Equality holds in every case computed, so there the translates span $\ker\mathbf{C}^{(k)}$: for $k = 1,\dots,5$ at $R = 8$ (kernel dimensions 102, 51, 16, 0 and 0) and for $k = 3$ at every $R = 2,\dots,11$. No proof is given here for general $k$ and $R$.
\item If $k\mid k'$ then $\ker\mathbf{C}^{(k')}\subseteq\ker\mathbf{C}^{(k)}$; in particular every translate of $\mathbf{z}_2$ or $\mathbf{z}_3$ lies in the span of the translates of $\mathbf{z}_1$. The kernels for $k = 2$ and $k = 3$ are not nested: each translate of $\mathbf{z}_3$ has 24 non-zero parity sums, and each translate of $\mathbf{z}_2$ has 24 non-zero mod-3 sums. The undilated component has 0, 24 and 24 non-zero sums under $\mathbf{C}^{(1)}$, $\mathbf{C}^{(2)}$ and $\mathbf{C}^{(3)}$.
\end{enumerate}
\end{theorem}
\begin{proof}
(1) The argument is the standard one~\cite{hajdu2001}. A line is invariant under translation by its direction $\mathbf{d}$, and translation by $k\mathbf{d}$ shifts the position along it by $k$, so it maps each mod-$k$ part of the line to itself. Hence $(1-T^{k\mathbf{d}})f$ has zero sum on every part of every line of direction $\mathbf{d}$ (the sum telescopes), and the product over the four directions has zero sums on every row of $\mathbf{C}^{(k)}$. Dilation preserves the identities $(0,1)+(1,0) = (1,1)$ and $(0,1)+(1,-1) = (1,0)$, so $\mathbf{z}_k$ has the 8 cells $kS_\pm$. In lexicographic order the first cell of $\mathbf{z}_k$ is $(0,0)$, with value $+1$, so the translate by $(a,b)$ has first cell $(a,b)$; distinct translates have distinct first cells and are therefore independent. The cells of $\mathbf{z}_k$ satisfy $0\le m\le 3k$, $-k\le n\le 2k$ and $0\le m+n\le 4k$, with every bound attained, so with $\rho = R-1$ the translate by $(a,b)$ lies in $I$ if and only if $-\rho\le a\le\rho-3k$, $-\rho+k\le b\le\rho-2k$ and $-\rho\le a+b\le\rho-4k$. With $i = a+\rho$, $j = b+\rho-k$ and $n = 2\rho-3k$ these read $0\le i,j\le n$ and $(n+k)/2\le i+j\le(3n-k)/2$, which has no solution for $R\le 2k$. Otherwise, removing from the $(n+1)^2$ square two corner triangles of $t(t+1)/2$ points each, $t = R-1-k$, leaves $T_k = (2R-1-3k)^2-(R-1-k)(R-k)$, which is~\eqref{eq:Tk}.

(2) The exact Bareiss rank of $\mathbf{C}^{(k)}$ equals $N_I-T_k(R)$ in every listed case (App.~\ref{app:code}).

(3) If $k\mid k'$, each mod-$k$ part of a line is a union of mod-$k'$ parts, so the rows of $\mathbf{C}^{(k)}$ are sums of rows of $\mathbf{C}^{(k')}$. The counts are computed by the listing.
\end{proof}

\begin{proposition}[Exact ranks with split chords at $R_{\text{hex}} = 8$]\label{prop:split}
For both values of $y_{\text{base}}$, $\operatorname{rank}\mathbf{C}^{(k)} = 67$, 118 and 153 and $\operatorname{rank}[\mathbf{A}_I;\mathbf{C}^{(k)}] = 112$, 163 and 169 for $k = 1$, 2 and 3. With the mod-3 split on a subset of the families, $\operatorname{rank}[\mathbf{A}_I;\mathbf{C}^{(3)}]$ is 90 for one of $m$, $n$, $m+n$ and 126 for $m-n$; 126 for two of $m$, $n$, $m+n$, 153 for $\{m+n,m-n\}$ and 162 for $\{m,m-n\}$ and $\{n,m-n\}$; 153 for $\{m,n,m+n\}$ and 169 for every three-family subset that contains $m-n$. At $R_{\text{hex}} = 8$ full rank therefore needs the $m-n$ family and two others.
\end{proposition}
\begin{proof}
In every case the mod-$p$ lower bound (two primes) equals $\min(N_I,\,45+\operatorname{rank}\mathbf{C})$, where $\operatorname{rank}\mathbf{C}$ is exact by Bareiss elimination and 45 bounds $\operatorname{rank}\mathbf{A}_I$ (Theorem~\ref{thm:ghost}(2)).
\end{proof}
The boundary block is needed: $\mathbf{C}^{(3)}$ alone leaves the 16 dimensions of Theorem~\ref{thm:splitker}, and $\mathbf{A}_I$ supplies them. Splitting further removes the need for the boundary at this size: $\mathbf{C}^{(4)}$ and $\mathbf{C}^{(5)}$ alone have rank 169.

\begin{theorem}[Closed loops screen the dilated components]\label{thm:splitscreen}
At the Revision-4.0 design point and $R_{\text{hex}} = 8$, two ridges that differ by a translate of $\mathbf{z}_k$ have identical data under $\mathbf{C}^{(k)}$, and:
\begin{enumerate}
\item $k = 2$: the 51 isolated pairs are separated by $75.5$ to $142\,\sigma_\Lambda$. Inside the closed radius-6 ring (36 cells) the 13 pairs that fit are separated by $5.40\times10^{-4}$ to $1.47\times10^{-3}\,\sigma_\Lambda$ ($4.2\times10^{-13}$ to $1.2\times10^{-12}$ S), so every decoder errs with probability at least $0.4997$ on one ridge of each pair.
\item $k = 3$: the 16 isolated pairs are separated by $99.9$ to $147\,\sigma_\Lambda$. A translate of $\mathbf{z}_3$ fits inside a closed ring only at radius 7, next to the perimeter; the 4 pairs there are separated by $1.45$ to $1.78\,\sigma_\Lambda$ ($2.5\times10^{-5}$ to $3.2\times10^{-5}$ S), so every decoder errs with probability at least $0.1868$ on one ridge of each pair.
\end{enumerate}
\end{theorem}
\begin{proof}
Identical data follow from Theorem~\ref{thm:splitker}(1). The radius-7 ring is bonded to $B$, so the formula of Theorem~\ref{thm:ghost}(3) does not apply; all separations are computed from the exact identity $\Lambda_\partial(\mathbf{s})-\Lambda_\partial(\mathbf{0}) = \mathbf{W}(\mathbf{s})^T(\mathbf{L}(\mathbf{s})-\mathbf{L}(\mathbf{0}))\mathbf{W}(\mathbf{0})$, which follows from $\mathbf{L}(\mathbf{s})\mathbf{W}(\mathbf{s}) = [\mathbf{0};\Lambda_\partial(\mathbf{s})]$ and involves no cancellation; the listing checks it against Theorem~\ref{thm:ghost}(3) on the test ridge. The error bound is the two-point bound of Theorem~\ref{thm:ghost}(3).
\end{proof}

\begin{theorem}[No fixed split order suffices for every array size]\label{thm:splitscale}
Let $d(R)$ be the number of distinct non-zero columns of $\mathbf{L}_{IB}$, so that $\operatorname{rank}\mathbf{A}_I\le d(R)\le18R$. Then
\begin{equation}\label{eq:splitbound}
\operatorname{rank}[\mathbf{A}_I;\mathbf{C}^{(k)}]\le N_I-T_k(R)+d(R),
\end{equation}
so for every fixed $k$ the stacked operator is rank-deficient once $T_k(R)>d(R)$, in particular for all sufficiently large $R$. For $k = 3$: $\operatorname{rank}[\mathbf{A}_I;\mathbf{C}^{(3)}] = N_I$ for $2\le R\le9$; it lies between 268 and 270 of 271 at $R = 10$ and equals 306 of 331 at $R = 11$; and it is below $N_I$ for every $R\ge10$.
\end{theorem}
\begin{proof}
For a basis $\mathbf{K}$ of $\ker\mathbf{C}$, $\operatorname{rank}[\mathbf{A}_I;\mathbf{C}] = \operatorname{rank}\mathbf{C}+\operatorname{rank}(\mathbf{A}_I\mathbf{K})$, and $\operatorname{rank}(\mathbf{A}_I\mathbf{K})\le\operatorname{rank}\mathbf{A}_I\le d(R)$ (Theorem~\ref{thm:ghost}(2)). With $\dim\ker\mathbf{C}^{(k)}\ge T_k(R)$ this gives~\eqref{eq:splitbound}, and $d(R)\le18R$ because $B$ has $18R$ nodes. For $k = 3$ the listing finds $d(R) = 6R-3$ for $2\le R\le15$, so $T_3(R)-d(R) = 3R^2-39R+91>0$ for $10\le R\le15$, while $T_3(R)-18R = 3R^2-51R+88>0$ for $R\ge16$. Full rank for $R\le9$ and the values at $R = 10, 11$ follow from mod-$p$ lower bounds and~\eqref{eq:splitbound}. At $R = 10$, $\mathbf{K}$ consists of the 58 translates, $\operatorname{rank}\mathbf{C}^{(3)} = 213$, and $\mathbf{A}_I\mathbf{K}$ has rank 55 modulo four primes and three singular values at round-off level in double precision, so the rank is most likely 268.
\end{proof}

\begin{corollary}[Critical radius for arbitrary split order $k$]\label{cor:rcrit}
Whenever $d(R) = 6R - 3$, the inequality $T_k(R) > d(R)$ holds for all
\begin{equation}
R_{\text{hex}} \ge R_{\text{crit}}(k) = \left\lfloor \frac{10k + 9 + \sqrt{4k^2 + 120k + 33}}{6} \right\rfloor + 1.
\end{equation}
For small $k$ this gives $R_{\text{crit}}(1)=6$, $R_{\text{crit}}(2)=8$,
$R_{\text{crit}}(3)=10$, $R_{\text{crit}}(4)=13$, $R_{\text{crit}}(5)=15$;
for $k=1$ and $k=3$ these are the onsets of rank deficiency observed in
Table~\ref{tab:family} (four families) and in Theorem~\ref{thm:splitscale}.
The root inside the floor is $2k + \tfrac{13}{2} + O(1/k)$, since
$\sqrt{4k^2 + 120k + 33} = 2k + 30 + O(1/k)$; exactly, $R_{\text{crit}}(k) = 2k + 7$
if and only if $k \ge 58$, so the additive constant is eventually $7$ rather
than the simpler $4$ one might guess.
Unconditionally (using only $d(R) \le 18R$),
$\operatorname{rank}[\mathbf{A}_I;\mathbf{C}^{(k)}] < N_I$ for all
$R_{\text{hex}} \ge \lfloor (10k + 21 + \sqrt{4k^2 + 360k + 429})/6 \rfloor + 1$,
which is $R_{\text{hex}}\ge16$ for $k=3$.
\end{corollary}
\begin{proof}
For $R\ge2k+1$ the condition $T_k(R) > d(R)$ reads $3R^2 - (10k+9)R + (8k^2+5k+4) > 0$,
whose larger root is $(10k+9+\sqrt{4k^2+120k+33})/6 > 2k+1$; since $R$ is an integer,
$R \ge \lfloor \cdot \rfloor + 1$ suffices, and Theorem~\ref{thm:splitscale} then gives the rank deficiency. The unconditional bound replaces
$d(R)$ by $18R$ and follows identically. The asymptotics follow from the
binomial expansion $\sqrt{4k^2+120k+33} = 2k\sqrt{1+30/k+O(1/k^2)} = 2k+30+O(1/k)$.
For the exact statement, $R_{\text{crit}}(k)=2k+7$ means $2k+27\le\sqrt{4k^2+120k+33}<2k+33$; the right inequality always holds and the left one is $12k\ge696$.
\end{proof}
Any linear probe set $\mathbf{P}$ satisfies $\operatorname{rank}[\mathbf{A}_I;\mathbf{P}]\le d(R)+\operatorname{rows}(\mathbf{P})$, so full rank needs at least $N_I-d(R)$ probe rows, quadratically many in $R$, whatever the design. Splitting to $k\ge R/2$ leaves no dilated translate ($T_k = 0$), and at $R = 8$ the chords alone then have full rank ($k = 4, 5$), at the cost of $k$ channels per line. How the separation in noise units of screened pairs behaves for such designs at larger $R$ is open.

\begin{remark}[Conditioning of the split operators]\label{rem:splitcond}
Rank is not resolution. At the design point the 163rd singular value of $[\mathbf{A}_I;\mathbf{C}^{(2)}]$ is $8.373\times10^{-10}$ against $\sigma_1 = 4.668$; at $y_{\text{base}} = 10^{-6}$ S it is below $10^{-14}$, at round-off level. For $[\mathbf{A}_I;\mathbf{C}^{(3)}]$ the condition number is $43{,}172$ at the design point ($43{,}171$ after column equilibration, $2{,}150$ after row equilibration) and $5.24\times10^{7}$ at $y_{\text{base}} = 10^{-6}$ S ($2.91\times10^{5}$ after row equilibration); the boundary block sees the 16 directions of $\ker\mathbf{C}^{(3)}$ with singular values from $3.99\times10^{-2}$ down to $8.80\times10^{-5}$ (down to $7.24\times10^{-8}$ at $10^{-6}$ S). The three-family subsets of Proposition~\ref{prop:split} reach full rank at a price: at the design point their condition numbers are $9.17\times10^{5}$ for $\{m,n,m-n\}$ and $2.78\times10^{7}$ for the other two. These numbers depend on how the dimensionless rows of $\mathbf{A}_I$ are weighted against the 0/1 chord rows, and $\mathbf{A}_I$ is a linearisation that misses the data at the operating point by a factor of up to 1115 (Remark~\ref{rem:mc}). The separations of Theorem~\ref{thm:splitscreen}, computed from the nonlinear model in noise units, are the relevant measure.
\end{remark}

\begin{remark}[Standard versus mod-3 chords on the same ridges]\label{rem:splitstudy}
For the two ridge classes of Remark~\ref{rem:classes}, the 100 ridges of each class are drawn here with a generator of their own (same seeds), so that both chord designs see the same ridges; 200 trials per ridge:
\begin{center}\small
\begin{tabular}{@{}lcccc@{}}
\toprule
& \multicolumn{2}{c}{standard chords (74 rows)} & \multicolumn{2}{c}{mod-3 split chords (216 rows)} \\
Class & unique & ML exact & unique & ML exact \\
\midrule
Strokes & 100/100 & 20000/20000 & 100/100 & 20000/20000 \\
Connected loop-free, 12 cells & 97/100 & 19928/20000 & 100/100 & 20000/20000 \\
\bottomrule
\end{tabular}
\end{center}
Under standard chords the three non-unique loop-free ridges have 4, 2 and 2 chord-consistent patterns, at separations $100.66$, $108.2$ and $110.6$, then $0.78$, then $5.18\,\sigma_\Lambda$; the second is decoded in 128/200 trials and the others in 200/200. The 97/100 here and the 95/100 of Remark~\ref{rem:classes} are two samples of the same class. Neither class contains the configurations of Theorem~\ref{thm:splitscreen} (a 42-cell ring around half of a translate of $\mathbf{z}_3$), so perfect decoding of these samples does not contradict the bound $0.1868$ for those pairs.
\end{remark}

\subsection{Convex Ridges and the Cross-Ratio Criterion}\label{sec:convex}
In axial coordinates the four chord directions are $(0,1)$, $(1,0)$, $(1,-1)$ and $(1,1)$. Their slopes $-1$, $0$, $1$, $\infty$ have cross ratio 2, one of the five values excluded by the Gardner--Gritzmann criterion~\cite{gardner1997}, so that theorem does not guarantee that convex ridges are determined by their chords. They are not.

\begin{proposition}[Convex ridges are not determined by the chords]\label{prop:convex}
Let $I_0 = \{(1,0),(1,1),(2,0),(2,1)\}$ and $P_\pm = S_\pm\cup I_0$, with $S_\pm$ as in Theorem~\ref{thm:ghost}(5).
\begin{enumerate}
\item $P_+$ and $P_-$ are convex lattice sets, mirror images of each other under the lattice reflection $(m,n)\mapsto(n,m)$ up to translation, with identical chord data in all four families; $P_+-P_-$ is the switching component $\mathbf{z}$. Their translates give 102 such pairs in the $R_{\text{hex}} = 8$ interior. No two ridges with identical chord data in four directions differ in fewer than 8 cells.
\item With three families, the smallest such convex pair found has 4 cells per ridge: a cell with three alternate neighbours, against the same cell with the other three.
\item At the design point the boundary separates the 70 isolated pairs (all cells at radius $\le6$) by $30.08$ to $89.84\,\sigma_\Lambda$. Inside the closed radius-6 ring the 44 pairs are separated by $1.44\times10^{-8}$ to $6.14\times10^{-4}\,\sigma_\Lambda$, so every decoder errs with probability at least $0.4998$ on one ridge of each pair.
\end{enumerate}
\end{proposition}
\begin{proof}
(1) The octagon with vertices $S_+\cup S_-$ is a lattice $U$-polygon in the sense of~\cite{gardner1997}: every line in one of the four directions through a vertex contains exactly one other vertex, of the opposite sign in $\mathbf{z}$. $P_+$ and $P_-$ are the lattice parallelograms spanned by the two sets of alternate vertices, with edge vectors $(2,-1),(1,2)$ and $(2,1),(-1,2)$; each has area 5 and primitive edges, so by Pick's theorem it contains exactly 4 interior lattice points, the points of $I_0$. Hence $P_\pm = \mathbb{Z}^2\cap\operatorname{conv}P_\pm$, and $P_+-P_- = \mathbf{z}$ has zero line sums (Theorem~\ref{thm:ghost}(5)). If $f\ne0$ has zero line sums in $k$ directions, each of the two supporting lines of $\operatorname{conv}(\operatorname{supp}f)$ parallel to a direction contains at least two points of the support, so the hull has two edges parallel to each direction and the support has at least $2k$ points. The listing verifies convexity, equal chord data and the placements, and finds (2) by searching all convex sets generated by up to four of the 37 cells within radius 3. (3) As in Theorem~\ref{thm:screen}.
\end{proof}
Theorem~\ref{thm:ghost}(4) is therefore a statement about the test ridge. A compact convex contact is in general not determined by its chords; the boundary data separate it from its mirror image by at least $30\,\sigma_\Lambda$ in isolation, but not when a closed ridge loop surrounds both. The four interior cells hardly matter at the boundary: the smallest isolated separation is $30.08\,\sigma_\Lambda$ against $30.09\,\sigma_\Lambda$ for the bare pair $S_\pm$ of Theorem~\ref{thm:ghost}(5), and the largest enclosed separation $6.14\times10^{-4}\,\sigma_\Lambda$ against $6.27\times10^{-4}\,\sigma_\Lambda$ (Theorem~\ref{thm:screen}).

% ============================================================================
\clearpage
\section{Design Targets and Numerical Verification}
% ============================================================================
Table~\ref{tab:specs} lists each verified specification with the result that establishes it.

\begin{table}[!p]\centering\footnotesize
\begin{tabularx}{\textwidth}{@{}>{\raggedright\arraybackslash}p{3.6cm}>{\raggedright\arraybackslash}X>{\raggedright\arraybackslash}p{2.6cm}@{}}
\toprule
\textbf{Quantity} & \textbf{Value} & \textbf{Source} \\
\midrule
Effective density $\rho_{\text{eff}}$ & 680.7 kg/m³ & 5-phase closure \\
Pitches $a_1/a_2$ & 2000/50 µm & $N=2$ hierarchy \\
Breathing asymmetry $\delta_b$ & 0.180 ($t_1 = 1.18$, $t_2 = 0.82$) & Thm.~\ref{thm:bloch} \\
Bulk gap at $\mathbf{K}$ & $3|t_1-t_2| = 1.08$ & Thm.~\ref{thm:bloch} \\
Phase at $\delta_b=+0.18$ & trivial: 217/217 states below gap, 0 in gap & Rem.~\ref{rem:phase} \\
$\Gamma$-degeneracy & $E^{(1)}_\Gamma = E^{(2)}_\Gamma = t_1+t_2$ & Thm.~\ref{thm:cls}(i) \\
CLS residual & $2.7\times10^{-17}$ & Thm.~\ref{thm:cls}(ii) \\
Composite permittivity & $\varepsilon_{\text{eff}} = 5.58$ & Maxwell--Garnett \\
Piezo $g_{33,\text{eff}}$ & 40.5 mV·m/N & $\phi_p=0.2$, $d_{33}=2$ pC/N \\
Wake / threshold & 70.2 / 50.0 mV & $F_z=1.5$ N, $h=4\,\mu$m \\
Intracell phase offsets & $211.22^\circ$, $334.44^\circ$ & Thm.~\ref{thm:jac}(1) \\
$\max\|\mathbf{E}_0^{-1}\Delta\mathbf{E}\|_2$ & $0.55089_\uparrow$, $0.52212_\downarrow$ (first-order bound $0.800_\uparrow$, $0.556_\downarrow$) & Thm.~\ref{thm:jac}(2,3) \\
Orientation margin & $\min\det(\mathbf{I}+\mathbf{E}_0^{-1}\Delta\mathbf{E}) = 0.57854_\uparrow$, $0.69278_\downarrow$ over 65{,}879 + 66{,}140 triangles & Thm.~\ref{thm:jac}(3) \\
Design changes (Rev.~4.0) & $y_{\text{base}} = 10^{-7}$ S (was $10^{-6}$); fourth chord family along $m-n$ (29 lines) & §\ref{sec:inv}.1 \\
Fréchet derivative check & complex step, $\le 10^{-9}$ ($y_{\text{base}} = 10^{-6}$ and $10^{-7}$ S, 169 columns) & Thm.~\ref{thm:frechet} \\
Chord FWER & $\le 2.44\times10^{-9}$ (45 chords), $\le 4.01\times10^{-9}$ (74 chords) & Thm.~\ref{thm:ghost}(1) \\
$\operatorname{rank}\widetilde{\mathbf{A}}_I$ & 87 (baseline), 112 (Rev.~4.0) of 169; kernel dimension 82 / 57; exact & Thm.~\ref{thm:ghost}(2) \\
Test ridge & baseline: 3 chord-consistent patterns, separation $0.5265\,\sigma_\Lambda$, any decoder errs w.p.\ $\ge 0.396$; Rev.~4.0: determined by the chords ($\operatorname{rank}\mathbf{C}_K = K = 16$) & Thm.~\ref{thm:ghost}(3,4) \\
Stage-B Monte Carlo (2000 trials, 34 dB) & estimator~\eqref{eq:ridge}: 0/2000 (every design); maximum likelihood: 1239/2000 (baseline), 2000/2000 (Rev.~4.0) & Rem.~\ref{rem:mc}, Table~\ref{tab:configs} \\
Closed ridge loop & two 40-cell ridges with identical chord data, separation $\le 6.3\times10^{-4}\,\sigma_\Lambda$; any decoder errs w.p.\ $\ge 0.4998$ & Thm.~\ref{thm:ghost}(5) \\
Monotonicity & $\partial\Lambda_{\partial,bb}/\partial s_c>0$ at the 69 active nodes for every $\mathbf{s}\in[0,1]^{169}$, $=0$ at the other 75; minimum $1.633\times10^{-16}$ S at $\mathbf{s}=\mathbf{0}$ & Lemma~\ref{lem:monotone} \\
Boundary Jacobian & $\operatorname{rank}\mathbf{J}(\mathbf{s})\le45$ at every state (24 node pairs carry identical data); $=45$ at the test ridge & Thm.~\ref{thm:collapse} \\
Switching component & 8 cells; moments of order $\le3$ vanish & Thm.~\ref{thm:screen}(1) \\
36-cell loop, 44 placements & separation $1.4445\times10^{-8}$ to $6.2732\times10^{-4}\,\sigma_\Lambda$ (both ends in ball arithmetic); error $\le1\%$ needs $\ge111.3$ / $204.1$ dB & Thm.~\ref{thm:screen}(2) \\
Ridge classes (100 ridges, 200 trials each) & strokes: 98/100 unique, 20000/20000; connected loop-free (12 cells): 95/100, 19978/20000 & Rem.~\ref{rem:classes} \\
Exact ranks, $R_{\text{hex}}=2,\dots,8$ & four families $16R_{\text{hex}}-16$ ($R_{\text{hex}}\ge6$); three families $12R_{\text{hex}}-9$ ($R_{\text{hex}}\ge5$); 35 of 37 at $R_{\text{hex}}=4$ & Prop.~\ref{prop:family} \\
Split chords ($R_{\text{hex}}=8$) & $\operatorname{rank}[\mathbf{A}_I;\mathbf{C}^{(k)}] = 112$, 163, 169 for $k=1,2,3$ (74, 146, 216 rows); $\ker\mathbf{C}^{(k)}$ spanned by 102, 51, 16 dilated translates & Thm.~\ref{thm:splitker}, Prop.~\ref{prop:split} \\
Dilated pairs inside loops & $k=2$: $\le1.47\times10^{-3}\,\sigma_\Lambda$, error $\ge0.4997$; $k=3$: $\le1.78\,\sigma_\Lambda$, error $\ge0.1868$ & Thm.~\ref{thm:splitscreen} \\
Array size, $k=3$ & full rank for $R_{\text{hex}}\le9$; deficient for every $R_{\text{hex}}\ge10$ (268--270 of 271 at 10, 306 of 331 at 11) & Thm.~\ref{thm:splitscale} \\
Same ridges, mod-3 chords & 100/100 unique, 20000/20000 for both classes (standard chords: 100/100 and 97/100 unique) & Rem.~\ref{rem:splitstudy} \\
Convex ridges & cross ratio 2; mirror-image 8-cell parallelograms share all chord data (102 placements); boundary separation $\ge30.08\,\sigma_\Lambda$ isolated, $\le6.14\times10^{-4}\,\sigma_\Lambda$ inside the 36-cell loop & Prop.~\ref{prop:convex} \\
\bottomrule
\end{tabularx}
\caption{Verified design specifications (Revision 5.5). Every computed entry is printed by the listings in App.~\ref{app:code}.}\label{tab:specs}
\end{table}

% ============================================================================
\section*{Revision 5.5: scope and critical radius}
\addcontentsline{toc}{section}{Revision 5.5 notes}
\begin{enumerate}[noitemsep,leftmargin=*]
\item Section~\ref{sec:related} opens with a scope paragraph: the results of Section~\ref{sec:inv} are substrate-agnostic, and the switching-component and cross-ratio results are placed within classical discrete tomography, with the contributions of this paper listed separately.
\item New Corollary~\ref{cor:rcrit}: a closed form for the radius $R_{\text{crit}}(k)$ beyond which $\mathbf{C}^{(k)}$ with the boundary block is rank-deficient, checked by the listing. Remarks and propositions after it are renumbered by one.
\item The listings are re-run on a second machine. Only round-off-level values differ between the two runs: the complex-step and remainder checks of Theorem~\ref{thm:frechet} and Lemma~\ref{lem:monotone} and $\sigma_{163}$ at $y_{\text{base}} = 10^{-6}$ S in Remark~\ref{rem:splitcond}. These are now stated as bounds that hold for both runs, and the listing asserts them; no stated result changes.
\end{enumerate}

% ============================================================================
\section*{Revision 5.4: editorial}
\addcontentsline{toc}{section}{Revision 5.4 notes}
\begin{enumerate}[noitemsep,leftmargin=*]
\item The abstract is shortened from about 650 to about 290 words; the per-revision summaries it carried remain in the body and in these notes, and the introduction ends with a section guide instead of a revision list. No result changes, and the listing is re-run.
\item Section~\ref{sec:convex} compares the convex pairs with the bare switching pairs of Theorem~\ref{thm:ghost}(5).
\end{enumerate}

\section*{Revision 5.3: related work and convex ridges}
\addcontentsline{toc}{section}{Revision 5.3 notes}
\begin{enumerate}[noitemsep,leftmargin=*]
\item New: Section~\ref{sec:related} (related work) and Section~\ref{sec:convex} (Proposition~\ref{prop:convex}). Everything else is Revision~5.2 and is re-verified by the listing.
\item Credit: the switching components and their dilations (Theorem~\ref{thm:ghost}(5), Theorem~\ref{thm:splitker}(1)) are instances of the theory of Hajdu and Tijdeman~\cite{hajdu2001,hajdu2001laa}; Theorem~\ref{thm:splitker} now says so and claims only the count $T_k(R)$ and the computed spanning. Lemma~\ref{lem:monotone} is the strict form of Rayleigh's monotonicity law~\cite{doyle1984}.
\item The four chord families have cross ratio 2, which the Gardner--Gritzmann criterion~\cite{gardner1997} excludes. The test asked for in discussion of Revision~5.2 finds that the obstruction is realised at $R_{\text{hex}} = 8$, by two mirror-image convex 8-cell ridges at 102 placements (Proposition~\ref{prop:convex}). Theorem~\ref{thm:ghost}(4), which concerns the test ridge, is unchanged.
\item Withdrawn from discussion drafts: that Theorem~\ref{thm:splitker} is a new characterisation of the kernel, that the loop separations are a rigorous two-sided bound for loops in general (they are computed for the listed loops, with ball-arithmetic enclosures at the extremes only), and that the paper solves a previously unsolved problem.
\end{enumerate}

\section*{Revision 5.2: split chords}
\addcontentsline{toc}{section}{Revision 5.2 notes}
\begin{enumerate}[noitemsep,leftmargin=*]
\item New: Section~\ref{sec:split} (Theorems~\ref{thm:splitker}, \ref{thm:splitscreen} and~\ref{thm:splitscale}, Proposition~\ref{prop:split}, Remarks~\ref{rem:splitcond} and~\ref{rem:splitstudy}). Everything else is Revision~5.1 and is re-verified by the listing.
\item A separately circulated draft, also labelled Revision~5.1, proposed parity and mod-3 split chords (``conjugate-basis probes''). Its exact ranks (112, 163 and 169, and the family-subset table) are confirmed by Proposition~\ref{prop:split}, and its prediction of 100/100 unique and 20{,}000/20{,}000 recoveries with mod-3 chords by Remark~\ref{rem:splitstudy}. Corrected:
  closed loops still screen, now the dilated components (Theorem~\ref{thm:splitscreen}), so splitting does not remove the obstruction;
  ``three families including $m-n$ are necessary and sufficient'' holds at $R_{\text{hex}} = 8$ only, and even all four families lose full rank at $R_{\text{hex}} = 10$ (Theorem~\ref{thm:splitscale});
  the kernel of $\mathbf{C}^{(k)}$ is shown here to be spanned by dilated translates only in the computed cases, while the inclusion and the count $T_k(R)$ hold for every $k$ and $R$;
  the kernels for $k = 2$ and $k = 3$ are not nested;
  the mod-3 obstruction does fit inside the $R_{\text{hex}} = 8$ flake (4 placements inside the outermost ring);
  the condition numbers 3756 and 350 and the ratio $8.9\times10^{-6}$ are not reproduced (Remark~\ref{rem:splitcond}), and the 163rd singular value is not resolvable at any realistic noise level;
  $m-n$ alone (rank 126) is not worth more than the other three families together (153);
  the undilated component has 24 non-zero mod-3 sums, not 40--60;
  whether a rational number vanishes modulo $p$ does not depend on its size, and a mod-$p$ rank is only a lower bound;
  the symbol $\mathbf{C}_3$ for the mod-3 split clashed with the three-family chords and is $\mathbf{C}^{(3)}$ here.
\end{enumerate}

\section*{Revision 5.1: corrections to Revision 5.0 and new results}
\addcontentsline{toc}{section}{Revision 5.1 notes}
\begin{enumerate}[noitemsep,leftmargin=*]
\item New: Section~\ref{sec:nonlin} (Lemma~\ref{lem:monotone}, Theorems~\ref{thm:collapse} and~\ref{thm:screen}, Remark~\ref{rem:classes}, Proposition~\ref{prop:family}, Table~\ref{tab:family}) and the supplementary checks of App.~\ref{app:supp}. Sections~1--5.3 are those of Revision~4.0 and are re-verified by the listing. Revision~5.0's Lemma~5.2, Theorem~5.3 and Theorem~5.4 correspond to Lemma~\ref{lem:monotone}, Theorem~\ref{thm:collapse} and Theorem~\ref{thm:screen}.
\item Revision~5.0 values that no computation produces, and what replaces them:
  $K = 14$ candidates: the four-family gate gives $K = 16$, and none of 45 family-subset and threshold variants gives 14;
  $\operatorname{rank}[\mathbf{J}(\mathbf{s}^\star)_K;w_c\mathbf{C}_K] = 14/14$: the rank is $16 = K$, and it adds nothing to $\operatorname{rank}\mathbf{C}_K = 16$;
  $\min\partial\Lambda_{\partial,bb}/\partial s_c = 1.054\times10^{-13}$: not reproduced by any of 28 conventions (closest $6.53\times10^{-14}$), replaced by the minima of Lemma~\ref{lem:monotone};
  $d_{\text{screen}} = 3.84\times10^{-6}\,\sigma_\Lambda$ for a 32-cell loop at depth 4--5: no hexagonal ring has 32 cells (rings have $6r$), replaced by Theorem~\ref{thm:screen}(3,4);
  $99/100$ unique ridges and $19{,}886/20{,}000$ recoveries: no sampler or code, replaced by the two seeded studies of Remark~\ref{rem:classes}, neither of which reproduces both values.
\item Revision~5.0 statements corrected:
  the separation inside a loop does not decay in noise units, only in siemens (Theorem~\ref{thm:screen}(3));
  the four-family switching component has 8 cells, not 16 (Revision~4.0 already noted the cancellation);
  $[1.4\times10^{-8},\,6.3\times10^{-4}]\,\sigma_\Lambda$ is the range of the separation over the 44 placements, not a ball-arithmetic enclosure of one quantity; both ends are now enclosed in ball arithmetic (and the ratio was written with $\sigma_\Lambda$ twice);
  ``$\text{SNR}\ge125$ dB for reliable resolution'' has no derivation; error probability $\le1\%$ needs at least 111.3 dB for the most separated placement and 204.1 dB for the least separated;
  the derivative $\mathbf{G}^T(\partial\mathbf{L}_{II}/\partial s_c)\mathbf{G}$ is correct only for cells not bonded to $b$; the general formula uses the harmonic extension (Lemma~\ref{lem:monotone});
  the orientation result keeps the Revision-3.9/4.0 statement (132{,}019 triangles; $0.5222$ for $\triangle^{\downarrow}$, rounded up), not the Revision-3.8 anchor count 130{,}592 and $0.5221$;
  Revision~5.0's appendix reproduced only the exact ranks, not every number of its Sections~2--3.
\item Not adopted from Revision~5.0: the re-framing of the material system around borax--water-activated cellulose-nanofibre LIG and BaTiO$_3$ paper. Its parameter table gives $g_1 = 7.227\times10^{-3}$ S, one third of the $2.168\times10^{-2}$ S used in every computation, and $V_{\text{wake}} = 168.6$ mV is not derived there; Sections~1--4 keep the Revision-4.0 model. Revision~5.0's conjecture on a topological characterisation of identifiable ridges is not carried over, since it does not define identifiability for a single pattern; its conjecture on loop-free ridges is Conjecture~\ref{conj:tree}.
\end{enumerate}

\section*{Revision 4.0: design changes and new results}
\addcontentsline{toc}{section}{Revision 4.0 notes}
\begin{enumerate}[noitemsep,leftmargin=*]
\item Design: $y_{\text{base}}$ lowered from $10^{-6}$ to $10^{-7}$ S, and a fourth chord family along $m-n$ (29 lines) added. Sections~2--4 do not depend on either change; the listing re-verifies all their values.
\item Theorem~\ref{thm:ghost} now covers both designs: exact ranks 87 and 112 (kernel dimensions 82 and 57). With four families the test ridge is determined by its chord sums and recovered in 2000/2000 trials (Table~\ref{tab:configs}). A closed ridge loop screens its interior, so uniqueness does not hold for every ridge (Theorem~\ref{thm:ghost}(5)).
\item Values that change with $y_{\text{base}}$: $\lambda_{\min}(\mathbf{L}_{II,0}) = 3.89\times10^{-8}$ S, the remainder bound at the operating point is 16.1 S, complex-step agreement $2.9\times10^{-10}$, linear-model mismatch 1115 (Remark~\ref{rem:mc}).
\item Bounds are now printed and quoted with safe-direction rounding: $p_{\text{chord}}\le5.42\times10^{-11}$ (Revision~3.9: $5.41\times10^{-11}$), $\max\|\mathbf{E}_0^{-1}\Delta\mathbf{E}\|_2\le0.5222$ for $\triangle^{\downarrow}$ (Revision~3.9 abstract: 0.5221), full-Schur decoder error $\ge0.345$ (Revision~3.9: 0.346), Bloch check $\le 2.3\times10^{-15}$ (Revision~3.9: $2.2\times10^{-15}$).
\item The title now names the three- and four-family identifiability analysis.
\end{enumerate}

\section*{Revision 3.9: corrections to Revision 3.8}
\addcontentsline{toc}{section}{Revision 3.9 notes}
\begin{enumerate}[noitemsep,leftmargin=*]
\item The console output printed in Revision~3.8 was not produced by its listing. Run as printed, the listing stops with an \texttt{AssertionError} in the Theorem~\ref{thm:jac} block; its later print statements hard-code ``[OK]'', ``(100.00\%)'' and ``(ker = \{0\})''.
\item Theorem~\ref{thm:ghost}(2) previously claimed $\operatorname{rank}\widetilde{\mathbf{A}}_I = 169$ and a trivial kernel, so that Stage-B ghost elimination was ``unconditional''. The exact rank is 87, and at 34 dB no decoder can succeed with probability above 0.604 on both the test ridge and its ghosts. The previous Monte Carlo claims (100\% success, $\langle K\rangle = 12$, $\min\sigma_K = 0.1418$, $K\le21$, error $\le0.204$) are replaced by the computed values (0/2000 for estimator~\eqref{eq:ridge}, $K=17$, $\sigma_{\min} = 4.3\times10^{-14}$). The title is changed accordingly.
\item Theorem~\ref{thm:jac}: the claimed analytic bound $2\eta_r/\alpha_b = 0.3902$ for $\triangle^{\uparrow}$ is violated (maximum $0.5509$). The intracell offsets are $211.22^\circ$ and $334.44^\circ$, not ``near $120^\circ$''. The identity $\|\mathbf{E}_0^{-1}\|_2 = \sqrt2/(\ell R/\rho)$ fails by up to $2.7\%$ because $\mathbf{\Phi}_0$ is not conformal. The triangle count was 60{,}552; the anchor rule actually evaluated 130{,}592 triangles, and the complete set of perturbed triangles is 132{,}019. The minimum ratios are 0.9988/0.9993, not 0.612/0.618. ``Ostrowski's determinant bound'' is replaced by the elementary eigenvalue argument.
\item Theorem~\ref{thm:frechet}: the finite-difference check with step $\epsilon_s = 10^{-4}$ gives $2.6\times10^{-2}$, not $3.19\times10^{-7}$. It is replaced by complex-step verification; the validity domain of the linearisation is stated in Remark~\ref{rem:mc}.
\item The HOTI designation is withdrawn for $\delta_b = +0.18$ (Remark~\ref{rem:phase}).
\item Smaller corrections: $\|\mathbf{q}_\varphi\|_2 = 4.5650/a_2$; FWER $\le 2.44\times10^{-9}$; the null space of $\mathbf{C}_{\text{chord}}$ is 127-dimensional (Revision~3.8 said ``124-dimensional cokernel''); off-diagonal Schur entries are not redundant (Remark~\ref{rem:mc}); ``LIGO'' $\to$ LIG; the hard-coded ``Remark 5.4'' cross-reference, the markdown emphasis and the math-mode thousands separators are fixed; the specification table no longer overflows the page.
\end{enumerate}

% ============================================================================
\begin{thebibliography}{99}\small
\bibitem{maxwell1864} J.C. Maxwell, Philos. Mag. 27, 294 (1864).
\bibitem{guest2003} S.D. Guest, J.W. Hutchinson, J. Mech. Phys. Solids 51, 383 (2003).
\bibitem{fleck2010} N.A. Fleck, V.S. Deshpande, M.F. Ashby, Proc. R. Soc. A 466, 2495 (2010).
\bibitem{gibson1997} L.J. Gibson, M.F. Ashby, \textit{Cellular Solids}, 2nd ed., Cambridge, 1997.
\bibitem{ezawa2018} M. Ezawa, Phys. Rev. Lett. 120, 026801 (2018).
\bibitem{kempkes2019} S.N. Kempkes \textit{et al.}, Nat. Mater. 18, 1292 (2019).
\bibitem{bergman2008} D.L. Bergman, C. Wu, L. Balents, Phys. Rev. B 78, 125104 (2008).
\bibitem{lin2014} J. Lin \textit{et al.}, Nat. Commun. 5, 5714 (2014).
\bibitem{danninger2022} D. Danninger \textit{et al.}, Sci. Adv. 8, eadd7118 (2022).
\bibitem{islam2017} M.R. Islam \textit{et al.}, Sci. Rep. 7, 13070 (2017).
\bibitem{yamada1982} T. Yamada, T. Ueda, T. Kitayama, J. Appl. Phys. 53, 4328 (1982).
\bibitem{simmons1963} J.G. Simmons, J. Appl. Phys. 34, 1793 (1963).
\bibitem{zeng2015} X. Zeng \textit{et al.}, ACS Nano 9, 9771 (2015).
\bibitem{schur1917} J. Schur, J. Reine Angew. Math. 147, 205 (1917).
\bibitem{calderon1980} A.P. Calder\'on, in \textit{Seminar on Numerical Analysis}, Rio de Janeiro, 1980, pp.~65--73.
\bibitem{szyld2006} D.B. Szyld, The many proofs of an identity on the norm of oblique projections, Numer. Algorithms 42, 309 (2006).
\bibitem{squire1998} W. Squire, G. Trapp, Using complex variables to estimate derivatives of real functions, SIAM Rev. 40, 110 (1998).
\bibitem{berman1994} A. Berman, R.J. Plemmons, \textit{Nonnegative Matrices in the Mathematical Sciences}, SIAM, Philadelphia, 1994, Ch.~6.
\bibitem{ryser1957} H.J. Ryser, Combinatorial properties of matrices of zeros and ones, Canad. J. Math. 9, 371 (1957).
\bibitem{hajdu2001} L. Hajdu, R. Tijdeman, Algebraic aspects of discrete tomography, J. Reine Angew. Math. 534, 119 (2001).
\bibitem{hajdu2001laa} L. Hajdu, R. Tijdeman, An algorithm for discrete tomography, Linear Algebra Appl. 339, 147 (2001).
\bibitem{stolk2010} A. Stolk, K.J. Batenburg, An algebraic framework for discrete tomography: revealing the structure of dependencies, SIAM J. Discrete Math. 24, 1056 (2010).
\bibitem{gardner1997} R.J. Gardner, P. Gritzmann, Discrete tomography: determination of finite sets by X-rays, Trans. Amer. Math. Soc. 349, 2271 (1997).
\bibitem{matej1998} S. Matej, G.T. Herman, A. Vardi, Binary tomography on the hexagonal grid using Gibbs priors, Int. J. Imaging Syst. Technol. 9, 126 (1998).
\bibitem{nagy2016} B. Nagy, T. Luki\'c, Dense projection tomography on the triangular tiling, Fund. Inform. 145, 125 (2016).
\bibitem{curtis1998} E.B. Curtis, D. Ingerman, J.A. Morrow, Circular planar graphs and resistor networks, Linear Algebra Appl. 283, 115 (1998).
\bibitem{colin1996} Y. Colin de Verdi\`ere, I. Gitler, D. Vertigan, R\'eseaux \'electriques planaires II, Comment. Math. Helv. 71, 144 (1996).
\bibitem{doyle1984} P.G. Doyle, J.L. Snell, \textit{Random Walks and Electric Networks}, Carus Math. Monographs 22, MAA, 1984.
\bibitem{tamburrino2002} A. Tamburrino, G. Rubinacci, A new non-iterative inversion method for electrical resistance tomography, Inverse Problems 18, 1809 (2002).
\bibitem{harrach2013} B. Harrach, M. Ullrich, Monotonicity-based shape reconstruction in electrical impedance tomography, SIAM J. Math. Anal. 45, 3382 (2013).
\bibitem{tsybakov2009} A.B. Tsybakov, \textit{Introduction to Nonparametric Estimation}, Springer, 2009, Ch.~2.
\end{thebibliography}

% ============================================================================
\appendix
\section{Verification Suite (Revision 5.5)}\label{app:code}
% ============================================================================
Requires only \texttt{numpy}; runtime $\approx$220 s on one CPU core. Every computed number in Sections~2--5 and in Tables~\ref{tab:configs}, \ref{tab:family} and~\ref{tab:specs} is printed by this listing, except the ball-arithmetic enclosures and the exact rank at $R_{\text{hex}} = 4$, which are printed by the supplementary script of App.~\ref{app:supp}. Each block raises \texttt{AssertionError} if a stated claim fails. Upper bounds are printed rounded up and lower bounds rounded down. Round-off-level digits (complex-step errors, residuals, singular values at round-off level) and timings depend on the numpy/BLAS build and can differ from the console output below; the text states bounds for them, which the listing asserts.

\begin{lstlisting}[language=Python,style=cleanstyle,caption={Verification suite for Sections~2--5, Theorems~\ref{thm:bloch}--\ref{thm:splitscale}, Corollary~\ref{cor:rcrit}, Propositions~\ref{prop:family}, \ref{prop:split} and~\ref{prop:convex} and Tables~\ref{tab:configs}, \ref{tab:family} and~\ref{tab:specs}.}]
"""Verification suite for 'Hierarchical Bio-Kagome Solid-State Metamaterials' (Revision 5.5).

Every number printed is computed below; no result string is hard-coded.  Only numpy is
required (exact ranks use a self-contained modular-arithmetic elimination).
Each block raises AssertionError if a stated claim fails.
"""
import math, time, itertools
from dataclasses import dataclass, replace
from fractions import Fraction
from decimal import Decimal, ROUND_CEILING, ROUND_FLOOR
import numpy as np


@dataclass(frozen=True)
class Spec:
    delta_b: float = 0.18
    t_sum: float = 2.0
    g1: float = 2.168e-2
    g2_base: float = 1.0e-7
    g2_on: float = 1.0e-3
    y_base: float = 1.0e-7          # Revision 4.0 design value (Revisions <= 3.9: 1.0e-6)
    y_ridge: float = 2.5e-3
    a2_mm: float = 0.050
    a1_mm: float = 2.000
    rx_mm: float = 10.0
    ry_mm: float = 12.5
    rho_c: float = 0.80
    eta_r: float = 0.080
    sigma_rel: float = 0.02
    R_hex: int = 8
    chord_families: int = 4         # Revision 4.0: m, n, m+n, m-n lines (Revisions <= 3.9: 3)


BASELINE = Spec(y_base=1.0e-6, chord_families=3)   # design analysed in Revision 3.9
# cell (m,n) owns nodes 3i+k (k=0,1,2) = up-triangle; down-triangle bonds:
# v2(m,n)-v1(m+1,n), v3(m,n)-v1(m,n+1), v3(m,n)-v2(m-1,n+1)
INTERCELL = (((1, 0), 1, 0), ((0, 1), 2, 0), ((-1, 1), 2, 1))
RIDGES = ((-3, 1), (-2, 1), (-1, 1), (0, 1), (1, 0), (2, 0), (3, -1),
          (-3, 0), (-2, 0), (-1, -1), (0, -1), (1, -1))


def Q(x):
    return 0.5 * math.erfc(x / math.sqrt(2.0))


def _directed(x, sig, mode):
    d = Decimal(repr(float(x)))
    if d == 0:
        return "0"
    q = d.quantize(Decimal(1).scaleb(d.adjusted() - sig + 1), rounding=mode)
    return f"{q:.{sig - 1}e}" if (q.adjusted() < -3 or q.adjusted() >= 4) else f"{q:f}"


def up(x, sig=3):        # upper bounds are printed rounded up ...
    return _directed(x, sig, ROUND_CEILING)


def dn(x, sig=3):        # ... and lower bounds rounded down
    return _directed(x, sig, ROUND_FLOOR)


def hex_cells(R):
    return [(m, n) for m in range(-R, R + 1) for n in range(-R, R + 1) if max(abs(m), abs(n), abs(m + n)) <= R]


# ============================== Section 3 ==============================
def kagome_H(cells, t1, t2):
    c2i = {c: i for i, c in enumerate(cells)}
    H = np.zeros((3 * len(cells), 3 * len(cells)))
    def hop(u, v, t):
        H[u, v] -= t; H[v, u] -= t
    for i, (m, n) in enumerate(cells):
        hop(3*i, 3*i+1, t1); hop(3*i+1, 3*i+2, t1); hop(3*i+2, 3*i, t1)
        for (dm, dn), ks, kd in INTERCELL:
            if (m + dm, n + dn) in c2i:
                hop(3*i + ks, 3*c2i[(m + dm, n + dn)] + kd, t2)
    return H, c2i


def verify_sec3(s):
    t1 = 0.5*s.t_sum*(1 + s.delta_b); t2 = 0.5*s.t_sum*(1 - s.delta_b)
    rng = np.random.default_rng(7)
    err = 0.0
    for _ in range(400):                       # Thm 3.1 at random k (k.a1, k.a2 uniform)
        k1, k2 = rng.uniform(-math.pi, math.pi, 2); k3 = k2 - k1
        Hk = -np.array([[0, t1 + t2*np.exp(-1j*k1), t1 + t2*np.exp(-1j*k2)],
                        [t1 + t2*np.exp(1j*k1), 0, t1 + t2*np.exp(-1j*k3)],
                        [t1 + t2*np.exp(1j*k2), t1 + t2*np.exp(1j*k3), 0]])
        Lam = 3 + 2*(math.cos(k1) + math.cos(k2) + math.cos(k3))
        r = math.sqrt(2.25*(t1 - t2)**2 + t1*t2*Lam)
        ana = np.sort([t1 + t2, -(t1 + t2)/2 + r, -(t1 + t2)/2 - r])
        err = max(err, float(np.abs(np.linalg.eigvalsh(Hk) - ana).max()))
    assert err < 1e-12
    gap_K = 3*abs(t1 - t2)
    cells = hex_cells(3)                         # Thm 3.2(ii): CLS on an internal hexagon
    H, c2i = kagome_H(cells, t1, t2)
    ring = [(3*c2i[(0, 0)]+1, 1), (3*c2i[(0, 0)]+2, -1), (3*c2i[(0, 1)]+0, 1),
            (3*c2i[(0, 1)]+1, -1), (3*c2i[(1, 0)]+2, 1), (3*c2i[(1, 0)]+0, -1)]
    psi = np.zeros(H.shape[0])
    for v, sg in ring:
        psi[v] = sg/math.sqrt(6)
    res = float(np.linalg.norm(H @ psi - (t1 + t2)*psi))
    assert res < 1e-13
    # phase of the lowest band on the R_hex flake cut along complete up-triangle cells
    flake = hex_cells(s.R_hex)
    lo = -(t1 + t2)/2 - 1.5*abs(t1 - t2); hi = -(t1 + t2)/2 + 1.5*abs(t1 - t2)
    phase = {}
    for db in (s.delta_b, -s.delta_b):
        a, b = 0.5*s.t_sum*(1 + db), 0.5*s.t_sum*(1 - db)
        ev = np.linalg.eigvalsh(kagome_H(flake, a, b)[0])
        phase[db] = (int(np.sum(ev <= lo + 1e-9)), int(np.sum((ev > lo + 1e-9) & (ev < hi - 1e-9))))
    below, ingap = phase[s.delta_b]
    assert below == len(flake) and ingap == 0         # trivial (atomic) phase for t1 > t2
    print(f"[Sec.3] Thm3.1 max|E_num-E_formula| over 400 k <= {up(err, 2)}; gap at K = {gap_K:.4f}; "
          f"Thm3.2 CLS residual <= {up(res, 2)}")
    print(f"[Sec.3] {len(flake)}-cell flake: delta_b={s.delta_b:+.2f}: {below} states below gap, {ingap} in gap "
          f"(trivial); delta_b={-s.delta_b:+.2f}: {phase[-s.delta_b][0]} below, {phase[-s.delta_b][1]} in gap (obstructed)")


# ============================== Section 4 ==============================
def verify_sec4(s):
    a2, rx, ry, rc, eta = s.a2_mm, s.rx_mm, s.ry_mm, s.rho_c, s.eta_r
    ab = 0.5*(1 - s.delta_b); d_rel = (s.a1_mm - a2)/a2
    phi_inv = (math.sqrt(5) - 1)/2; phi_inv2 = phi_inv**2; s32 = math.sqrt(3)/2
    U = 2.0**-53
    # (1) intracell phase offsets are cell-independent constants
    D21 = (2*math.pi*(ab*phi_inv + 1/3)) % (2*math.pi)
    D31 = (2*math.pi*(ab*(phi_inv/2 + s32*phi_inv2) + 2/3)) % (2*math.pi)
    C_up = np.array([[math.cos(D21) - 1, math.cos(D31) - 1], [math.sin(D21), math.sin(D31)]])
    qnorm = 2*math.pi*math.hypot(phi_inv, phi_inv2)

    def ref(m, n, k):
        x = a2*(m + 0.5*n + np.where(k == 2, ab, np.where(k == 3, 0.5*ab, 0.0)))
        y = a2*s32*(n + np.where(k == 3, ab, 0.0))
        return x, y

    def vmap(m, n, k):
        xr, yr = ref(m, n, k)
        rho = np.hypot(xr/rx, yr/ry); th = np.arctan2(yr/ry, xr/rx)
        u = np.clip((rho - rc)/(1 - rc), 0.0, 1.0)
        R = np.where(rho <= rc, rho, np.where(rho <= 1.0, rho + d_rel*(1 - rc)*(u**6 - 3*u**5 + 2.5*u**4),
                                              1 + 0.5*d_rel*(1 - rc) + (1 + d_rel)*(rho - 1)))
        inside = (u > 0) & (u < 1)
        w = np.where(inside, 16*u**2*(1 - u)**2, 0.0)
        psi = 2*math.pi*phi_inv*(m*m + m*n + n*n) + (2*math.pi/a2)*(phi_inv*xr + phi_inv2*yr) + 2*math.pi*(k - 1)/3
        f = np.where(rho > 0, R/np.where(rho > 0, rho, 1.0), 1.0)
        sc = np.where(inside, eta*a2*f*w, 0.0)
        X0, Y0 = rx*R*np.cos(th), ry*R*np.sin(th)
        return dict(X0=X0, Y0=Y0, dx=sc*np.cos(psi), dy=sc*np.sin(psi), inside=inside, f=f,
                    big=np.maximum(np.abs(X0), np.abs(Y0)))

    mm = int(math.ceil(rx/a2 + ry/(a2*math.sqrt(3)))) + 3
    nn = int(math.ceil(ry/(a2*s32))) + 3
    M, N = np.meshgrid(np.arange(-mm, mm + 1), np.arange(-nn, nn + 1), indexing='ij')
    M, N = M.ravel(), N.ravel()
    one = np.ones_like(M)
    xa, ya = ref(M, N, one); rho_a = np.hypot(xa/rx, ya/ry)
    v38_sel = (rho_a >= rc) & (rho_a <= 1.0)       # anchor-vertex rule used in v3.8
    ell = {'up': ab*a2, 'dn': (1 - ab)*a2}
    res = {}
    for fam in ('up', 'dn'):
        V = ([vmap(M, N, one*k) for k in (1, 2, 3)] if fam == 'up' else
             [vmap(M, N, 3*one), vmap(M, N + 1, one), vmap(M - 1, N + 1, 2*one)])
        E0 = np.array([[V[1]['X0'] - V[0]['X0'], V[2]['X0'] - V[0]['X0']],
                       [V[1]['Y0'] - V[0]['Y0'], V[2]['Y0'] - V[0]['Y0']]])
        dE = np.array([[V[1]['dx'] - V[0]['dx'], V[2]['dx'] - V[0]['dx']],
                       [V[1]['dy'] - V[0]['dy'], V[2]['dy'] - V[0]['dy']]])
        sel = V[0]['inside'] | V[1]['inside'] | V[2]['inside']       # every triangle with a perturbed vertex
        det = lambda A: A[0, 0]*A[1, 1] - A[0, 1]*A[1, 0]
        d0 = det(E0)
        inv = np.array([[E0[1, 1], -E0[0, 1]], [-E0[1, 0], E0[0, 0]]])/d0
        Mx = np.einsum('ijt,jkt->ikt', inv, dE)
        fr = (Mx**2).sum((0, 1)); dm = det(Mx)
        nrm = np.sqrt(0.5*(fr + np.sqrt(np.maximum(0, fr**2 - 4*dm**2))))
        dIM = (1 + Mx[0, 0])*(1 + Mx[1, 1]) - Mx[0, 1]*Mx[1, 0]
        E = E0 + dE; dE_ = det(E)
        fr0 = (E0**2).sum((0, 1)); smin0 = np.abs(d0)/np.sqrt(0.5*(fr0 + np.sqrt(np.maximum(0, fr0**2 - 4*d0**2))))
        fbar = (V[0]['f'] + V[1]['f'] + V[2]['f'])/3
        ratio = (1/smin0)/(math.sqrt(2)/(ell[fam]*fbar))
        big = np.maximum.reduce([V[i]['big'] for i in range(3)])
        ce = 2*64*U*big + 64*U*ell[fam]                       # generous per-entry error of E
        Emax = np.abs(E).max((0, 1))
        derr = 2*(2*Emax*ce + ce**2) + 4*U*Emax**2
        dref = s32*ell[fam]**2
        res[fam] = dict(n=int(sel.sum()), n_v38=int(v38_sel.sum()), missed=int((sel & ~v38_sel).sum()),
                        maxM=float(nrm[sel].max()), min_dIM=float(dIM[sel].min()),
                        min_ratio=float((dE_/dref)[sel].min()), min_d0=float((d0/dref)[sel].min()),
                        max_inv=float(ratio[sel].max()), margin=float((dE_/derr)[sel].min()))
        assert res[fam]['maxM'] < 1 and res[fam]['min_dIM'] > 0 and res[fam]['min_d0'] > 0 and res[fam]['margin'] > 1e6

    # first-order (a2 -> 0) supremum: M ~ eta*W*(a2/ell)*T*(I - beta*D e e^T D^-1)*[chords]
    Dm, Di = np.array([rx, ry]), np.array([1/rx, 1/ry])
    T = {'up': np.linalg.inv(np.array([[1, 0.5], [0, s32]])), 'dn': np.linalg.inv(np.array([[0.5, -0.5], [s32, s32]]))}
    def n2(A):  # spectral norm of stacked 2x2 matrices A[...,2,2]
        F = (A**2).sum((-1, -2)); d = A[..., 0, 0]*A[..., 1, 1] - A[..., 0, 1]*A[..., 1, 0]
        return np.sqrt(0.5*(F + np.sqrt(np.maximum(0, F**2 - 4*d**2))))
    def P_of(u, th):
        rho = rc + u*(1 - rc)
        R = rho + d_rel*(1 - rc)*(u**6 - 3*u**5 + 2.5*u**4); Rp = 1 + d_rel*(6*u**5 - 15*u**4 + 10*u**3)
        beta = 1 - (R/rho)/Rp
        e = np.stack([np.cos(th), np.sin(th)], -1)
        Pm = np.eye(2) - beta[..., None, None]*np.einsum('...i,...j->...ij', Dm*e, Di*e)
        return 16*u**2*(1 - u)**2, Pm
    def sup_dn(us, ths):
        uu, tt = np.meshgrid(us, ths, indexing='ij'); W, Pm = P_of(uu, tt)
        val = W*n2(T['dn'] @ Pm)*2*math.sqrt(2)*eta*a2/ell['dn']
        i = np.unravel_index(np.argmax(val), val.shape); return float(val[i]), us[i[0]], ths[i[1]]
    def sup_up(us, ths, ps):
        uu, tt, pp = np.meshgrid(us, ths, ps, indexing='ij'); W, Pm = P_of(uu, tt)
        Rot = np.stack([np.stack([np.cos(pp), -np.sin(pp)], -1), np.stack([np.sin(pp), np.cos(pp)], -1)], -2)
        val = W*n2(T['up'] @ Pm @ Rot @ C_up)*eta*a2/ell['up']
        i = np.unravel_index(np.argmax(val), val.shape); return float(val[i]), us[i[0]], ths[i[1]], ps[i[2]]
    v, u0, t0 = sup_dn(np.linspace(0.01, 0.99, 491), np.linspace(-math.pi, math.pi, 1441))
    for h in (0.01, 0.001):
        v, u0, t0 = sup_dn(np.linspace(u0 - h, u0 + h, 201), np.linspace(t0 - 5*h, t0 + 5*h, 201))
    fo_dn = v
    v, u0, t0, p0 = sup_up(np.linspace(0.05, 0.95, 61), np.linspace(-math.pi, math.pi, 145), np.linspace(-math.pi, math.pi, 145))
    for h in (0.02, 0.002, 0.0002):
        v, u0, t0, p0 = sup_up(np.linspace(u0 - h, u0 + h, 41), np.linspace(t0 - 2*h, t0 + 2*h, 41), np.linspace(p0 - 2*h, p0 + 2*h, 41))
    fo_up = v
    mu = 0.5*(ry/rx + rx/ry)
    print(f"[Sec.4] |q_phi| a2 = {qnorm:.4f}; intracell offsets psi2-psi1 = {math.degrees(D21):.2f} deg, "
          f"psi3-psi1 = {math.degrees(D31):.2f} deg; ||C_up||_2 = {np.linalg.norm(C_up, 2):.4f}")
    for fam, name in (('up', 'up  '), ('dn', 'down')):
        r = res[fam]
        print(f"[Sec.4] {name}: {r['n']} perturbed triangles ({r['missed']} not in v3.8 anchor rule of {r['n_v38']}); "
              f"max||E0^-1 dE|| = {r['maxM']:.5f}; min det(I+M) = {r['min_dIM']:.5f}; "
              f"min det(E)/det_ref = {r['min_ratio']:.4f}; min det(E0)/det_ref = {r['min_d0']:.4f}")
        print(f"[Sec.4] {name}: max ||E0^-1|| / (sqrt2/(ell R/rho)) = {r['max_inv']:.4f}; "
              f"min det(E)/float-error bound >= {dn(r['margin'], 2)}")
    print(f"[Sec.4] first-order sup: up {fo_up:.4f}, down {fo_dn:.4f}; rigorous first-order bound 4*mu*eta*a2/ell "
          f"(mu={mu:.3f}): up {4*mu*eta/ab:.4f}, down {4*mu*eta/(1 - ab):.4f}")
    return res


# ============================== Section 5 ==============================
class Network:
    def __init__(self, s):
        self.s = s
        self.cells = hex_cells(s.R_hex)
        self.c2i = {c: i for i, c in enumerate(self.cells)}
        self.N = len(self.cells)
        depth = lambda c: max(abs(c[0]), abs(c[1]), abs(c[0] + c[1]))
        self.int_cells = [i for i, c in enumerate(self.cells) if depth(c) < s.R_hex]
        self.bnd_cells = [i for i, c in enumerate(self.cells) if depth(c) == s.R_hex]
        self.I = np.array([3*i + k for i in self.int_cells for k in range(3)])
        self.B = np.array([3*i + k for i in self.bnd_cells for k in range(3)])
        self.N_I = len(self.int_cells)
        self.edges = []
        for i, (m, n) in enumerate(self.cells):
            self.edges += [(3*i, 3*i + 1, i, i, 1), (3*i + 1, 3*i + 2, i, i, 1), (3*i + 2, 3*i, i, i, 1)]
            for (dm, dn), ks, kd in INTERCELL:
                if (m + dm, n + dn) in self.c2i:
                    j = self.c2i[(m + dm, n + dn)]
                    self.edges.append((3*i + ks, 3*j + kd, i, j, 2))
        self.m = np.array([self.cells[i][0] for i in self.int_cells])
        self.n = np.array([self.cells[i][1] for i in self.int_cells])
        self.p, self.q = self.m + self.n, self.m - self.n
        self.pos = {(int(a), int(b)): c for c, (a, b) in enumerate(zip(self.m, self.n))}
        # chord families: lines of constant m, n, m+n and (Revision 4.0) m-n
        self.fam = [self.m, self.n, self.p, self.q][:s.chord_families]
        self.off, rows = [], 0
        for f in self.fam:
            self.off.append(rows - int(f.min())); rows += int(f.max() - f.min()) + 1
        self.C = np.zeros((rows, self.N_I))
        for f, o in zip(self.fam, self.off):
            self.C[f + o, np.arange(self.N_I)] = 1
        self.line_len = self.C.sum(1)
        self.cell_bonds = {i: [] for i in range(self.N)}            # intercell bonds touching each cell
        for (u, v, i, j, kind) in self.edges:
            if kind == 2:
                self.cell_bonds[i].append((u, v)); self.cell_bonds[j].append((u, v))

    def laplacian(self, state, g2_switch=True):
        """state: per-cell values in [0, 1] (length N); complex values allowed (complex-step)"""
        s = self.s
        state = np.asarray(state)
        L = np.zeros((3*self.N, 3*self.N), dtype=np.result_type(state.dtype, float))
        for (u, v, i, j, kind) in self.edges:
            g = s.g1 if kind == 1 else s.g2_base + (0.5*(state[i] + state[j])*(s.g2_on - s.g2_base) if g2_switch else 0.0)
            L[u, u] += g; L[v, v] += g; L[u, v] -= g; L[v, u] -= g
        for i in range(self.N):
            for k in range(3):
                L[3*i + k, 3*i + k] += (s.y_base + state[i]*s.y_ridge)/3.0
        return L

    def schur(self, L):
        I, B = self.I, self.B
        return L[np.ix_(B, B)] - L[np.ix_(B, I)] @ np.linalg.solve(L[np.ix_(I, I)], L[np.ix_(I, B)])

    def A_of(self, G):
        Gr = G.reshape(self.N_I, 3, -1)
        return (Gr**2).sum(1).T/3.0                     # A[b, c] = (1/3) sum_k G[3c+k, b]^2

    def gate(self, P, thr):
        on = P > thr
        ok = np.ones(self.N_I, dtype=bool)
        for f, o in zip(self.fam, self.off):
            ok &= on[f + o]
        return [int(c) for c in np.flatnonzero(ok)]

    def interior_vector(self, cells):
        v = np.zeros(self.N_I)
        for c in cells:
            v[self.pos[c]] = 1.0
        return v

    def cells_of(self, positions):
        return [(int(self.m[c]), int(self.n[c])) for c in positions]

    def state(self, cells):
        x = np.zeros(self.N)
        for c in cells:
            x[self.c2i[c]] = 1.0
        return x


def linear_model(net):
    L0 = net.laplacian(np.zeros(net.N), g2_switch=False)
    G0 = -np.linalg.solve(L0[np.ix_(net.I, net.I)], L0[np.ix_(net.I, net.B)])
    return L0, G0, net.A_of(G0)


def forward(net, L0, G0, cells, full=False):
    """Lambda(s) - Lambda(0) for ridge cells (shunt + g2 switching), without cancellation:
    G(s)^T dL_II G0, valid because interior ridges leave L_IB and L_BB unchanged."""
    I, B = net.I, net.B
    Ls = net.laplacian(net.state(cells), g2_switch=True)
    dL = Ls - L0
    assert not dL[np.ix_(I, B)].any() and not dL[np.ix_(B, B)].any()
    Gs = -np.linalg.solve(Ls[np.ix_(I, I)], Ls[np.ix_(I, B)])
    if full:
        return Gs.T @ dL[np.ix_(I, I)] @ G0
    return np.einsum('ib,ij,jb->b', Gs, dL[np.ix_(I, I)], G0)


def verify_frechet(net, L0, G0, A):
    s = net.s
    I, B = net.I, net.B
    worst = 0.0                                    # complex-step derivative of diag(Lambda), every interior cell
    h = 1e-30
    for pos, ci in enumerate(net.int_cells):
        z = np.zeros(3*net.N, dtype=complex); z[3*ci:3*ci + 3] = 1j*h/3
        d = np.diag(net.schur(L0.astype(complex) + np.diag(z))).imag/h
        nz = A[:, pos] != 0
        assert np.all(d[~nz] == 0)
        worst = max(worst, float(np.max(np.abs(d[nz] - A[nz, pos])/A[nz, pos])))
    assert worst < 1e-9                             # bound stated in Sec. 5.2 (holds on both reference machines)
    rng = np.random.default_rng(3)                  # exact second-order remainder identity (Thm 5.1(3))
    dy = rng.uniform(0, 1e-7, net.N_I)
    dY = np.diag(np.repeat(dy/3, 3))
    st = np.zeros(net.N); st[net.int_cells] = dy/s.y_ridge
    dLam = net.schur(net.laplacian(st, g2_switch=False)) - net.schur(L0)
    LII = L0[np.ix_(I, I)]
    R2 = -G0.T @ dY @ np.linalg.solve(LII + dY, dY @ G0)
    ident = float(np.linalg.norm(dLam - (G0.T @ dY @ G0 + R2))/np.linalg.norm(dLam))
    ev = np.linalg.eigvalsh(LII)
    bnd = np.linalg.norm(G0, 2)**2*(dy.max()/3)**2/ev[0]
    assert ident < 2e-8 and np.linalg.norm(R2, 2) <= bnd   # 2e-8: bound stated in Sec. 5.2
    print(f"[Sec.5] y_base={s.y_base:.0e}: Thm5.1 complex-step max rel. err over {net.N_I} columns <= {up(worst, 2)}; "
          f"remainder identity residual <= {up(ident, 2)}; ||R2|| = {np.linalg.norm(R2, 2):.2e} <= {bnd:.2e}")
    print(f"[Sec.5] y_base={s.y_base:.0e}: lambda_min(L_II) = {ev[0]:.3e} S, cond(L_II) = {ev[-1]/ev[0]:.2e}; "
          f"remainder bound at dy = y_ridge: {np.linalg.norm(G0, 2)**2*(s.y_ridge/3)**2/ev[0]:.2e} S")


def modp_eliminate(M, p, ncols):
    """Gauss-Jordan over F_p on the first ncols columns; returns (rank, reduced matrix)."""
    M = M % p
    r = 0
    for c in range(ncols):
        nz = np.flatnonzero(M[r:, c])
        if nz.size == 0:
            continue
        piv = r + nz[0]
        if piv != r:
            M[[r, piv]] = M[[piv, r]]
        M[r] = (M[r]*pow(int(M[r, c]), p - 2, p)) % p
        f = M[:, c].copy(); f[r] = 0
        M = (M - np.outer(f, M[r])) % p
        r += 1
        if r == M.shape[0]:
            break
    return r, M


def rank_exact(M):
    """exact rank of an integer matrix: fraction-free (Bareiss) elimination, exact divisions asserted"""
    A = [[int(x) for x in row] for row in M]
    rows, cols = len(A), len(A[0])
    r, prev = 0, 1
    for c in range(cols):
        piv = next((i for i in range(r, rows) if A[i][c] != 0), None)
        if piv is None:
            continue
        A[r], A[piv] = A[piv], A[r]
        for i in range(r + 1, rows):
            if A[i][c] == 0:
                for j in range(c + 1, cols):
                    num = A[i][j]*A[r][c]
                    assert num % prev == 0
                    A[i][j] = num//prev
                continue
            for j in range(c + 1, cols):
                num = A[i][j]*A[r][c] - A[i][c]*A[r][j]
                assert num % prev == 0
                A[i][j] = num//prev
            A[i][c] = 0
        prev = A[r][c]; r += 1
        if r == rows:
            break
    return r


def modular_model(net):
    """exact integer Laplacian (all parameters are exact decimals) and A_I mod p; structural row bound"""
    s = net.s
    gq, g2q, shq = Fraction(str(s.g1)), Fraction(str(s.g2_base)), Fraction(str(s.y_base))/3
    scale = math.lcm(gq.denominator, g2q.denominator, shq.denominator)
    gi, g2i, shi = int(gq*scale), int(g2q*scale), int(shq*scale)
    Lz = np.zeros((3*net.N, 3*net.N), dtype=np.int64)
    for (u, v, i, j, kind) in net.edges:
        g = gi if kind == 1 else g2i
        Lz[u, u] += g; Lz[v, v] += g; Lz[u, v] -= g; Lz[v, u] -= g
    Lz[np.diag_indices_from(Lz)] += shi
    assert np.allclose(Lz/scale, net.laplacian(np.zeros(net.N), g2_switch=False), rtol=1e-15, atol=0)
    LII, LIB = Lz[np.ix_(net.I, net.I)], Lz[np.ix_(net.I, net.B)]
    nzcols = [c for c in range(LIB.shape[1]) if LIB[:, c].any()]
    n_distinct = len({tuple(LIB[:, c]) for c in nzcols})            # rank A_I <= n_distinct
    p = 2147483647
    r, X = modp_eliminate(np.concatenate([LII, LIB], 1).astype(np.int64), p, LII.shape[0])
    assert r == LII.shape[0]
    G = X[:, LII.shape[0]:]                                           # = LII^{-1} LIB mod p
    Gr = G.reshape(net.N_I, 3, -1)
    Ap = ((Gr*Gr % p).sum(1) % p).T*pow(3, p - 2, p) % p
    rA = modp_eliminate(Ap.copy(), p, Ap.shape[1])[0]
    assert rA == n_distinct                                           # bounds meet: rank A_I exact over Q
    return dict(p=p, Ap=Ap, G=G, rA=rA, n_zero=len(net.B) - len(nzcols), n_nonzero=len(nzcols), n_distinct=n_distinct)


def verify_rank(net, mod, label):
    p, Ap, rA = mod['p'], mod['Ap'], mod['rA']
    Ci = net.C.astype(np.int64)
    rC = rank_exact(Ci)
    assert modp_eliminate(Ci.copy(), p, Ci.shape[1])[0] == rC
    rAug = modp_eliminate(np.concatenate([Ap, Ci]), p, Ap.shape[1])[0]
    assert rAug == rA + rC                                            # = upper bound rank A_I + rank C
    print(f"[Sec.5] {label}: {len(net.B)} boundary nodes, {mod['n_zero']} without interior neighbour, "
          f"{mod['n_distinct']} distinct rows -> rank A_I = {rA}; {len(net.fam)} chord families, {Ci.shape[0]} lines, "
          f"rank C = {rC}; rank [A_I; w_c C] = {rAug} of {net.N_I}; dim ker = {net.N_I - rAug} (exact)")
    return Ci, rC, rAug


def stage_b(net, L0, G0, A, mod, Ci, label, n_trials=2000, full_dtn=False):
    s = net.s
    p, Ap, Gp = mod['p'], mod['Ap'], mod['G']
    y_r, sig = s.y_ridge, s.sigma_rel
    st = np.zeros(net.N_I)
    for rc in RIDGES:
        st[net.pos[rc]] = 1.0
    cache = {}
    def fwd(on):
        key = tuple(sorted(int(c) for c in on))
        if key not in cache:
            cache[key] = forward(net, L0, G0, [net.cells[net.int_cells[c]] for c in key])
        return cache[key]
    dL = fwd(np.flatnonzero(st))
    nscale = float(np.abs(dL).max())
    w_c = nscale/y_r
    nlines = len(net.line_len)
    zmax = 1/(2*sig*math.sqrt(net.line_len.max()))
    tight = sum(Q(1/(2*sig*math.sqrt(L))) for L in net.line_len)
    cand0 = net.gate(net.C @ (st*y_r), 0.5*y_r)
    K0 = len(cand0)
    CK = net.C[:, cand0]
    cnt = CK @ st[cand0]
    pats = [np.array(b) for b in itertools.product((0, 1), repeat=K0) if np.array_equal(CK @ np.array(b), cnt)]
    i_true = [i for i, b in enumerate(pats) if np.array_equal(b, st[cand0])][0]
    rCK = modp_eliminate(Ci[:, cand0].copy(), p, K0)[0]
    ring = ((1, 0), (0, 1), (-1, 1), (-1, 0), (0, -1), (1, -1))
    loc = {(int(net.m[c]), int(net.n[c])): j for j, c in enumerate(cand0)}
    rings = []
    for (a, b) in loc:                                              # hexagonal-ring kernel vectors of C_K
        cells = [(a + dm, b + dn) for dm, dn in ring]
        if all(c in loc for c in cells):
            v = np.zeros(K0)
            for t, c in enumerate(cells):
                v[loc[c]] = (-1)**t
            if not (CK @ v).any():
                rings.append(v)
    nring = int(np.linalg.matrix_rank(np.array(rings))) if rings else 0
    assert rCK == K0 - nring                                        # rank C_K exact (mod-p lower = ring upper bound)
    rK = modp_eliminate(np.concatenate([Ap[:, cand0], Ci[:, cand0]]), p, K0)[0]
    Aaug = np.vstack([A, w_c*net.C])
    sv = np.linalg.svd(Aaug[:, cand0], compute_uv=False)
    sigs = [fwd([cand0[j] for j in np.flatnonzero(b)]) for b in pats]
    seps = [float(np.linalg.norm(sigs[i] - dL)/(sig*nscale)) for i in range(len(pats)) if i != i_true]
    lin = A @ st*y_r
    print(f"[Sec.5] {label}: chord gate {nlines} lines, p_chord <= Q({zmax:.3f}) <= {up(Q(zmax))}, FWER <= {up(nlines*Q(zmax))} "
          f"(line-length-aware <= {up(tight)}); max|dLambda| = {nscale:.4e} S, w_c = {w_c:.2e}; "
          f"||dL - A s y_r||/||dL|| = {np.linalg.norm(dL - lin)/np.linalg.norm(dL):.1f}")
    print(f"[Sec.5] {label}: K = {K0}, rank C_K = {rCK}, rank [A_K; w_c C_K] = {rK} (exact), "
          f"sigma(A_K) ratio {sv[-1]/sv[0]:.1e}; {len(pats)} binary pattern(s) share the chord data" +
          (f"; ghost separation {', '.join(f'{x:.4f}' for x in seps)} std -> any decoder errs w.p. >= "
           f"{dn(Q(min(seps)/2))}" if seps else "; the chord data determine the ridge"))
    if full_dtn:
        iu = np.triu_indices(len(net.B))
        Ff = [forward(net, L0, G0, [net.cells[net.int_cells[cand0[j]]] for j in np.flatnonzero(b)], full=True)[iu] for b in pats]
        sdf = sig*np.abs(Ff[i_true]).max()
        sepf = [np.linalg.norm(Ff[i] - Ff[i_true])/sdf for i in range(len(pats)) if i != i_true]
        Gr = Gp.reshape(net.N_I, 3, -1)
        J = np.zeros((len(iu[0]), net.N_I), dtype=np.int64)
        for c in range(net.N_I):
            acc = np.zeros(len(iu[0]), dtype=np.int64)
            for k in range(3):
                acc = (acc + Gr[c, k, iu[0]]*Gr[c, k, iu[1]] % p) % p
            J[:, c] = acc*pow(3, p - 2, p) % p
        rJ = modp_eliminate(J, p, net.N_I)[0]
        print(f"[Sec.5] {label}: full Schur matrix: linearised Jacobian rank {rJ} of {net.N_I} (exact); ghost "
              f"separation {', '.join(f'{x:.3f}' for x in sepf)} std -> decoder error >= {dn(Q(min(sepf)/2))}; "
              f"sigma(A_K) = [{sv[0]:.2e}, {sv[-1]:.2e}]")
    # Monte Carlo: same draws as Revision 3.8/3.9 (np.random.seed(0); dLambda noise then cell noise)
    def exact(cand, on_mask):
        rec = np.zeros(net.N_I); rec[[c for c, o in zip(cand, on_mask) if o]] = 1
        return bool(np.array_equal(rec, st))
    np.random.seed(0)
    ok_ls = ok_ml = 0
    Ks = []; t_ml = 0.0; max_est = 0.0
    pat_cache = {}
    for _ in range(n_trials):
        dLn = dL + np.random.normal(0, sig*nscale, size=dL.shape)
        vals = st*y_r + np.random.normal(0, sig*y_r, net.N_I)
        P = net.C @ vals
        cand = net.gate(P, 0.5*y_r); K = len(cand); Ks.append(K)
        if K == 0:
            continue
        AK = Aaug[:, cand]
        b = np.concatenate([dLn/y_r, w_c*P/y_r])
        est = np.linalg.solve(AK.T @ AK + 1e-10*np.eye(K), AK.T @ b)      # Revision-3.8 ridge least squares
        ok_ls += exact(cand, est >= 0.45)
        max_est = max(max_est, float(np.abs(est).max()))
        key = (tuple(cand), tuple(np.rint(P/y_r).astype(int)))             # maximum likelihood over
        if key not in pat_cache:                                           # chord-consistent binary patterns
            CKk = net.C[:, cand]; target = np.array(key[1])
            on = [[cand[j] for j in np.flatnonzero(bb)] for bb in itertools.product((0, 1), repeat=K)
                  if np.array_equal(CKk @ np.array(bb), target)]
            pat_cache[key] = (on, np.array([fwd(o) for o in on]))
        on, S = pat_cache[key]
        t0 = time.perf_counter()
        best = on[int(np.argmin(((S - dLn)**2).sum(1)))]
        t_ml += time.perf_counter() - t0
        rec = np.zeros(net.N_I); rec[best] = 1
        ok_ml += bool(np.array_equal(rec, st))
    Ks = np.array(Ks)
    print(f"[Sec.5] {label}: MC {n_trials} @ {20*math.log10(1/sig):.1f} dB: K in [{Ks.min()}, {Ks.max()}]; "
          f"ridge-LS (lambda=1e-10) {ok_ls}/{n_trials} (max |s_hat| <= {up(max_est, 2)}); maximum likelihood {ok_ml}/{n_trials} "
          f"({1e3*t_ml/n_trials:.3f} ms/trial select)")
    return dict(K=K0, patterns=len(pats), seps=seps, ok_ml=ok_ml, ok_ls=ok_ls)


def switching_component(net):
    """prod_d (1 - T^d) over the line directions of the active chord families: +-1 cells with zero line sums"""
    dirs = [(0, 1), (1, 0), (1, -1), (1, 1)][:net.s.chord_families]   # keep m, n, m+n, m-n constant
    comp = {(0, 0): 1}
    for d in dirs:
        new = dict(comp)
        for (a, b), v in comp.items():
            new[(a + d[0], b + d[1])] = new.get((a + d[0], b + d[1]), 0) - v
        comp = {k: v for k, v in new.items() if v != 0}
    return sorted(k for k, v in comp.items() if v > 0), sorted(k for k, v in comp.items() if v < 0)


def ghost_pairs(net, L0, G0, label):
    """Theorem 5.2(5): binary ridge pairs with identical chord data in every active family."""
    s = net.s
    plus, minus = switching_component(net)
    depth = lambda c: max(abs(c[0]), abs(c[1]), abs(c[0] + c[1]))
    def sep(RA, RB):
        va, vb = np.zeros(net.N_I), np.zeros(net.N_I)
        for c in RA: va[net.pos[c]] = 1
        for c in RB: vb[net.pos[c]] = 1
        assert np.array_equal(net.C @ va, net.C @ vb)                  # identical chord sums, all families
        sa, sb = forward(net, L0, G0, RA), forward(net, L0, G0, RB)
        return float(np.linalg.norm(sa - sb)/(s.sigma_rel*np.abs(sa).max()))
    shifts = [(a, b) for a in range(-8, 9) for b in range(-8, 9)]
    deep = lambda c, r: c in net.pos and depth(c) <= r            # r <= R_hex-2: no ridge bond touches B
    iso = [(sep([(x + a, y + b) for x, y in plus], [(x + a, y + b) for x, y in minus]), (a, b)) for a, b in shifts
           if all(deep((x + a, y + b), s.R_hex - 2) for x, y in plus + minus)]
    ring = [c for c in net.pos if depth(c) == s.R_hex - 2]           # closed loop of ridge cells at depth 2
    framed = [(sep(ring + [(x + a, y + b) for x, y in plus], ring + [(x + a, y + b) for x, y in minus]), (a, b))
              for a, b in shifts
              if all((x + a, y + b) in net.pos and depth((x + a, y + b)) < s.R_hex - 2 for x, y in plus + minus)]
    iso.sort(); framed.sort()
    print(f"[Sec.5] Thm5.2(5) {label}: {len(plus)}+{len(minus)}-cell switching component {plus} vs {minus}; "
          f"isolated (all cells >= 2 rings deep): {len(iso)} placements, min separation >= {dn(iso[0][0], 4)} std at shift {iso[0][1]}")
    print(f"[Sec.5] Thm5.2(5) {label}: inside a closed {len(ring)}-cell ridge loop: {len(framed)} placements, separation "
          f"{framed[0][0]:.1e} to <= {up(framed[-1][0], 2)} std -> any decoder errs w.p. >= {dn(Q(framed[-1][0]/2), 4)} on one of each pair")
    assert Q(framed[-1][0]/2) > 0.49 and iso[0][0] > 10
    return iso, framed




# ======================= Section 5.4 (Revision 5.1) =======================
forward_diag = forward


def verify_sec2():
    """Sections 2 and 3.1: five-phase closure, Maxwell-Garnett permittivity, g33 and wake voltage."""
    V = (0.220, 0.450, 0.050, 0.044, 0.236); rho = (1480, 180, 4200, 1450, 1.2)
    assert abs(sum(V) - 1) < 1e-12
    rho_eff = sum(v*r for v, r in zip(V, rho))
    W = [100*v*r/rho_eff for v, r in zip(V, rho)]
    eps_bt, eps_m, phi, d33, eps0 = 1700.0, 3.2, 0.20, 2.0e-12, 8.8541878128e-12
    beta = (eps_bt - eps_m)/(eps_bt + 2*eps_m)
    eps_eff = eps_m*(1 + 2*phi*beta)/(1 - phi*beta)
    g33 = d33/(eps0*eps_eff)
    V_wake = g33*1.5/3.46e-6*4e-6
    assert round(rho_eff, 1) == 680.7 and round(eps_eff, 2) == 5.58 and round(1e3*V_wake, 1) == 70.2
    print(f"[Sec.2] rho_eff = {rho_eff:.4f} kg/m^3; W_i = {', '.join(f'{w:.2f}' for w in W)} %; "
          f"[Sec.3.1] beta = {beta:.4f}, eps_eff = {eps_eff:.4f}, g33 = {1e3*g33:.2f} mV m/N, V_wake = {1e3*V_wake:.2f} mV")


def Q_inv(p):
    """inverse of the Gaussian tail Q, by bisection"""
    lo, hi = 0.0, 40.0
    for _ in range(200):
        mid = 0.5*(lo + hi)
        lo, hi = (mid, hi) if Q(mid) > p else (lo, mid)
    return 0.5*(lo + hi)


def radius(c):
    return max(abs(c[0]), abs(c[1]), abs(c[0] + c[1]))


def rank_modp(M, p=2147483647):
    M = np.asarray(M, dtype=np.int64)
    return modp_eliminate(M.copy(), p, M.shape[1])[0]


def frac_modp(x, p=2147483647):
    x = Fraction(x)
    return x.numerator % p * pow(x.denominator % p, p - 2, p) % p


def exact_laplacian(net, binary_state):
    """integer matrix scale*L(state) for a 0/1 state; all parameters are exact decimals"""
    s = net.s
    g1, g2b, g2on = Fraction(str(s.g1)), Fraction(str(s.g2_base)), Fraction(str(s.g2_on))
    yb, yr = Fraction(str(s.y_base)), Fraction(str(s.y_ridge))
    bond, vals = {}, []
    for (u, v, i, j, kind) in net.edges:
        g = g1 if kind == 1 else g2b + Fraction(int(binary_state[i]) + int(binary_state[j]), 2)*(g2on - g2b)
        bond[(u, v)] = g; vals.append(g)
    shunt = [(yb + int(binary_state[i])*yr)/3 for i in range(net.N)]
    scale = 1
    for x in vals + shunt:
        scale = math.lcm(scale, x.denominator)
    Lz = np.zeros((3*net.N, 3*net.N), dtype=np.int64)
    for (u, v), g in bond.items():
        gi = int(g*scale)
        Lz[u, u] += gi; Lz[v, v] += gi; Lz[u, v] -= gi; Lz[v, u] -= gi
    for i in range(net.N):
        for k in range(3):
            Lz[3*i + k, 3*i + k] += int(shunt[i]*scale)
    return Lz


def harmonic_modp(net, binary_state, p=2147483647):
    """X = L_II^{-1} L_IB mod p (interior potentials are G = -X)"""
    Lz = exact_laplacian(net, binary_state)
    LII, LIB = Lz[np.ix_(net.I, net.I)], Lz[np.ix_(net.I, net.B)]
    r, Mx = modp_eliminate(np.concatenate([LII, LIB], 1), p, LII.shape[0])
    assert r == LII.shape[0]
    return Mx[:, LII.shape[0]:]


def A_modp(net, p=2147483647):
    X = harmonic_modp(net, np.zeros(net.N, dtype=int), p)
    Xr = X.reshape(net.N_I, 3, -1)
    return ((Xr*Xr % p).sum(1) % p).T*pow(3, p - 2, p) % p


def structural_rows(net):
    LIB = exact_laplacian(net, np.zeros(net.N, dtype=int))[np.ix_(net.I, net.B)]
    nz = [c for c in range(LIB.shape[1]) if LIB[:, c].any()]
    return len(net.B) - len(nz), len(nz), len({tuple(LIB[:, c]) for c in nz})


def diag_jacobian(net, state_all):
    """D[b, c] = d Lambda_bb / d s_c at a (continuous) state, D = w_b^T M_c w_b (Lemma 5.5)"""
    s = net.s
    L = net.laplacian(state_all, g2_switch=True)
    I, B = net.I, net.B
    G = -np.linalg.solve(L[np.ix_(I, I)], L[np.ix_(I, B)])
    W = np.zeros((3*net.N, len(B))); W[I, :] = G; W[B, np.arange(len(B))] = 1.0
    D = np.zeros((len(B), net.N_I))
    half = 0.5*(s.g2_on - s.g2_base)
    for pos, ci in enumerate(net.int_cells):
        d = s.y_ridge*(W[3*ci:3*ci + 3, :]**2).sum(0)/3.0
        for (u, v) in net.cell_bonds[ci]:
            d = d + half*(W[u, :] - W[v, :])**2
        D[:, pos] = d
    return D, G, L


def jacobian_modp(net, binary_state, p=2147483647):
    """the same Jacobian exactly mod p at a 0/1 state"""
    s = net.s
    X = harmonic_modp(net, binary_state, p)
    W = np.zeros((3*net.N, len(net.B)), dtype=np.int64)
    W[net.I, :] = (-X) % p
    W[net.B, np.arange(len(net.B))] = 1
    yr3 = frac_modp(Fraction(str(s.y_ridge))/3, p)
    half = frac_modp((Fraction(str(s.g2_on)) - Fraction(str(s.g2_base)))/2, p)
    J = np.zeros((len(net.B), net.N_I), dtype=np.int64)
    for pos, ci in enumerate(net.int_cells):
        acc = np.zeros(len(net.B), dtype=np.int64)
        for k in range(3):
            w = W[3*ci + k, :]; acc = (acc + w*w % p) % p
        acc = acc*yr3 % p
        for (u, v) in net.cell_bonds[ci]:
            d = (W[u, :] - W[v, :]) % p
            acc = (acc + (d*d % p)*half % p) % p
        J[:, pos] = acc
    return J


def active_nodes(net):
    L0 = net.laplacian(np.zeros(net.N), g2_switch=False)
    return np.array([b for b in range(len(net.B)) if L0[np.ix_(net.I, [net.B[b]])].any()])


def consistent_patterns(CK, target, cap=64, node_limit=5_000_000):
    """all 0/1 vectors x with CK x = target (exact DFS with line-sum propagation); returns (solutions, complete?)"""
    K = CK.shape[1]
    lines_of = [np.flatnonzero(CK[:, j]) for j in range(K)]
    order = sorted(range(K), key=lambda j: -len(lines_of[j]))
    rem = np.array(target, dtype=np.int64).copy()
    free = CK.sum(1).astype(np.int64)
    if (rem < 0).any() or (rem > free).any():
        return [], True
    sols, x, nodes = [], np.zeros(K, dtype=np.int64), [0]
    def dfs(i):
        if len(sols) >= cap or nodes[0] > node_limit:
            return
        nodes[0] += 1
        if i == K:
            if not rem.any():
                sols.append(x.copy())
            return
        j = order[i]; Lj = lines_of[j]
        for v in (1, 0):
            rem[Lj] -= v; free[Lj] -= 1
            if (rem[Lj] >= 0).all() and (rem[Lj] <= free[Lj]).all():
                x[j] = v; dfs(i + 1); x[j] = 0
            rem[Lj] += v; free[Lj] += 1
    dfs(0)
    return sols, nodes[0] <= node_limit and len(sols) < cap


def gate_variants():
    """Revision 5.0 reported K = 14 for the four-family gate on the test ridge: try every family subset and threshold."""
    net = Network(Spec())
    st = net.interior_vector(RIDGES)
    fams = {"m": net.m, "n": net.n, "m+n": net.p, "m-n": net.q}
    found = {}
    for k in range(1, 5):
        for sub in itertools.combinations(fams, k):
            for thr in (0.5, 1.5, 2.5):                        # at least 1, 2 or 3 ridge cells per line
                ok = np.ones(net.N_I, dtype=bool)
                for name in sub:
                    f = fams[name]
                    counts = {v: st[f == v].sum() for v in np.unique(f)}
                    ok &= np.array([counts[v] > thr for v in f])
                found[(sub, thr)] = int(ok.sum())
    full = found[(("m", "n", "m+n", "m-n"), 0.5)]
    assert full == 16
    print(f"[Sec.5.4] gate variants: K = {full} for the four-family gate; {len(found)} family-subset/threshold variants "
          f"give K in {sorted(set(found.values()))}; K = 14 occurs: {14 in found.values()}")


def verify_monotonicity(lin):
    """Lemma 5.5: d Lambda_bb / d s_c = w_b^T M_c w_b > 0 on active nodes, = 0 on inactive ones."""
    for spec, label in ((Spec(), "design 1e-7 S"), (BASELINE, "baseline 1e-6 S")):
        net = Network(spec)
        act = active_nodes(net)
        inact = np.setdiff1d(np.arange(len(net.B)), act)
        rng = np.random.default_rng(11)
        states = {"s=0": np.zeros(net.N), "s*": net.state(RIDGES)}
        for t in range(6):
            x = np.zeros(net.N); x[net.int_cells] = rng.uniform(0, 1, net.N_I); states[f"cont{t}"] = x
        for t in range(4):
            x = np.zeros(net.N); x[net.int_cells] = (rng.uniform(0, 1, net.N_I) < 0.3); states[f"bin{t}"] = x
        mins, zero_ok, dmp_ok = {}, True, True
        for name, x in states.items():
            D, G, L = diag_jacobian(net, x)
            LII = L[np.ix_(net.I, net.I)]
            off = LII - np.diag(np.diag(LII))
            dmp_ok &= bool((off <= 0).all() and (LII.sum(1) > 0).all() and (G[:, act] > 0).all())
            zero_ok &= bool(np.abs(D[inact, :]).max() == 0)
            mins[name] = float(D[act, :].min())
        adj = {k: set() for k in range(len(net.I))}
        idx = {int(v): k for k, v in enumerate(net.I)}
        for (u, v, i, j, kind) in net.edges:
            if u in idx and v in idx:
                adj[idx[u]].add(idx[v]); adj[idx[v]].add(idx[u])
        seen, stack = {0}, [0]
        while stack:
            for w in adj[stack.pop()] - seen:
                seen.add(w); stack.append(w)
        assert dmp_ok and len(seen) == len(net.I) and min(mins.values()) > 0 and zero_ok
        print(f"[Sec.5.4] Lemma 5.5 {label}: Z-matrix, strictly diag. dominant, connected, G[:, active] > 0 at "
              f"{len(states)} states; min dLambda_bb/ds_c over {len(act)} active nodes: s=0 {mins['s=0']:.3e}, "
              f"s* {mins['s*']:.3e}, all states {min(mins.values()):.3e} S; exactly 0 on {len(inact)} inactive nodes")
    net = Network(Spec())
    rng = np.random.default_rng(5)
    worst = 0.0
    for name, x in (("s*", net.state(RIDGES)),
                    ("random", np.where(np.isin(np.arange(net.N), net.int_cells), rng.uniform(0, 1, net.N), 0.0))):
        D = diag_jacobian(net, x)[0]
        for pos in rng.choice(net.N_I, 12, replace=False):
            ci = net.int_cells[pos]
            xc = x.astype(complex); xc[ci] += 1e-30j
            d = np.diag(net.schur(net.laplacian(xc))).imag/1e-30
            nz = D[:, pos] != 0
            worst = max(worst, float(np.max(np.abs(d[nz] - D[nz, pos])/np.abs(D[nz, pos]))))
    assert worst < 1e-9                             # bound stated after Lemma 5.5
    print(f"[Sec.5.4] Lemma 5.5 formula w_b^T M_c w_b vs complex-step derivative (24 columns, 2 states): "
          f"max rel. error <= {up(worst, 2)}")
    claimed, tried = 1.054e-13, []
    for spec, lab in ((Spec(), "design"), (BASELINE, "baseline")):
        net = Network(spec); act = active_nodes(net)
        n4 = Network(replace(spec, chord_families=4)); c4 = n4.gate(n4.C @ n4.interior_vector(RIDGES), 0.5)
        rp = [net.pos[c] for c in RIDGES]
        for sname, x in (("s=0", np.zeros(net.N)), ("s*", net.state(RIDGES))):
            Da = diag_jacobian(net, x)[0][act]
            for k, v in {"min all": Da.min(), "min candidates": Da[:, c4].min(), "min ridge cells": Da[:, rp].min(),
                         "min_c max_b": Da.max(0).min(), "min all / y_ridge": Da.min()/spec.y_ridge,
                         "min candidates / y_ridge": Da[:, c4].min()/spec.y_ridge,
                         "min_c max_b / y_ridge": Da.max(0).min()/spec.y_ridge}.items():
                tried.append((abs(v/claimed - 1), f"{lab} {sname} {k}", v))
    tried.sort()
    print(f"[Sec.5.4] Rev.5.0 value 1.054e-13: not reproduced by {len(tried)} conventions; closest {tried[0][1]} = "
          f"{tried[0][2]:.4e} ({100*tried[0][0]:.0f}% off)")
    net = Network(Spec()); L0, G0 = lin[Spec().y_base][:2]
    act = active_nodes(net)
    rng = np.random.default_rng(3)
    deep = [c for c in net.pos if radius(c) <= net.s.R_hex - 2]
    integrand_ok, diff_ok, n_res = True, True, 0
    for _ in range(10):
        base = [deep[i] for i in rng.choice(len(deep), 8, replace=False)]
        extra = [c for c in deep if c not in base][rng.integers(len(deep) - 8)]
        for tt in (1e-6, 1e-4, 1e-2, 0.25, 0.5, 1.0):
            x = net.state(base); x[net.c2i[extra]] = tt
            integrand_ok &= bool((diag_jacobian(net, x)[0][act, net.pos[extra]] > 0).all())
        lo_, hi_ = forward_diag(net, L0, G0, base), forward_diag(net, L0, G0, base + [extra])
        res = np.abs(hi_[act] - lo_[act]) > 1e-8*np.maximum(np.abs(hi_[act]), np.abs(lo_[act]))
        n_res += int(res.sum()); diff_ok &= bool((hi_[act][res] > lo_[act][res]).all())
    assert integrand_ok and diff_ok
    print(f"[Sec.5.4] Lemma 5.5(2) (10 random s <= s'): integrand > 0 along every segment; "
          f"{n_res} resolvable increases, all positive")


def verify_nonlinear_rank(lin):
    """Theorem 5.6: data collapse to 45 functions; rank J(s) <= 45 everywhere, = 45 at s*."""
    for spec, label in ((Spec(), "4 fam"), (replace(Spec(), chord_families=3), "3 fam")):
        net = Network(spec)
        L0, G0 = lin[spec.y_base][:2]
        st = net.interior_vector(RIDGES)
        cand = net.gate(net.C @ st, 0.5)
        Jp = jacobian_modp(net, net.state(RIDGES).astype(int))
        Jf = diag_jacobian(net, net.state(RIDGES))[0]
        CKi = net.C[:, cand].astype(np.int64)
        r_aug, r_J, r_C = rank_modp(np.concatenate([Jp[:, cand], CKi])), rank_modp(Jp[:, cand]), rank_modp(CKi)
        w_c = float(np.abs(forward_diag(net, L0, G0, RIDGES)).max())/spec.y_ridge
        sv = np.linalg.svd(np.vstack([Jf[:, cand]/spec.y_ridge, w_c*net.C[:, cand]]), compute_uv=False)
        assert r_aug == len(cand)
        print(f"[Sec.5.4] Thm 5.6 {label}: K = {len(cand)}, rank C_K = {r_C}, rank J(s*)_K = {r_J}, "
              f"rank [J(s*)_K; w_c C_K] = {r_aug} (exact); sigma_min/sigma_max = {sv[-1]/sv[0]:.2e}")
        if label == "4 fam":
            r_full = rank_modp(Jp)
            LIB = net.laplacian(np.zeros(net.N), g2_switch=False)[np.ix_(net.I, net.B)]
            groups = {}
            for b in range(len(net.B)):
                if LIB[:, b].any():
                    groups.setdefault(tuple(LIB[:, b]), []).append(b)
            same_rows = all(np.allclose(Jf[g[0]], Jf[g[i]], rtol=1e-12, atol=0) for g in groups.values() for i in range(len(g)))
            rngJ = np.random.default_rng(9); same_data = True
            for _ in range(3):
                cells_r = [c for c in net.pos if radius(c) <= spec.R_hex - 2 and rngJ.random() < 0.15]
                fd = forward_diag(net, L0, G0, cells_r)
                same_data &= all(np.allclose(fd[g[0]], fd[g[i]], rtol=1e-10, atol=0) for g in groups.values() for i in range(len(g)))
            assert r_full == len(groups) == 45 and same_rows and same_data
            sizes = sorted(len(g) for g in groups.values())
            pair_ok = all(np.count_nonzero(LIB[:, g[0]]) == 1 for g in groups.values() if len(g) == 2)
            const_ok = True
            for _ in range(2):
                cells_r = [c for c in net.pos if radius(c) <= spec.R_hex - 2 and rngJ.random() < 0.15]
                Lr = net.laplacian(net.state(cells_r))
                X = np.linalg.solve(Lr[np.ix_(net.I, net.I)], Lr[np.ix_(net.I, net.B)])
                diagr = np.diag(Lr[np.ix_(net.B, net.B)]) - np.einsum('ib,ib->b', Lr[np.ix_(net.I, net.B)], X)
                inactive = [b for b in range(len(net.B)) if not LIB[:, b].any()]
                const_ok &= bool(np.array_equal(diagr[inactive], np.diag(net.laplacian(np.zeros(net.N))[np.ix_(net.B, net.B)])[inactive]))
            assert max(sizes) <= 2 and pair_ok and const_ok
            print(f"[Sec.5.4] Thm 5.6: the 69 active nodes form {sizes.count(2)} pairs bonded to the same single interior node "
                  f"and {sizes.count(1)} singletons; Lambda_bb of the 75 inactive nodes is independent of s (2 random ridges)")
            print(f"[Sec.5.4] Thm 5.6: rank J(s*) over all {net.N_I} cells = {r_full} (exact: mod-p lower bound = "
                  f"{len(groups)} groups of identical rows); Delta Lambda_bb identical within each group for 3 random ridges")


def verify_screening(lin):
    """Theorem 5.7: four-family switching pair inside closed hexagonal ridge loops."""
    net = Network(Spec()); L0, G0 = lin[Spec().y_base][:2]
    plus, minus = switching_component(net)
    comp = {c: 1 for c in plus}; comp.update({c: -1 for c in minus})
    mom = {(a, b): sum(v*(m**a)*(n**b) for (m, n), v in comp.items()) for a in range(5) for b in range(5 - a)}
    assert all(mom[k] == 0 for k in mom if sum(k) <= 3) and any(mom[k] != 0 for k in mom if sum(k) == 4)
    print(f"[Sec.5.4] Thm 5.7(1): switching component has {len(plus) + len(minus)} cells ({len(plus)} at +1, "
          f"{len(minus)} at -1, not 16: the shifts satisfy (0,1)+(1,0) = (1,1)); moments of order <= 3 vanish, order 4 does not")
    def sep_abs(RA, RB):
        sa, sb = forward_diag(net, L0, G0, RA), forward_diag(net, L0, G0, RB)
        return float(np.linalg.norm(sa - sb)), float(np.linalg.norm(sa - sb)/(net.s.sigma_rel*np.abs(sa).max()))
    rows = []
    for r in (6, 5, 4, 3):
        ring = [c for c in net.pos if radius(c) == r]
        recs = []
        for a in range(-8, 9):
            for b in range(-8, 9):
                A_ = [(x + a, y + b) for x, y in plus]; B_ = [(x + a, y + b) for x, y in minus]
                if all(c in net.pos and radius(c) < r for c in A_ + B_):
                    fa, fr = sep_abs(ring + A_, ring + B_); ia, ir = sep_abs(A_, B_)
                    recs.append((min(r - radius(c) for c in A_ + B_), fa, fr, fa/ia, a, b))
        recs = np.array(recs)
        rows.append((r, np.median(recs[:, 1]), np.median(recs[:, 2])))
        print(f"[Sec.5.4] Thm 5.7 loop radius {r} ({len(ring)} cells): {len(recs)} placements; median separation "
              f"{np.median(recs[:, 1]):.2e} S = {np.median(recs[:, 2]):.2e} std (max {recs[:, 2].max():.2e} std); "
              f"loop factor {np.median(recs[:, 3]):.1e}")
        if r == 6:
            print("[Sec.5.4] Thm 5.7 radius-6 loop, by depth of the pair inside the wall: " + "; ".join(
                f"{int(d)} cell(s): {int((recs[:, 0] == d).sum())} placements, median {np.median(recs[recs[:, 0] == d, 2]):.2e} std"
                for d in sorted(set(recs[:, 0]))))
            z = 2*Q_inv(0.01)                                   # d/sigma needed for P_err <= 1 %
            snr0 = 20*np.log10(1/net.s.sigma_rel)
            for tag, i in (("largest", int(np.argmax(recs[:, 2]))), ("smallest", int(np.argmin(recs[:, 2])))):
                print(f"[Sec.5.4] Thm 5.7 radius-6 loop, {tag} separation: {recs[i, 2]:.3e} std ({recs[i, 1]:.2e} S) at shift "
                      f"({int(recs[i, 4])}, {int(recs[i, 5])}), depth {int(recs[i, 0])}; P_err <= 1% needs d >= {z:.4f} std, "
                      f"i.e. SNR >= {dn(snr0 + 20*np.log10(z/recs[i, 2]), 4)} dB (now {snr0:.2f} dB)")
    absd = np.array([x[1] for x in rows]); reld = np.array([x[2] for x in rows])
    assert np.all(np.diff(absd) < 0)
    print(f"[Sec.5.4] Thm 5.7: absolute separation falls monotonically inward (log10 slope "
          f"{np.polyfit(np.arange(len(absd)), np.log10(absd), 1)[0]:.2f} per ring); in noise units monotone: "
          f"{bool(np.all(np.diff(reld) < 0))}")


def clopper_pearson(k, n, alpha=0.05):
    """exact two-sided (1 - alpha) confidence interval for a binomial proportion with k of n (k small)"""
    cdf = lambda j, p: sum(math.comb(n, i)*p**i*(1 - p)**(n - i) for i in range(j + 1))
    def solve(f):                                   # f decreasing in p on (0, 1)
        lo, hi = 0.0, 1.0
        for _ in range(100):
            mid = 0.5*(lo + hi)
            lo, hi = (mid, hi) if f(mid) > 0 else (lo, mid)
        return 0.5*(lo + hi)
    p_lo = 0.0 if k == 0 else solve(lambda p: alpha/2 - (1 - cdf(k - 1, p)))
    p_hi = solve(lambda p: cdf(k, p) - alpha/2)
    return p_lo, p_hi


DIRS6 = [(1, 0), (0, 1), (-1, 1), (-1, 0), (0, -1), (1, -1)]


def is_tree(cells):
    """connected and loop-free: the 6-neighbour adjacency graph induced on the cells is a tree"""
    cs = set(cells)
    edges = sum((c[0] + d[0], c[1] + d[1]) in cs for c in cs for d in DIRS6)//2
    seen, todo = set(), [next(iter(cs))]
    while todo:
        c = todo.pop()
        if c not in seen:
            seen.add(c); todo += [(c[0] + d[0], c[1] + d[1]) for d in DIRS6 if (c[0] + d[0], c[1] + d[1]) in cs]
    return len(seen) == len(cs) and edges == len(cs) - 1


def decode_study(net, L0, G0, rng, sampler, n_ridges, trials, label):
    """four-family gate, exact enumeration of chord-consistent patterns, ML decoding over them"""
    s = net.s
    n_unique = n_ok = n_trials = 0
    failures = []
    for k in range(n_ridges):
        ridge = sampler()
        st = net.interior_vector(ridge)
        cand = net.gate(net.C @ st, 0.5)
        sols, complete = consistent_patterns(net.C[:, cand], net.C[:, cand] @ st[cand], cap=256)
        if not complete:
            failures.append((k, ridge, len(cand), f">={len(sols)}", 0, [], []))
            n_trials += trials
            continue
        n_unique += (len(sols) == 1)
        dL = forward_diag(net, L0, G0, ridge)
        sd = s.sigma_rel*np.abs(dL).max()
        sigs = {tuple(np.flatnonzero(x)): forward_diag(net, L0, G0, net.cells_of([cand[j] for j in np.flatnonzero(x)])) for x in sols}
        keys = list(sigs); S = np.array([sigs[q] for q in keys])
        true_key = tuple(j for j, c in enumerate(cand) if st[c] > 0.5)
        ok = 0
        for _ in range(trials):
            dLn = dL + rng.normal(0, sd, size=dL.shape)
            vals = st*s.y_ridge + rng.normal(0, s.sigma_rel*s.y_ridge, net.N_I)
            P = net.C @ vals
            if net.gate(P, 0.5*s.y_ridge) != cand or not np.array_equal(np.rint(P/s.y_ridge), net.C @ st):
                continue
            ok += keys[int(np.argmin(((S - dLn)**2).sum(1)))] == true_key
        n_ok += ok; n_trials += trials
        if len(sols) > 1 or ok < trials:
            alts = [net.cells_of([cand[j] for j in q]) for q in keys if q != true_key]
            d = [float(np.linalg.norm(sigs[q] - dL)/sd) for q in keys if q != true_key]
            dS = [float(np.linalg.norm(sigs[q] - dL)) for q in keys if q != true_key]
            failures.append((k, ridge, len(cand), len(sols), ok, alts, d, dS))
    a_lo, a_hi = clopper_pearson(n_ridges - n_unique, n_ridges)
    b_lo, b_hi = clopper_pearson(n_trials - n_ok, n_trials)
    print(f"[Sec.5.4] Rem 5.8 ({label}): {n_ridges} seeded ridges: unique chord solution {n_unique}/{n_ridges} "
          f"(non-unique rate 95% CI [{dn(100*a_lo)}, {up(100*a_hi)}] %); ML exact {n_ok}/{n_trials} "
          f"(error rate 95% CI [{dn(100*b_lo)}, {up(100*b_hi)}] %)")
    for (k, ridge, K, ns, ok, alts, d, *dS) in failures:
        trees = [is_tree(a) for a in alts]
        print(f"[Sec.5.4]   ridge {k}: {len(ridge)} cells, K = {K}, {ns} chord-consistent patterns, ML {ok}/{trials}, "
              f"separations {[round(x, 3) for x in d]} std" + (f" = {[float(f'{x:.2e}') for x in dS[0]]} S; alternatives "
              f"connected and loop-free: {sum(trees)}/{len(trees)}; {trials} Q(d_min/2) = {trials*Q(min(d)/2):.1f}"
              if dS else ""))
    return n_unique, n_ok, n_trials


def ridge_study(lin, n_ridges=100, trials=200):
    """Remark 5.8: (a) seeded stroke ridges; (b) Revision 5.0's class: connected, loop-free 12-cell ridges."""
    net = Network(Spec()); L0, G0 = lin[Spec().y_base][:2]
    inner = [c for c in net.pos if radius(c) <= 5]
    inner_set = set(inner)

    rng = np.random.default_rng(2026)
    def sample_strokes():
        cells = set()
        for _ in range(int(rng.integers(2, 4))):
            c = inner[int(rng.integers(len(inner)))]; d = int(rng.integers(6))
            for _ in range(int(rng.integers(4, 9))):
                if c in net.pos and radius(c) <= 5:
                    cells.add(c)
                if rng.random() < 0.25:
                    d = (d + int(rng.choice([-1, 1]))) % 6
                c = (c[0] + DIRS6[d][0], c[1] + DIRS6[d][1])
        return sorted(cells)
    decode_study(net, L0, G0, rng, sample_strokes, n_ridges, trials, "strokes")

    rng_t = np.random.default_rng(1205)
    def sample_tree(size=12):
        while True:
            cells = {inner[int(rng_t.integers(len(inner)))]}
            while len(cells) < size:
                front = sorted({(c[0] + d[0], c[1] + d[1]) for c in cells for d in DIRS6} - cells)
                front = [f for f in front if f in inner_set
                         and sum((f[0] + d[0], f[1] + d[1]) in cells for d in DIRS6) == 1]
                if not front:
                    break
                cells.add(front[int(rng_t.integers(len(front)))])
            if len(cells) == size:
                assert is_tree(cells)
                return sorted(cells)
    decode_study(net, L0, G0, rng_t, sample_tree, n_ridges, trials, "connected loop-free, 12 cells")


def rank_family():
    """Proposition 5.9: exact ranks for R_hex = 2..8."""
    print("[Sec.5.4] Prop 5.9:  R | interior | active | distinct | rank A_I | rank C3 | rank C4 | rank aug3 | rank aug4")
    fmt = lambda lo, hi: f"{lo}" if lo == hi else f"[{lo},{hi}]"
    for Rh in range(2, 9):
        net3, net4 = Network(replace(BASELINE, R_hex=Rh)), Network(replace(BASELINE, R_hex=Rh, chord_families=4))
        Ap = A_modp(net3)
        n0, nnz, ndist = structural_rows(net3)
        rA = rank_modp(Ap)
        rC3, rC4 = rank_exact(net3.C.astype(int)), rank_exact(net4.C.astype(int))
        r3 = rank_modp(np.concatenate([Ap, net3.C.astype(np.int64)])); r4 = rank_modp(np.concatenate([Ap, net4.C.astype(np.int64)]))
        hiA = min(ndist, net3.N_I)
        hi3, hi4 = min(hiA + rC3, net3.N_I), min(hiA + rC4, net3.N_I)
        if r3 < hi3:
            Ap2 = A_modp(net3, 2147483629)
            r3 = max(r3, rank_modp(np.concatenate([Ap2, net3.C.astype(np.int64)]), 2147483629))
        assert rA == hiA and r4 == hi4
        print(f"[Sec.5.4] Prop 5.9: {Rh:2d} | {net3.N_I:8d} | {nnz:6d} | {ndist:8d} | {fmt(rA, hiA):>8s} | {rC3:7d} | "
              f"{rC4:7d} | {fmt(r3, hi3):>9s} | {fmt(r4, hi4):>9s}")
        assert (nnz, ndist, rC3, rC4) == (9*Rh - 3, 6*Rh - 3, 6*Rh - 6, 10*Rh - 13) and net3.N_I == 3*Rh*Rh - 3*Rh + 1
        assert (Rh < 5 or r3 == 12*Rh - 9) and (Rh < 6 or r4 == 16*Rh - 16)
    print("[Sec.5.4] Prop 5.9: closed forms hold for 2 <= R <= 8: N_I = 3R^2-3R+1, active = 9R-3, distinct = 6R-3, "
          "rank C3 = 6R-6, rank C4 = 10R-13; rank aug3 = 12R-9 (R >= 5), rank aug4 = 16R-16 (R >= 6)")


# ======================= Revision 5.2: per-line split chords (Section 5.5) =======================
FAM4 = ("m", "n", "p", "q")
PRIMES = (2147483647, 2147483629, 2147483587, 2147483579)          # 2^31 - 1, -19, -61, -69


def split_rows(net, fams, k):
    """mod-k split chords: one 0/1 row per (family, line, residue r) with the cells of the line whose
    position along it (a unit-step coordinate: n on m-lines, m on the other three) is r mod k"""
    val = {"m": net.m, "n": net.n, "p": net.p, "q": net.q}
    pos = {"m": net.n, "n": net.m, "p": net.m, "q": net.m}
    rows = []
    for f in fams:
        for v in np.unique(val[f]):
            on = val[f] == v
            for r in range(k):
                sel = on & (pos[f] % k == r)
                if sel.any():
                    rows.append(sel.astype(np.int64))
    return np.array(rows)


def T_count(k, R):
    """number of translates of the k-dilated switching component inside the interior (Theorem 5.11)"""
    return 3*R*R - (10*k + 3)*R + 8*k*k + 5*k + 1 if R >= 2*k + 1 else 0


def dilated_translates(net, k):
    plus, minus = switching_component(Network(replace(net.s, chord_families=4)))
    out = []
    for a in range(-4*net.s.R_hex, 4*net.s.R_hex + 1):
        for b in range(-4*net.s.R_hex, 4*net.s.R_hex + 1):
            P_ = [(k*x + a, k*y + b) for x, y in plus]; M_ = [(k*x + a, k*y + b) for x, y in minus]
            if all(c in net.pos for c in P_ + M_):
                out.append((P_, M_))
    return out


def as_matrix(net, trans):
    Z = np.zeros((len(trans), net.N_I), dtype=np.int64)
    for i, (P_, M_) in enumerate(trans):
        for c in P_: Z[i, net.pos[c]] += 1
        for c in M_: Z[i, net.pos[c]] -= 1
    return Z


def verify_split_kernels():
    """Theorem 5.11: kernels of the split chords are spanned by translates of the dilated component."""
    net = Network(replace(Spec(), chord_families=4))
    Z = {}
    for k in (1, 2, 3, 4, 5):
        C = split_rows(net, FAM4, k)
        Zk = as_matrix(net, dilated_translates(net, k)); Z[k] = Zk
        rc = rank_exact(C.astype(int))
        if len(Zk):
            assert not (C @ Zk.T).any() and (np.abs(Zk).sum(1) == 8).all() and rank_exact(Zk.astype(int)) == len(Zk)
        assert len(Zk) == T_count(k, 8) == net.N_I - rc
        print(f"[Sec.5.5] Thm 5.11 k={k}: {C.shape[0]} split-chord rows, rank C^(k) = {rc} (exact); {len(Zk)} = T_k(8) translates "
              f"of the 8-cell k-dilated component lie in the kernel, are independent and span it (dim {net.N_I - rc})")
    C1, C2, C3 = (split_rows(net, FAM4, k) for k in (1, 2, 3))
    r1 = rank_exact(Z[1].astype(int))
    assert rank_exact(np.vstack([Z[1], Z[2], Z[3]]).astype(int)) == r1
    nz32 = sorted({int((C2 @ z != 0).sum()) for z in Z[3]}); nz23 = sorted({int((C3 @ z != 0).sum()) for z in Z[2]})
    assert nz32 == [24] and nz23 == [24]
    z = Z[1][0]
    vis = [int((C @ z != 0).sum()) for C in (C1, C2, C3)]
    print(f"[Sec.5.5] Thm 5.11: k=2 and k=3 translates lie in the span of the k=1 translates; kernels for k=2 and k=3 are not "
          f"nested (each z_3 translate has {nz32[0]} non-zero parity sums, each z_2 translate {nz23[0]} non-zero mod-3 sums); "
          f"the undilated component has {vis[0]}/{C1.shape[0]}, {vis[1]}/{C2.shape[0]} and {vis[2]}/{C3.shape[0]} non-zero split sums")


def verify_split_ranks():
    """Proposition 5.12: exact ranks of [A_I; C^(k)] at R_hex = 8, and the mod-3 family subsets."""
    for spec, lab in ((Spec(), "design"), (BASELINE, "baseline")):
        net = Network(replace(spec, chord_families=4))
        Ap, Ap2 = A_modp(net, PRIMES[0]), A_modp(net, PRIMES[1])
        rA = rank_modp(Ap)
        res = []
        for k in (1, 2, 3):
            C = split_rows(net, FAM4, k)
            rc = rank_exact(C.astype(int))
            lo = max(rank_modp(np.concatenate([Ap, C]), PRIMES[0]), rank_modp(np.concatenate([Ap2, C]), PRIMES[1]))
            assert lo == min(net.N_I, rA + rc)
            res.append(f"k={k}: rank C^(k) = {rc}, rank [A_I; C^(k)] = {lo}")
        print(f"[Sec.5.5] Prop 5.12 {lab} (y_base={spec.y_base:g}): " + "; ".join(res) + " (exact)")
    net = Network(replace(Spec(), chord_families=4))
    Ap, Ap2 = A_modp(net, PRIMES[0]), A_modp(net, PRIMES[1]); rA = rank_modp(Ap)
    out = []
    for size in (1, 2, 3, 4):
        for sub in itertools.combinations(FAM4, size):
            C = split_rows(net, sub, 3)
            rc = rank_exact(C.astype(int))
            lo = max(rank_modp(np.concatenate([Ap, C]), PRIMES[0]), rank_modp(np.concatenate([Ap2, C]), PRIMES[1]))
            assert lo == min(net.N_I, rA + rc)
            out.append(f"{''.join(sub)} {rc}/{lo}")
    print("[Sec.5.5] Prop 5.12 mod-3 family subsets, rank C | rank [A_I; C] (exact): " + ", ".join(out))


def verify_rcrit():
    """Corollary 5.15: R_crit(k) formula matches observed transitions."""
    import math
    def R_crit(k):
        return int((10*k + 9 + math.sqrt(4*k*k + 120*k + 33)) / 6) + 1
    def R_unc(k):
        return int((10*k + 21 + math.sqrt(4*k*k + 360*k + 429)) / 6) + 1
    def T(k, R):
        return 3*R*R - (10*k + 3)*R + 8*k*k + 5*k + 1 if R >= 2*k + 1 else 0
    def brute(k, d):  # smallest R0 with T_k(R) > d(R) for every R >= R0 (quadratic: beyond 4k+60 nothing changes)
        return 1 + max(R for R in range(1, 4*k + 60) if T(k, R) <= d(R))
    expected = {1: 6, 2: 8, 3: 10, 4: 13, 5: 15}
    for k, want in expected.items():
        got = R_crit(k)
        assert got == want, f"R_crit({k}) = {got}, expected {want}"
    for k in range(1, 501):
        assert R_crit(k) == brute(k, lambda R: 6*R - 3), k
        assert R_unc(k) == brute(k, lambda R: 18*R), k
        assert (R_crit(k) - 2*k == 7) == (k >= 58), k
    assert R_unc(3) == 16  # matches the R >= 16 of Theorem 5.14
    # asymptotic check: additive constant approaches 7
    for k in (100, 1000, 10000):
        gap = R_crit(k) - 2 * k
        assert 6 <= gap <= 7, f"R_crit({k}) - 2k = {gap}, expected 6..7"
    print(f"[Sec.5.5] Cor 5.15: R_crit(k) verified for k=1..5 "
          f"(1:6, 2:8, 3:10, 4:13, 5:15); R_crit(10000) - 20000 = {R_crit(10000)-20000} (asymptote 7)")


def verify_split_conditioning():
    """Remark 5.16: singular values of the stacked linear operators."""
    sv = lambda M: np.linalg.svd(M, compute_uv=False)
    for yb in (1e-7, 1e-6):
        net = Network(replace(Spec(), chord_families=4, y_base=yb))
        A = np.asarray(linear_model(net)[2], float)
        C2, C3 = split_rows(net, FAM4, 2), split_rows(net, FAM4, 3)
        s2 = sv(np.vstack([A, C2]))
        M = np.vstack([A, C3]); s3 = sv(M)
        Mn = M[np.linalg.norm(M, axis=1) > 0]
        sr = sv(Mn/np.linalg.norm(Mn, axis=1)[:, None]); scol = sv(M/np.linalg.norm(M, axis=0))
        _, _, Vt = np.linalg.svd(C3.astype(float)); sa = sv(A @ Vt[153:].T)
        if yb == 1e-6:
            assert s2[162] < 1e-14                  # Remark 5.16: sigma_163 at round-off level for y_base = 1e-6 S
        print(f"[Sec.5.5] Rem 5.16 y_base={yb:g}: [A_I; C^(2)] sigma_1 = {s2[0]:.3e}, sigma_163 = {s2[162]:.3e}, "
              f"sigma_164 = {s2[163]:.3e}; [A_I; C^(3)] sigma_max = {s3[0]:.3e}, sigma_min = {s3[-1]:.3e}, "
              f"kappa = {s3[0]/s3[-1]:.0f} (column-equilibrated {scol[0]/scol[-1]:.0f}, row-equilibrated {sr[0]/sr[-1]:.0f}); "
              f"A_I on ker C^(3): singular values {sa[0]:.2e} .. {sa[-1]:.2e}")
        trip = []
        for sub in (("m", "n", "q"), ("m", "p", "q"), ("n", "p", "q")):
            st = sv(np.vstack([A, split_rows(net, sub, 3)]))
            trip.append(f"{''.join(sub)} {st[0]/st[-1]:.2e}")
        print(f"[Sec.5.5] Rem 5.16 y_base={yb:g}: kappa of [A_I; C^(3)] on the full-rank three-family subsets: " + ", ".join(trip))


def verify_split_screening():
    """Theorem 5.13: dilated-component pairs, isolated and inside closed ridge loops (exact identity, any ridge)."""
    net = Network(replace(Spec(), chord_families=4))
    L0f = net.laplacian(np.zeros(net.N))
    def W_of(L):
        W = np.zeros((L.shape[0], len(net.B)))
        W[np.ix_(net.I, np.arange(len(net.B)))] = -np.linalg.solve(L[np.ix_(net.I, net.I)], L[np.ix_(net.I, net.B)])
        W[net.B, np.arange(len(net.B))] = 1.0
        return W
    W0 = W_of(L0f)
    def fwd(cells):                          # diag(Lambda(s) - Lambda(0)) = diag(W(s)^T (L(s) - L(0)) W(0))
        L = net.laplacian(net.state(cells))
        return np.einsum('ib,ib->b', W_of(L), (L - L0f) @ W0)
    L0, G0 = linear_model(net)[:2]
    ref = forward_diag(net, L0, G0, RIDGES)
    assert np.abs(fwd(RIDGES) - ref).max() < 1e-9*np.abs(ref).max()
    def sep(RA, RB):
        a, b = fwd(RA), fwd(RB)
        return float(np.linalg.norm(a - b)), float(np.linalg.norm(a - b)/(net.s.sigma_rel*np.abs(a).max()))
    for k, rr in ((2, 6), (3, 7)):
        ring = [c for c in net.pos if radius(c) == rr]
        tr = dilated_translates(net, k)
        iso = np.array([sep(P_, M_)[1] for P_, M_ in tr])
        enc = np.array([sep(ring + P_, ring + M_) for P_, M_ in tr if all(radius(c) < rr for c in P_ + M_)])
        print(f"[Sec.5.5] Thm 5.13 k={k}: {len(tr)} isolated pairs separated by {dn(iso.min())} .. {up(iso.max())} std; inside the "
              f"closed radius-{rr} ring ({len(ring)} cells) {len(enc)} pairs at {dn(enc[:, 1].min())} .. {up(enc[:, 1].max())} std "
              f"({dn(enc[:, 0].min(), 2)} .. {up(enc[:, 0].max(), 2)} S) -> every decoder errs w.p. >= "
              f"{dn(Q(enc[:, 1].max()/2), 4)} on one ridge of each pair")


def verify_split_scaling():
    """Theorem 5.14: mod-3 split chords plus boundary diagonal, R_hex = 2..15."""
    rows, keep = [], {}
    for R in range(2, 12):
        net = Network(replace(Spec(), R_hex=R, chord_families=4))
        C = split_rows(net, FAM4, 3)
        rc = rank_exact(C.astype(int))
        d = structural_rows(net)[2]
        Ap = A_modp(net)
        lo, hi = rank_modp(np.concatenate([Ap, C])), min(net.N_I, d + rc)
        T = len(dilated_translates(net, 3))
        assert T == T_count(3, R) == net.N_I - rc and d == 6*R - 3
        assert (R <= 9 and lo == net.N_I) or (R >= 10 and hi < net.N_I)
        rows.append(f"R={R}: {lo}" + ("" if lo == hi else f"..{hi}") + f"/{net.N_I}")
        keep[R] = (rc, d)
    print("[Sec.5.5] Thm 5.14 rank [A_I; C^(3)] / N_I: " + ", ".join(rows) + "  (full for R <= 9, deficient for R = 10, 11)")
    for R in range(12, 16):
        net = Network(replace(Spec(), R_hex=R, chord_families=4))
        d = structural_rows(net)[2]
        assert d == 6*R - 3 and T_count(3, R) - d == 3*R*R - 39*R + 91 > 0
    assert all(T_count(3, R) - 18*R == 3*R*R - 51*R + 88 > 0 for R in range(16, 200))
    print("[Sec.5.5] Thm 5.14: d(R) = 6R - 3 distinct L_IB columns for 2 <= R <= 15, and T_3(R) - d(R) = 3R^2 - 39R + 91 > 0 "
          "for 10 <= R <= 15; T_3(R) - 18R = 3R^2 - 51R + 88 > 0 for R >= 16 (18R boundary nodes)")
    net = Network(replace(Spec(), R_hex=10, chord_families=4))
    K = as_matrix(net, dilated_translates(net, 3)).T
    ranks = []
    for p in PRIMES:
        Ap = A_modp(net, p)
        AK = np.zeros((Ap.shape[0], K.shape[1]), dtype=np.int64)
        for j in range(K.shape[1]):
            for i in np.flatnonzero(K[:, j]):
                AK[:, j] = (AK[:, j] + Ap[:, i]*(int(K[i, j]) % p)) % p
        ranks.append(rank_modp(AK, p))
    A = np.asarray(linear_model(net)[2], float)
    s = np.linalg.svd(A @ K, compute_uv=False)
    rc10, d10 = keep[10]
    lo10, hi10 = rc10 + max(ranks), rc10 + min(d10, K.shape[1])
    print(f"[Sec.5.5] Thm 5.14 R=10: rank [A_I; C^(3)] = rank C^(3) ({rc10}) + rank(A_I K) with K the {K.shape[1]} kernel translates; "
          f"rank(A_I K) mod four primes = {ranks}; float64 singular values of A_I K {s[0]:.2e} .. {s[54]:.2e}, then "
          f"{s[55]:.1e}, {s[56]:.1e}, {s[57]:.1e} -> rank [A_I; C^(3)] in [{lo10}, {hi10}] of {net.N_I}")


def split_decoding_study():
    """Remark 5.17: standard versus mod-3 split chords on the same seeded ridges."""
    net = Network(Spec()); s = net.s
    L0, G0 = linear_model(net)[:2]
    inner = [c for c in net.pos if radius(c) <= 5]; inner_set = set(inner)
    def strokes(rng):
        cells = set()
        for _ in range(int(rng.integers(2, 4))):
            c = inner[int(rng.integers(len(inner)))]; d = int(rng.integers(6))
            for _ in range(int(rng.integers(4, 9))):
                if c in net.pos and radius(c) <= 5: cells.add(c)
                if rng.random() < 0.25: d = (d + int(rng.choice([-1, 1]))) % 6
                c = (c[0] + DIRS6[d][0], c[1] + DIRS6[d][1])
        return sorted(cells)
    def tree(rng, size=12):
        while True:
            cells = {inner[int(rng.integers(len(inner)))]}
            while len(cells) < size:
                front = sorted({(c[0] + d[0], c[1] + d[1]) for c in cells for d in DIRS6} - cells)
                front = [f for f in front if f in inner_set and sum((f[0] + d[0], f[1] + d[1]) in cells for d in DIRS6) == 1]
                if not front: break
                cells.add(front[int(rng.integers(len(front)))])
            if len(cells) == size: return sorted(cells)
    def gate(C, P, thr):
        ok = np.ones(C.shape[1], bool)
        for r in np.flatnonzero(~(P > thr)):
            ok &= ~C[r].astype(bool)
        return [int(c) for c in np.flatnonzero(ok)]
    designs = (("standard", split_rows(net, FAM4, 1)), ("mod-3", split_rows(net, FAM4, 3)))
    for cname, sampler, seed in (("strokes", strokes, 2026), ("connected loop-free 12", tree, 1205)):
        rng_r = np.random.default_rng(seed)
        ridges = [sampler(rng_r) for _ in range(100)]          # the same ridges for both designs
        for dname, C in designs:
            rng = np.random.default_rng(77)
            uniq = n_ok = 0; multi = []
            for k, ridge in enumerate(ridges):
                st = net.interior_vector(ridge)
                cand = gate(C, C @ st, 0.5)
                sols, complete = consistent_patterns(C[:, cand], C[:, cand] @ st[cand], cap=256)
                assert complete
                uniq += len(sols) == 1
                dL = forward_diag(net, L0, G0, ridge); sd = s.sigma_rel*np.abs(dL).max()
                sigs = {tuple(np.flatnonzero(x)): forward_diag(net, L0, G0, net.cells_of([cand[j] for j in np.flatnonzero(x)])) for x in sols}
                keys = list(sigs); S = np.array([sigs[q] for q in keys])
                tk = tuple(j for j, c in enumerate(cand) if st[c] > 0.5)
                ok = 0
                for _ in range(200):
                    dLn = dL + rng.normal(0, sd, size=dL.shape)
                    vals = st*s.y_ridge + rng.normal(0, s.sigma_rel*s.y_ridge, net.N_I)
                    P = C @ vals
                    if gate(C, P, 0.5*s.y_ridge) != cand or not np.array_equal(np.rint(P/s.y_ridge), C @ st):
                        continue
                    ok += keys[int(np.argmin(((S - dLn)**2).sum(1)))] == tk
                n_ok += ok
                if len(sols) > 1:
                    multi.append(f"{len(sols)} patterns, ML {ok}/200, separations "
                                 f"{[round(float(np.linalg.norm(sigs[q] - dL)/sd), 2) for q in keys if q != tk]} std")
            print(f"[Sec.5.5] Rem 5.17 {cname}, {dname} chords ({C.shape[0]} rows): unique chord data {uniq}/100, "
                  f"ML exact {n_ok}/20000" + (f"; non-unique: {'; '.join(multi)}" if multi else ""))


# ======================= Revision 5.3: convex ridges and the cross-ratio criterion (Section 5.6) =======================
def _cross(o, a, b):
    return (a[0] - o[0])*(b[1] - o[1]) - (a[1] - o[1])*(b[0] - o[0])


def _hull(pts):
    pts = sorted(set(pts))
    if len(pts) <= 2:
        return pts
    lo, hi = [], []
    for p in pts:
        while len(lo) >= 2 and _cross(lo[-2], lo[-1], p) <= 0: lo.pop()
        lo.append(p)
    for p in reversed(pts):
        while len(hi) >= 2 and _cross(hi[-2], hi[-1], p) <= 0: hi.pop()
        hi.append(p)
    return lo[:-1] + hi[:-1]


def _inside(H, q):
    if len(H) == 1:
        return q == H[0]
    if len(H) == 2:
        a, b = H
        return _cross(a, b, q) == 0 and min(a[0], b[0]) <= q[0] <= max(a[0], b[0]) and min(a[1], b[1]) <= q[1] <= max(a[1], b[1])
    return all(_cross(H[i], H[(i + 1) % len(H)], q) >= 0 for i in range(len(H)))


def lattice_closure(pts):
    """Z^2 intersected with conv(pts): the smallest convex lattice set containing pts"""
    H = _hull(pts)
    xs, ys = [p[0] for p in pts], [p[1] for p in pts]
    return frozenset((x, y) for x in range(min(xs), max(xs) + 1) for y in range(min(ys), max(ys) + 1) if _inside(H, (x, y)))


def chord_data(cells, fams):
    lines = {"m": lambda c: c[0], "n": lambda c: c[1], "p": lambda c: c[0] + c[1], "q": lambda c: c[0] - c[1]}
    out = []
    for f in fams:
        cnt = {}
        for c in cells:
            cnt[lines[f](c)] = cnt.get(lines[f](c), 0) + 1
        out.append(tuple(sorted(cnt.items())))
    return tuple(out)


def verify_convex(lin):
    """Proposition 5.18: convex ridges with identical chord data, and how well the boundary separates them."""
    slopes = (-1.0, 0.0, 1.0)                       # with infinity: cross ratio (c-a)/(c-b) for a<b<c<inf
    cr = (slopes[2] - slopes[0])/(slopes[2] - slopes[1])
    assert cr == 2.0
    net = Network(Spec()); s = net.s
    W = [c for c in net.pos if radius(c) <= 3]
    convex = set()
    for k in (1, 2, 3, 4):
        for pts in itertools.combinations(W, k):
            convex.add(lattice_closure(pts))
    res = {}
    for fams in (("m", "n", "p"), ("m", "n", "p", "q")):
        groups = {}
        for F in convex:
            groups.setdefault(chord_data(F, fams), []).append(F)
        coll = [g for g in groups.values() if len(g) > 1]
        res[len(fams)] = (len(coll), min(len(g[0]) for g in coll))
    assert res[3][1] == 4 and res[4][1] == 8
    print(f"[Sec.5.6] Prop 5.18: cross ratio of the four chord slopes = {cr:g}; {len(convex)} convex lattice sets generated by "
          f"<= 4 of the {len(W)} cells within radius 3: three families: {res[3][0]} chord-data classes with distinct convex "
          f"sets, smallest {res[3][1]} cells; four families: {res[4][0]} classes, smallest {res[4][1]} cells")
    plus, minus = switching_component(net)
    interior = [(1, 0), (1, 1), (2, 0), (2, 1)]
    P0, M0 = frozenset(plus + interior), frozenset(minus + interior)
    assert lattice_closure(sorted(plus)) == P0 and lattice_closure(sorted(minus)) == M0
    pairs = []
    for a in range(-2*s.R_hex, 2*s.R_hex + 1):
        for b in range(-2*s.R_hex, 2*s.R_hex + 1):
            F1 = frozenset((x + a, y + b) for x, y in P0); F2 = frozenset((x + a, y + b) for x, y in M0)
            if all(c in net.pos for c in F1 | F2):
                assert lattice_closure(sorted(F1)) == F1 and lattice_closure(sorted(F2)) == F2
                assert chord_data(F1, FAM4) == chord_data(F2, FAM4) and F1 != F2
                pairs.append((F1, F2))
    L0, G0 = lin[s.y_base][:2]
    def sep(A, B):
        sa, sb = forward_diag(net, L0, G0, sorted(A)), forward_diag(net, L0, G0, sorted(B))
        return float(np.linalg.norm(sa - sb)/(s.sigma_rel*np.abs(sa).max()))
    iso = np.array([sep(F1, F2) for F1, F2 in pairs if all(radius(c) <= s.R_hex - 2 for c in F1 | F2)])
    ring = [c for c in net.pos if radius(c) == s.R_hex - 2]
    enc = np.array([sep(set(ring) | F1, set(ring) | F2) for F1, F2 in pairs if all(radius(c) < s.R_hex - 2 for c in F1 | F2)])
    print(f"[Sec.5.6] Prop 5.18: the two lattice parallelograms S+ u I and S- u I (I = 4 interior points) are convex, have "
          f"identical chord data in all four families and differ by the switching component; {len(pairs)} placements in the "
          f"interior; isolated ({len(iso)} placements, radius <= {s.R_hex - 2}): separation {dn(iso.min(), 4)} .. "
          f"{up(iso.max(), 4)} std; inside the closed radius-{s.R_hex - 2} ring ({len(enc)} placements): {dn(enc.min(), 3)} .. "
          f"{up(enc.max(), 3)} std -> every decoder errs w.p. >= {dn(Q(float(up(enc.max(), 3))/2), 4)} on one ridge of each pair")


if __name__ == "__main__":
    t_start = time.perf_counter()
    design = Spec()
    verify_sec2()
    verify_sec3(design)
    verify_sec4(design)
    configs = ((BASELINE, "R3.9 baseline (y_base=1e-6, 3 fam)"),
               (replace(BASELINE, y_base=1.0e-7), "y_base=1e-7, 3 fam"),
               (replace(BASELINE, chord_families=4), "y_base=1e-6, 4 fam"),
               (design, "R4.0 design (y_base=1e-7, 4 fam)"))
    lin, mods, summary = {}, {}, {}
    for spec, label in configs:
        net = Network(spec)
        if spec.y_base not in lin:
            lin[spec.y_base] = linear_model(net)
            verify_frechet(net, *lin[spec.y_base])
            mods[spec.y_base] = modular_model(net)
        L0, G0, A = lin[spec.y_base]
        Ci, rC, rAug = verify_rank(net, mods[spec.y_base], label)
        summary[label] = stage_b(net, L0, G0, A, mods[spec.y_base], Ci, label, full_dtn=(spec == BASELINE))
        summary[label]['rank'] = rAug
    for spec, label in configs[2:]:
        ghost_pairs(Network(spec), *lin[spec.y_base][:2], label)
    print("[Sec.5] summary: config | rank [A_I; C] | K | chord-consistent patterns | ML exact / 2000")
    for label, r in summary.items():
        print(f"[Sec.5]   {label:36s} | {r['rank']:3d} | {r['K']:2d} | {r['patterns']} | {r['ok_ml']}")
    gate_variants()
    verify_monotonicity(lin)
    verify_nonlinear_rank(lin)
    verify_screening(lin)
    ridge_study(lin)
    rank_family()
    verify_split_kernels()
    verify_split_ranks()
    verify_split_screening()
    verify_split_scaling()
    verify_rcrit()
    verify_split_conditioning()
    split_decoding_study()
    verify_convex(lin)
    print(f"[done] {time.perf_counter() - t_start:.0f} s")
\end{lstlisting}

\noindent Console output of the listing (verbatim):
\begin{lstlisting}[style=consolestyle]
[Sec.2] rho_eff = 680.6832 kg/m^3; W_i = 47.83, 11.90, 30.85, 9.37, 0.04 %; [Sec.3.1] beta = 0.9944, eps_eff = 5.5831, g33 = 40.46 mV m/N, V_wake = 70.16 mV
[Sec.3] Thm3.1 max|E_num-E_formula| over 400 k <= 2.3e-15; gap at K = 1.0800; Thm3.2 CLS residual <= 2.7e-17
[Sec.3] 217-cell flake: delta_b=+0.18: 217 states below gap, 0 in gap (trivial); delta_b=-0.18: 204 below, 21 in gap (obstructed)
[Sec.4] |q_phi| a2 = 4.5650; intracell offsets psi2-psi1 = 211.22 deg, psi3-psi1 = 334.44 deg; ||C_up||_2 = 1.9382
[Sec.4] up  : 65879 perturbed triangles (586 not in v3.8 anchor rule of 65296); max||E0^-1 dE|| = 0.55089; min det(I+M) = 0.57854; min det(E)/det_ref = 0.9988; min det(E0)/det_ref = 1.0000
[Sec.4] up  : max ||E0^-1|| / (sqrt2/(ell R/rho)) = 1.0266; min det(E)/float-error bound >= 1.7e+10
[Sec.4] down: 66140 perturbed triangles (1126 not in v3.8 anchor rule of 65296); max||E0^-1 dE|| = 0.52212; min det(I+M) = 0.69278; min det(E)/det_ref = 0.9993; min det(E0)/det_ref = 1.0000
[Sec.4] down: max ||E0^-1|| / (sqrt2/(ell R/rho)) = 1.0089; min det(E)/float-error bound >= 2.9e+10
[Sec.4] first-order sup: up 0.5451, down 0.5424; rigorous first-order bound 4*mu*eta*a2/ell (mu=1.025): up 0.8000, down 0.5559
[Sec.5] y_base=1e-06: Thm5.1 complex-step max rel. err over 169 columns <= 1.1e-10; remainder identity residual <= 8.4e-9; ||R2|| = 6.66e-11 <= 2.08e-10
[Sec.5] y_base=1e-06: lambda_min(L_II) = 3.389e-07 S, cond(L_II) = 1.92e+05; remainder bound at dy = y_ridge: 1.30e-01 S
[Sec.5] R3.9 baseline (y_base=1e-6, 3 fam): 144 boundary nodes, 75 without interior neighbour, 45 distinct rows -> rank A_I = 45; 3 chord families, 45 lines, rank C = 42; rank [A_I; w_c C] = 87 of 169; dim ker = 82 (exact)
[Sec.5] R3.9 baseline (y_base=1e-6, 3 fam): chord gate 45 lines, p_chord <= Q(6.455) <= 5.42e-11, FWER <= 2.44e-9 (line-length-aware <= 2.48e-10); max|dLambda| = 5.9003e-14 S, w_c = 2.36e-11; ||dL - A s y_r||/||dL|| = 21.9
[Sec.5] R3.9 baseline (y_base=1e-6, 3 fam): K = 17, rank C_K = 13, rank [A_K; w_c C_K] = 17 (exact), sigma(A_K) ratio 6.8e-05; 3 binary pattern(s) share the chord data; ghost separation 0.5265, 0.5265 std -> any decoder errs w.p. >= 0.396
[Sec.5] R3.9 baseline (y_base=1e-6, 3 fam): full Schur matrix: linearised Jacobian rank 169 of 169 (exact); ghost separation 0.794, 0.794 std -> decoder error >= 0.345; sigma(A_K) = [6.30e-10, 4.31e-14]
[Sec.5] R3.9 baseline (y_base=1e-6, 3 fam): MC 2000 @ 34.0 dB: K in [17, 17]; ridge-LS (lambda=1e-10) 0/2000 (max |s_hat| <= 2.3e-10); maximum likelihood 1239/2000 (0.007 ms/trial select)
[Sec.5] y_base=1e-07: Thm5.1 complex-step max rel. err over 169 columns <= 3.5e-10; remainder identity residual <= 8.4e-10; ||R2|| = 2.91e-09 <= 2.58e-08
[Sec.5] y_base=1e-07: lambda_min(L_II) = 3.888e-08 S, cond(L_II) = 1.67e+06; remainder bound at dy = y_ridge: 1.61e+01 S
[Sec.5] y_base=1e-7, 3 fam: 144 boundary nodes, 75 without interior neighbour, 45 distinct rows -> rank A_I = 45; 3 chord families, 45 lines, rank C = 42; rank [A_I; w_c C] = 87 of 169; dim ker = 82 (exact)
[Sec.5] y_base=1e-7, 3 fam: chord gate 45 lines, p_chord <= Q(6.455) <= 5.42e-11, FWER <= 2.44e-9 (line-length-aware <= 2.48e-10); max|dLambda| = 1.6549e-10 S, w_c = 6.62e-08; ||dL - A s y_r||/||dL|| = 1115.4
[Sec.5] y_base=1e-7, 3 fam: K = 17, rank C_K = 13, rank [A_K; w_c C_K] = 17 (exact), sigma(A_K) ratio 8.9e-04; 3 binary pattern(s) share the chord data; ghost separation 4.4595, 4.4598 std -> any decoder errs w.p. >= 0.0128
[Sec.5] y_base=1e-7, 3 fam: MC 2000 @ 34.0 dB: K in [17, 17]; ridge-LS (lambda=1e-10) 0/2000 (max |s_hat| <= 0.0014); maximum likelihood 1973/2000 (0.006 ms/trial select)
[Sec.5] y_base=1e-6, 4 fam: 144 boundary nodes, 75 without interior neighbour, 45 distinct rows -> rank A_I = 45; 4 chord families, 74 lines, rank C = 67; rank [A_I; w_c C] = 112 of 169; dim ker = 57 (exact)
[Sec.5] y_base=1e-6, 4 fam: chord gate 74 lines, p_chord <= Q(6.455) <= 5.42e-11, FWER <= 4.01e-9 (line-length-aware <= 2.48e-10); max|dLambda| = 5.9003e-14 S, w_c = 2.36e-11; ||dL - A s y_r||/||dL|| = 21.9
[Sec.5] y_base=1e-6, 4 fam: K = 16, rank C_K = 16, rank [A_K; w_c C_K] = 16 (exact), sigma(A_K) ratio 3.1e-02; 1 binary pattern(s) share the chord data; the chord data determine the ridge
[Sec.5] y_base=1e-6, 4 fam: MC 2000 @ 34.0 dB: K in [16, 16]; ridge-LS (lambda=1e-10) 0/2000 (max |s_hat| <= 2.4e-10); maximum likelihood 2000/2000 (0.006 ms/trial select)
[Sec.5] R4.0 design (y_base=1e-7, 4 fam): 144 boundary nodes, 75 without interior neighbour, 45 distinct rows -> rank A_I = 45; 4 chord families, 74 lines, rank C = 67; rank [A_I; w_c C] = 112 of 169; dim ker = 57 (exact)
[Sec.5] R4.0 design (y_base=1e-7, 4 fam): chord gate 74 lines, p_chord <= Q(6.455) <= 5.42e-11, FWER <= 4.01e-9 (line-length-aware <= 2.48e-10); max|dLambda| = 1.6549e-10 S, w_c = 6.62e-08; ||dL - A s y_r||/||dL|| = 1115.4
[Sec.5] R4.0 design (y_base=1e-7, 4 fam): K = 16, rank C_K = 16, rank [A_K; w_c C_K] = 16 (exact), sigma(A_K) ratio 1.3e-03; 1 binary pattern(s) share the chord data; the chord data determine the ridge
[Sec.5] R4.0 design (y_base=1e-7, 4 fam): MC 2000 @ 34.0 dB: K in [16, 16]; ridge-LS (lambda=1e-10) 0/2000 (max |s_hat| <= 0.0015); maximum likelihood 2000/2000 (0.006 ms/trial select)
[Sec.5] Thm5.2(5) y_base=1e-6, 4 fam: 4+4-cell switching component [(0, 0), (1, 2), (2, -1), (3, 1)] vs [(0, 1), (1, -1), (2, 2), (3, 0)]; isolated (all cells >= 2 rings deep): 70 placements, min separation >= 44.62 std at shift (0, -3)
[Sec.5] Thm5.2(5) y_base=1e-6, 4 fam: inside a closed 36-cell ridge loop: 44 placements, separation 2.1e-09 to <= 0.0019 std -> any decoder errs w.p. >= 0.4996 on one of each pair
[Sec.5] Thm5.2(5) R4.0 design (y_base=1e-7, 4 fam): 4+4-cell switching component [(0, 0), (1, 2), (2, -1), (3, 1)] vs [(0, 1), (1, -1), (2, 2), (3, 0)]; isolated (all cells >= 2 rings deep): 70 placements, min separation >= 30.09 std at shift (-3, 0)
[Sec.5] Thm5.2(5) R4.0 design (y_base=1e-7, 4 fam): inside a closed 36-cell ridge loop: 44 placements, separation 1.4e-08 to <= 6.3e-4 std -> any decoder errs w.p. >= 0.4998 on one of each pair
[Sec.5] summary: config | rank [A_I; C] | K | chord-consistent patterns | ML exact / 2000
[Sec.5]   R3.9 baseline (y_base=1e-6, 3 fam)   |  87 | 17 | 3 | 1239
[Sec.5]   y_base=1e-7, 3 fam                   |  87 | 17 | 3 | 1973
[Sec.5]   y_base=1e-6, 4 fam                   | 112 | 16 | 1 | 2000
[Sec.5]   R4.0 design (y_base=1e-7, 4 fam)     | 112 | 16 | 1 | 2000
[Sec.5.4] gate variants: K = 16 for the four-family gate; 45 family-subset/threshold variants give K in [0, 3, 6, 7, 8, 9, 10, 12, 13, 15, 16, 17, 18, 19, 21, 24, 25, 30, 42, 43, 49, 59, 68, 69, 81, 93]; K = 14 occurs: False
[Sec.5.4] Lemma 5.5 design 1e-7 S: Z-matrix, strictly diag. dominant, connected, G[:, active] > 0 at 12 states; min dLambda_bb/ds_c over 69 active nodes: s=0 1.633e-16, s* 1.908e-21, all states 5.624e-34 S; exactly 0 on 75 inactive nodes
[Sec.5.4] Lemma 5.5 baseline 1e-6 S: Z-matrix, strictly diag. dominant, connected, G[:, active] > 0 at 12 states; min dLambda_bb/ds_c over 69 active nodes: s=0 6.280e-33, s* 3.919e-34, all states 6.631e-35 S; exactly 0 on 75 inactive nodes
[Sec.5.4] Lemma 5.5 formula w_b^T M_c w_b vs complex-step derivative (24 columns, 2 states): max rel. error <= 1.7e-10
[Sec.5.4] Rev.5.0 value 1.054e-13: not reproduced by 28 conventions; closest design s=0 min all / y_ridge = 6.5313e-14 (38% off)
[Sec.5.4] Lemma 5.5(2) (10 random s <= s'): integrand > 0 along every segment; 381 resolvable increases, all positive
[Sec.5.4] Thm 5.6 4 fam: K = 16, rank C_K = 16, rank J(s*)_K = 16, rank [J(s*)_K; w_c C_K] = 16 (exact); sigma_min/sigma_max = 7.21e-03
[Sec.5.4] Thm 5.6: the 69 active nodes form 24 pairs bonded to the same single interior node and 21 singletons; Lambda_bb of the 75 inactive nodes is independent of s (2 random ridges)
[Sec.5.4] Thm 5.6: rank J(s*) over all 169 cells = 45 (exact: mod-p lower bound = 45 groups of identical rows); Delta Lambda_bb identical within each group for 3 random ridges
[Sec.5.4] Thm 5.6 3 fam: K = 17, rank C_K = 13, rank J(s*)_K = 17, rank [J(s*)_K; w_c C_K] = 17 (exact); sigma_min/sigma_max = 4.92e-08
[Sec.5.4] Thm 5.7(1): switching component has 8 cells (4 at +1, 4 at -1, not 16: the shifts satisfy (0,1)+(1,0) = (1,1)); moments of order <= 3 vanish, order 4 does not
[Sec.5.4] Thm 5.7 loop radius 6 (36 cells): 44 placements; median separation 2.34e-13 S = 2.99e-04 std (max 6.27e-04 std); loop factor 1.3e-04
[Sec.5.4] Thm 5.7 radius-6 loop, by depth of the pair inside the wall: 1 cell(s): 20 placements, median 3.70e-04 std; 2 cell(s): 14 placements, median 2.99e-04 std; 3 cell(s): 8 placements, median 7.66e-08 std; 4 cell(s): 2 placements, median 1.44e-08 std
[Sec.5.4] Thm 5.7 radius-6 loop, largest separation: 6.273e-04 std (4.91e-13 S) at shift (-1, -4), depth 1; P_err <= 1% needs d >= 4.6527 std, i.e. SNR >= 111.3 dB (now 33.98 dB)
[Sec.5.4] Thm 5.7 radius-6 loop, smallest separation: 1.445e-08 std (1.13e-17 S) at shift (-1, -1), depth 4; P_err <= 1% needs d >= 4.6527 std, i.e. SNR >= 204.1 dB (now 33.98 dB)
[Sec.5.4] Thm 5.7 loop radius 5 (30 cells): 24 placements; median separation 2.83e-14 S = 2.83e-04 std (max 5.10e-04 std); loop factor 1.1e-04
[Sec.5.4] Thm 5.7 loop radius 4 (24 cells): 10 placements; median separation 6.58e-15 S = 3.64e-04 std (max 5.17e-04 std); loop factor 1.0e-04
[Sec.5.4] Thm 5.7 loop radius 3 (18 cells): 2 placements; median separation 2.29e-15 S = 6.26e-04 std (max 6.26e-04 std); loop factor 1.0e-04
[Sec.5.4] Thm 5.7: absolute separation falls monotonically inward (log10 slope -0.67 per ring); in noise units monotone: False
[Sec.5.4] Rem 5.8 (strokes): 100 seeded ridges: unique chord solution 98/100 (non-unique rate 95% CI [0.243, 7.04] %); ML exact 20000/20000 (error rate 95% CI [0, 0.0185] %)
[Sec.5.4]   ridge 48: 13 cells, K = 28, 2 chord-consistent patterns, ML 200/200, separations [10.687] std = [9.08e-10] S; alternatives connected and loop-free: 0/1; 200 Q(d_min/2) = 0.0
[Sec.5.4]   ridge 60: 20 cells, K = 47, 3 chord-consistent patterns, ML 200/200, separations [13.443, 39.609] std = [1.15e-09, 3.4e-09] S; alternatives connected and loop-free: 0/2; 200 Q(d_min/2) = 0.0
[Sec.5.4] Rem 5.8 (connected loop-free, 12 cells): 100 seeded ridges: unique chord solution 95/100 (non-unique rate 95% CI [1.64, 11.3] %); ML exact 19978/20000 (error rate 95% CI [0.0689, 0.167] %)
[Sec.5.4]   ridge 3: 12 cells, K = 21, 2 chord-consistent patterns, ML 200/200, separations [26.044] std = [2.59e-09] S; alternatives connected and loop-free: 0/1; 200 Q(d_min/2) = 0.0
[Sec.5.4]   ridge 10: 12 cells, K = 20, 2 chord-consistent patterns, ML 197/200, separations [4.742] std = [4.03e-10] S; alternatives connected and loop-free: 0/1; 200 Q(d_min/2) = 1.8
[Sec.5.4]   ridge 22: 12 cells, K = 26, 2 chord-consistent patterns, ML 181/200, separations [2.376] std = [2.34e-10] S; alternatives connected and loop-free: 0/1; 200 Q(d_min/2) = 23.5
[Sec.5.4]   ridge 28: 12 cells, K = 25, 2 chord-consistent patterns, ML 200/200, separations [67.632] std = [5.87e-09] S; alternatives connected and loop-free: 0/1; 200 Q(d_min/2) = 0.0
[Sec.5.4]   ridge 42: 12 cells, K = 21, 2 chord-consistent patterns, ML 200/200, separations [49.173] std = [4.23e-09] S; alternatives connected and loop-free: 0/1; 200 Q(d_min/2) = 0.0
[Sec.5.4] Prop 5.9:  R | interior | active | distinct | rank A_I | rank C3 | rank C4 | rank aug3 | rank aug4
[Sec.5.4] Prop 5.9:  2 |        7 |     15 |        9 |        7 |       6 |       7 |         7 |         7
[Sec.5.4] Prop 5.9:  3 |       19 |     24 |       15 |       15 |      12 |      17 |        19 |        19
[Sec.5.4] Prop 5.9:  4 |       37 |     33 |       21 |       21 |      18 |      27 |   [35,37] |        37
[Sec.5.4] Prop 5.9:  5 |       61 |     42 |       27 |       27 |      24 |      37 |        51 |        61
[Sec.5.4] Prop 5.9:  6 |       91 |     51 |       33 |       33 |      30 |      47 |        63 |        80
[Sec.5.4] Prop 5.9:  7 |      127 |     60 |       39 |       39 |      36 |      57 |        75 |        96
[Sec.5.4] Prop 5.9:  8 |      169 |     69 |       45 |       45 |      42 |      67 |        87 |       112
[Sec.5.4] Prop 5.9: closed forms hold for 2 <= R <= 8: N_I = 3R^2-3R+1, active = 9R-3, distinct = 6R-3, rank C3 = 6R-6, rank C4 = 10R-13; rank aug3 = 12R-9 (R >= 5), rank aug4 = 16R-16 (R >= 6)
[Sec.5.5] Thm 5.11 k=1: 74 split-chord rows, rank C^(k) = 67 (exact); 102 = T_k(8) translates of the 8-cell k-dilated component lie in the kernel, are independent and span it (dim 102)
[Sec.5.5] Thm 5.11 k=2: 146 split-chord rows, rank C^(k) = 118 (exact); 51 = T_k(8) translates of the 8-cell k-dilated component lie in the kernel, are independent and span it (dim 51)
[Sec.5.5] Thm 5.11 k=3: 216 split-chord rows, rank C^(k) = 153 (exact); 16 = T_k(8) translates of the 8-cell k-dilated component lie in the kernel, are independent and span it (dim 16)
[Sec.5.5] Thm 5.11 k=4: 284 split-chord rows, rank C^(k) = 169 (exact); 0 = T_k(8) translates of the 8-cell k-dilated component lie in the kernel, are independent and span it (dim 0)
[Sec.5.5] Thm 5.11 k=5: 350 split-chord rows, rank C^(k) = 169 (exact); 0 = T_k(8) translates of the 8-cell k-dilated component lie in the kernel, are independent and span it (dim 0)
[Sec.5.5] Thm 5.11: k=2 and k=3 translates lie in the span of the k=1 translates; kernels for k=2 and k=3 are not nested (each z_3 translate has 24 non-zero parity sums, each z_2 translate 24 non-zero mod-3 sums); the undilated component has 0/74, 24/146 and 24/216 non-zero split sums
[Sec.5.5] Prop 5.12 design (y_base=1e-07): k=1: rank C^(k) = 67, rank [A_I; C^(k)] = 112; k=2: rank C^(k) = 118, rank [A_I; C^(k)] = 163; k=3: rank C^(k) = 153, rank [A_I; C^(k)] = 169 (exact)
[Sec.5.5] Prop 5.12 baseline (y_base=1e-06): k=1: rank C^(k) = 67, rank [A_I; C^(k)] = 112; k=2: rank C^(k) = 118, rank [A_I; C^(k)] = 163; k=3: rank C^(k) = 153, rank [A_I; C^(k)] = 169 (exact)
[Sec.5.5] Prop 5.12 mod-3 family subsets, rank C | rank [A_I; C] (exact): m 45/90, n 45/90, p 45/90, q 81/126, mn 81/126, mp 81/126, mq 117/162, np 81/126, nq 117/162, pq 108/153, mnp 108/153, mnq 144/169, mpq 135/169, npq 135/169, mnpq 153/169
[Sec.5.5] Thm 5.13 k=2: 51 isolated pairs separated by 75.5 .. 142 std; inside the closed radius-6 ring (36 cells) 13 pairs at 5.40e-4 .. 0.00147 std (4.2e-13 .. 1.2e-12 S) -> every decoder errs w.p. >= 0.4997 on one ridge of each pair
[Sec.5.5] Thm 5.13 k=3: 16 isolated pairs separated by 99.9 .. 147 std; inside the closed radius-7 ring (42 cells) 4 pairs at 1.45 .. 1.78 std (2.5e-5 .. 3.2e-5 S) -> every decoder errs w.p. >= 0.1868 on one ridge of each pair
[Sec.5.5] Thm 5.14 rank [A_I; C^(3)] / N_I: R=2: 7/7, R=3: 19/19, R=4: 37/37, R=5: 61/61, R=6: 91/91, R=7: 127/127, R=8: 169/169, R=9: 217/217, R=10: 268..270/271, R=11: 306/331  (full for R <= 9, deficient for R = 10, 11)
[Sec.5.5] Thm 5.14: d(R) = 6R - 3 distinct L_IB columns for 2 <= R <= 15, and T_3(R) - d(R) = 3R^2 - 39R + 91 > 0 for 10 <= R <= 15; T_3(R) - 18R = 3R^2 - 51R + 88 > 0 for R >= 16 (18R boundary nodes)
[Sec.5.5] Thm 5.14 R=10: rank [A_I; C^(3)] = rank C^(3) (213) + rank(A_I K) with K the 58 kernel translates; rank(A_I K) mod four primes = [55, 55, 55, 55]; float64 singular values of A_I K 1.22e-01 .. 1.85e-11, then 1.9e-16, 1.8e-19, 1.1e-19 -> rank [A_I; C^(3)] in [268, 270] of 271
[Sec.5.5] Cor 5.15: R_crit(k) verified for k=1..5 (1:6, 2:8, 3:10, 4:13, 5:15); R_crit(10000) - 20000 = 7 (asymptote 7)
[Sec.5.5] Rem 5.16 y_base=1e-07: [A_I; C^(2)] sigma_1 = 4.668e+00, sigma_163 = 8.373e-10, sigma_164 = 9.192e-16; [A_I; C^(3)] sigma_max = 3.797e+00, sigma_min = 8.794e-05, kappa = 43172 (column-equilibrated 43171, row-equilibrated 2150); A_I on ker C^(3): singular values 3.99e-02 .. 8.80e-05
[Sec.5.5] Rem 5.16 y_base=1e-07: kappa of [A_I; C^(3)] on the full-rank three-family subsets: mnq 9.17e+05, mpq 2.78e+07, npq 2.78e+07
[Sec.5.5] Rem 5.16 y_base=1e-06: [A_I; C^(2)] sigma_1 = 4.668e+00, sigma_163 = 9.776e-16, sigma_164 = 9.157e-16; [A_I; C^(3)] sigma_max = 3.797e+00, sigma_min = 7.241e-08, kappa = 52431411 (column-equilibrated 52431388, row-equilibrated 291401); A_I on ker C^(3): singular values 4.91e-03 .. 7.24e-08
[Sec.5.5] Rem 5.16 y_base=1e-06: kappa of [A_I; C^(3)] on the full-rank three-family subsets: mnq 5.78e+09, mpq 1.18e+12, npq 1.18e+12
[Sec.5.5] Rem 5.17 strokes, standard chords (74 rows): unique chord data 100/100, ML exact 20000/20000
[Sec.5.5] Rem 5.17 strokes, mod-3 chords (216 rows): unique chord data 100/100, ML exact 20000/20000
[Sec.5.5] Rem 5.17 connected loop-free 12, standard chords (74 rows): unique chord data 97/100, ML exact 19928/20000; non-unique: 4 patterns, ML 200/200, separations [100.66, 108.2, 110.6] std; 2 patterns, ML 128/200, separations [0.78] std; 2 patterns, ML 200/200, separations [5.18] std
[Sec.5.5] Rem 5.17 connected loop-free 12, mod-3 chords (216 rows): unique chord data 100/100, ML exact 20000/20000
[Sec.5.6] Prop 5.18: cross ratio of the four chord slopes = 2; 48853 convex lattice sets generated by <= 4 of the 37 cells within radius 3: three families: 344 chord-data classes with distinct convex sets, smallest 4 cells; four families: 15 classes, smallest 8 cells
[Sec.5.6] Prop 5.18: the two lattice parallelograms S+ u I and S- u I (I = 4 interior points) are convex, have identical chord data in all four families and differ by the switching component; 102 placements in the interior; isolated (70 placements, radius <= 6): separation 30.08 .. 89.84 std; inside the closed radius-6 ring (44 placements): 1.44e-8 .. 6.14e-4 std -> every decoder errs w.p. >= 0.4998 on one ridge of each pair
[done] 157 s
\end{lstlisting}

\subsection{Supplementary exact-arithmetic checks}\label{app:supp}
Requires \texttt{numpy} and \texttt{python-flint}, and imports the listing above; runtime $\approx$40 s. Parts (a), (b) and (d) evaluate the boundary signatures directly as differences of Schur complements in 256-bit ball arithmetic and print rigorous enclosures: (a) the test ridge and its two ghosts at the baseline (Theorem~\ref{thm:ghost}(3)), and (b), (d) the most and least separated enclosed pairs of Theorem~\ref{thm:screen}(2). Part (c) computes the three-family augmented rank at $R_{\text{hex}} = 4$ exactly over $\mathbb{Q}$ (Proposition~\ref{prop:family}).

\begin{lstlisting}[language=Python,style=cleanstyle,caption={Supplementary exact-arithmetic checks.}]
"""Supplementary exact-arithmetic checks for Revision 5.5 (independent of numpy linear algebra).

(a) Theorem 5.2(3), Revision-3.9 baseline (y_base = 1e-6, three chord families): boundary signatures
    diag(Lambda(s) - Lambda(0)) of the test ridge and its two chord-consistent ghosts.
(b) Theorem 5.2(5), Revision-4.0 design (y_base = 1e-7, four families): the ridge pair formed by a
    closed 36-cell ridge loop plus either half of the four-family switching component (the placement
    with the LARGEST float64 separation among the 44 tested), and (d) the placement (-1, -1) with the SMALLEST.
(c) Proposition 5.9: the three-family augmented rank at R_hex = 4, computed exactly over Q.
(a), (b) and (d) use 256-bit ball arithmetic (python-flint) directly on differences of Schur complements and
report rigorous enclosures; (c) uses exact rational matrices.

Usage (same folder as biokagome_v55_verify.py):  pip install python-flint numpy
                                                  python supplementary_exact_checks.py
Runtime about 40 s.
"""
import itertools, math
import numpy as np
from flint import arb, arb_mat, ctx
from biokagome_v55_verify import Spec, BASELINE, Network, RIDGES, Q, switching_component, exact_laplacian, replace
import flint

ctx.prec = 256


def use(spec):
    global s, net, I, B, g1, g2b, g2on, yb, yr
    s = spec
    net = Network(s)
    I, B = [int(x) for x in net.I], [int(x) for x in net.B]
    g1, g2b, g2on = arb(str(s.g1)), arb(str(s.g2_base)), arb(str(s.g2_on))
    yb, yr = arb(str(s.y_base)), arb(str(s.y_ridge))


use(BASELINE)


def laplacian(state):
    L = [dict() for _ in range(3*net.N)]
    def add(a, b, v):
        L[a][b] = L[a].get(b, arb(0)) + v
    for (u, v, i, j, kind) in net.edges:
        g = g1 if kind == 1 else g2b + arb(int(state[i]) + int(state[j]))/2*(g2on - g2b)
        add(u, u, g); add(v, v, g); add(u, v, -g); add(v, u, -g)
    for i in range(net.N):
        for k in range(3):
            add(3*i + k, 3*i + k, (yb + arb(int(state[i]))*yr)/3)
    return L


def schur_diag(state):
    L = laplacian(state)
    X = arb_mat([[L[a].get(b, arb(0)) for b in I] for a in I]).solve(
        arb_mat([[L[a].get(b, arb(0)) for b in B] for a in I]))
    out = []
    for c, b in enumerate(B):
        acc = L[b].get(b, arb(0))
        for r, a in enumerate(I):
            if a in L[b]:
                acc -= L[b][a]*X[r, c]
        out.append(acc)
    return out


st = np.zeros(net.N, dtype=int)
for rc in RIDGES:
    st[net.c2i[rc]] = 1
sti = st[net.int_cells]
cand = net.gate(net.C @ sti, 0.5)
CK = net.C[:, cand]
cnt = CK @ sti[cand]
pats = [np.array(b) for b in itertools.product((0, 1), repeat=len(cand)) if np.array_equal(CK @ np.array(b), cnt)]
i_true = [i for i, b in enumerate(pats) if np.array_equal(b, sti[cand])][0]
base = schur_diag(np.zeros(net.N, dtype=int))
sig = []
for b in pats:
    x = np.zeros(net.N, dtype=int)
    for c, v in zip(cand, b):
        x[net.int_cells[c]] = int(v)
    d = schur_diag(x)
    sig.append([d[k] - base[k] for k in range(len(B))])


def norm_enclosure(vec):
    mid = math.sqrt(sum(float(v.mid())**2 for v in vec))
    rad = math.sqrt(sum(float(v.rad())**2 for v in vec))
    return mid - rad, mid + rad


nscale = max(abs(v) for v in sig[i_true])
sd = arb(str(s.sigma_rel))*nscale
lo_sd, hi_sd = float(sd.mid()) - float(sd.rad()), float(sd.mid()) + float(sd.rad())
print("(a) baseline, three families")
print(f"max|dLambda_bb| = {nscale.str(10, radius=True)} S")
for i in range(len(pats)):
    if i != i_true:
        lo, hi = norm_enclosure([sig[i][k] - sig[i_true][k] for k in range(len(B))])
        print(f"ghost {i}: separation / noise std in [{lo/hi_sd:.6f}, {hi/lo_sd:.6f}]")
g = [i for i in range(len(pats)) if i != i_true]
lo, hi = norm_enclosure([sig[g[0]][k] - sig[g[1]][k] for k in range(len(B))])
print(f"ghost vs ghost: separation / noise std in [{lo/hi_sd:.3e}, {hi/lo_sd:.3e}]")

# ---------------- (b) closed ridge loop screens a four-family ghost pair ----------------
use(Spec())
plus, minus = switching_component(net)
depth = lambda c: max(abs(c[0]), abs(c[1]), abs(c[0] + c[1]))
ring = [c for c in net.pos if depth(c) == s.R_hex - 2]
def state(cells):
    x = np.zeros(net.N, dtype=int)
    for c in cells:
        x[net.c2i[c]] = 1
    return x
base = schur_diag(np.zeros(net.N, dtype=int))
for (a, b), tag in (((-1, -4), "(b) largest"), ((-1, -1), "(d) smallest")):   # float64 extremes among the 44 placements
    RA = ring + [(x + a, y + b) for x, y in plus]
    RB = ring + [(x + a, y + b) for x, y in minus]
    dA = schur_diag(state(RA)); dB = schur_diag(state(RB))
    sigA = [dA[k] - base[k] for k in range(len(B))]
    diff = [dA[k] - dB[k] for k in range(len(B))]
    nscale = max(abs(v) for v in sigA)
    sd = arb(str(s.sigma_rel))*nscale
    lo_sd, hi_sd = float(sd.mid()) - float(sd.rad()), float(sd.mid()) + float(sd.rad())
    lo, hi = norm_enclosure(diff)
    print(f"{tag} separation, Revision-4.0 design, closed 36-cell loop + switching component at shift ({a}, {b})")
    print(f"max|dLambda_bb| = {nscale.str(10, radius=True)} S; separation in [{lo:.6e}, {hi:.6e}] S")
    print(f"pair separation / noise std in [{lo/hi_sd:.6e}, {hi/lo_sd:.6e}] -> any decoder errs w.p. >= {math.floor(Q(hi/lo_sd/2)*1e6)/1e6:.6f}")

# ---------------- (c) exact rank of [A_I; C3] at R_hex = 4 over Q ----------------
net4 = Network(replace(BASELINE, R_hex=4))
Lz = exact_laplacian(net4, np.zeros(net4.N, dtype=int))
X = flint.fmpz_mat(Lz[np.ix_(net4.I, net4.I)].tolist()).solve(flint.fmpz_mat(Lz[np.ix_(net4.I, net4.B)].tolist()))
A4 = [[sum((X[3*c + k, b]**2 for k in range(3)), flint.fmpq(0))/3 for c in range(net4.N_I)] for b in range(len(net4.B))]
M4 = flint.fmpq_mat(A4 + [[flint.fmpq(int(x)) for x in row] for row in net4.C])
print(f"(c) R_hex = 4, three families: rank over Q of [A_I; C] = {M4.rank()} of {net4.N_I} (exact rational arithmetic)")
\end{lstlisting}

\noindent Console output (verbatim):
\begin{lstlisting}[style=consolestyle]
(a) baseline, three families
max|dLambda_bb| = [5.900262477e-14 +/- 1.68e-24] S
ghost 1: separation / noise std in [0.526493, 0.526493]
ghost 2: separation / noise std in [0.526535, 0.526535]
ghost vs ghost: separation / noise std in [2.228e-04, 2.228e-04]
(b) largest separation, Revision-4.0 design, closed 36-cell loop + switching component at shift (-1, -4)
max|dLambda_bb| = [3.914459884e-8 +/- 1.37e-18] S; separation in [4.911171e-13, 4.911171e-13] S
pair separation / noise std in [6.273115e-04, 6.273115e-04] -> any decoder errs w.p. >= 0.499874
(d) smallest separation, Revision-4.0 design, closed 36-cell loop + switching component at shift (-1, -1)
max|dLambda_bb| = [3.914433371e-8 +/- 2.41e-18] S; separation in [1.130908e-17, 1.130908e-17] S
pair separation / noise std in [1.444536e-08, 1.444536e-08] -> any decoder errs w.p. >= 0.499999
(c) R_hex = 4, three families: rank over Q of [A_I; C] = 35 of 37 (exact rational arithmetic)
\end{lstlisting}

\end{document}
